% taller-qmd-template.tex — plantilla pandoc/quarto para TALLERES del curso AED
% Basada en el driver de tesis main-es.tex (clase mitthesis parcheada UR), con
% portada y resumen a medida para trabajos de curso. Compilar con XeLaTeX
% (build.sh usa latexmk -pdfxe).
%
%   Metadatos por taller (YAML del .qmd): title, subtitle, fecha, resumen, palabras
%   (fecha: clave propia; quarto reformatea `date:` a ISO-8601)
%   Datos del curso (compartidos):         common/meta-taller-es.tex
%   Preliminares (portada/resumen/índice): common/frontmatter-taller.tex
\def\oaslang{es}
\DocumentMetadata{lang=es, pdfversion=1.7, pdfstandard=a-2b}
% Guardarraíl de motor: fontset=libertinus exige XeLaTeX (ver main-es.tex).
\RequirePackage{iftex}\RequireXeTeX
\documentclass[oneside, openany, fontset=libertinus]{mitthesis}
\input{common/preamble}

% --- Compatibilidad con la salida de pandoc ---
\usepackage{fancyvrb}   % entornos Shaded/Highlighting del código R
\usepackage{upquote}    % comillas rectas en verbatim
\usepackage{float}      % opción de flotante [H] que quarto usa en las figuras
\usepackage{color}
\usepackage{fancyvrb}
\newcommand{\VerbBar}{|}
\newcommand{\VERB}{\Verb[commandchars=\\\{\}]}
\DefineVerbatimEnvironment{Highlighting}{Verbatim}{commandchars=\\\{\}}
% Add ',fontsize=\small' for more characters per line
\usepackage{framed}
\definecolor{shadecolor}{RGB}{241,243,245}
\newenvironment{Shaded}{\begin{snugshade}}{\end{snugshade}}
\newcommand{\AlertTok}[1]{\textcolor[rgb]{0.68,0.00,0.00}{#1}}
\newcommand{\AnnotationTok}[1]{\textcolor[rgb]{0.37,0.37,0.37}{#1}}
\newcommand{\AttributeTok}[1]{\textcolor[rgb]{0.40,0.45,0.13}{#1}}
\newcommand{\BaseNTok}[1]{\textcolor[rgb]{0.68,0.00,0.00}{#1}}
\newcommand{\BuiltInTok}[1]{\textcolor[rgb]{0.00,0.23,0.31}{#1}}
\newcommand{\CharTok}[1]{\textcolor[rgb]{0.13,0.47,0.30}{#1}}
\newcommand{\CommentTok}[1]{\textcolor[rgb]{0.37,0.37,0.37}{#1}}
\newcommand{\CommentVarTok}[1]{\textcolor[rgb]{0.37,0.37,0.37}{\textit{#1}}}
\newcommand{\ConstantTok}[1]{\textcolor[rgb]{0.56,0.35,0.01}{#1}}
\newcommand{\ControlFlowTok}[1]{\textcolor[rgb]{0.00,0.23,0.31}{\textbf{#1}}}
\newcommand{\DataTypeTok}[1]{\textcolor[rgb]{0.68,0.00,0.00}{#1}}
\newcommand{\DecValTok}[1]{\textcolor[rgb]{0.68,0.00,0.00}{#1}}
\newcommand{\DocumentationTok}[1]{\textcolor[rgb]{0.37,0.37,0.37}{\textit{#1}}}
\newcommand{\ErrorTok}[1]{\textcolor[rgb]{0.68,0.00,0.00}{#1}}
\newcommand{\ExtensionTok}[1]{\textcolor[rgb]{0.00,0.23,0.31}{#1}}
\newcommand{\FloatTok}[1]{\textcolor[rgb]{0.68,0.00,0.00}{#1}}
\newcommand{\FunctionTok}[1]{\textcolor[rgb]{0.28,0.35,0.67}{#1}}
\newcommand{\ImportTok}[1]{\textcolor[rgb]{0.00,0.46,0.62}{#1}}
\newcommand{\InformationTok}[1]{\textcolor[rgb]{0.37,0.37,0.37}{#1}}
\newcommand{\KeywordTok}[1]{\textcolor[rgb]{0.00,0.23,0.31}{\textbf{#1}}}
\newcommand{\NormalTok}[1]{\textcolor[rgb]{0.00,0.23,0.31}{#1}}
\newcommand{\OperatorTok}[1]{\textcolor[rgb]{0.37,0.37,0.37}{#1}}
\newcommand{\OtherTok}[1]{\textcolor[rgb]{0.00,0.23,0.31}{#1}}
\newcommand{\PreprocessorTok}[1]{\textcolor[rgb]{0.68,0.00,0.00}{#1}}
\newcommand{\RegionMarkerTok}[1]{\textcolor[rgb]{0.00,0.23,0.31}{#1}}
\newcommand{\SpecialCharTok}[1]{\textcolor[rgb]{0.37,0.37,0.37}{#1}}
\newcommand{\SpecialStringTok}[1]{\textcolor[rgb]{0.13,0.47,0.30}{#1}}
\newcommand{\StringTok}[1]{\textcolor[rgb]{0.13,0.47,0.30}{#1}}
\newcommand{\VariableTok}[1]{\textcolor[rgb]{0.07,0.07,0.07}{#1}}
\newcommand{\VerbatimStringTok}[1]{\textcolor[rgb]{0.13,0.47,0.30}{#1}}
\newcommand{\WarningTok}[1]{\textcolor[rgb]{0.37,0.37,0.37}{\textit{#1}}}
\providecommand{\tightlist}{\setlength{\itemsep}{0pt}\setlength{\parskip}{0pt}}
\providecommand{\pandocbounded}[1]{#1}
% pandoc marca las tablas sin caption con \def\LTcaptype{none}: necesita el contador
\makeatletter\@ifundefined{c@none}{\newcounter{none}}{}\makeatother
% Código y resultado lado a lado (salidas cortas): \begin{codres} código \codressep resultado \end{codres}
\newenvironment{codres}%
  {\par\smallskip\noindent\begin{minipage}[t]{0.57\linewidth}\vspace{0pt}}%
  {\end{minipage}\par\medskip}
\newcommand{\codressep}{\end{minipage}\hfill\begin{minipage}[t]{0.40\linewidth}\vspace{0pt}}
% Bibliografía: \imprimirbibliografia la ubica donde se llame (p. ej. antes de
% \appendix); si el taller no la llama, se imprime al final como siempre.
\newif\ifbibimpresa
\newcommand{\imprimirbibliografia}{\printbibliography[heading=bibintoc]\global\bibimpresatrue}
% Paginación: \Needspace{alto} empieza en la página siguiente lo que no
% cabe entero (lo emite paginar() de R/taller.R antes de cada chunk).
\usepackage{needspace}
% Notas para la sesión: recuadro rojo pastel titulado «Notas» que se puede
% partir entre páginas. En el .qmd van como `::: notas` (filtro notas.lua).
% Interruptor global: \mostrarnotasfalse (o `ocultar-notas: true` en el
% YAML) las quita todas sin borrarlas del cuaderno.
\usepackage{comment}
\newif\ifmostrarnotas
\mostrarnotastrue

\definecolor{notasfondo}{HTML}{F8E1E3}
\definecolor{notasmarco}{HTML}{A6192E}
\ifmostrarnotas
  \newtcolorbox{notas}{breakable, enhanced, colback=notasfondo,
    colframe=notasmarco, coltitle=white, fonttitle=\bfseries,
    title={Notas}, boxrule=0.6pt, arc=2pt, left=7pt, right=7pt,
    top=5pt, bottom=5pt, before skip=12pt, after skip=12pt,
    before upper={\setlength{\parindent}{0pt}\setlength{\parskip}{4pt}}}
\else
  \excludecomment{notas}
\fi

% --- Resolución de recursos hacia la plantilla (layout taller-XX junto a template/) ---
\graphicspath{{assets/pics/}{figures/}{../template/assets/pics/}{../template/figures/}}
\addbibresource{../template/bib/refs-taller.bib}

% --- Metadatos del taller (vienen del YAML del .qmd vía pandoc) ---
\newcommand{\tallertitulo}{Taller 2 --- Componentes principales y
agrupamiento de países (EIA)}
\newcommand{\tallersubtitulo}{Análisis Estadístico de Datos}
\newcommand{\tallerfecha}{Bogotá D.C., 7 de octubre de 2026}
\newcommand{\tallerresumen}{Se estudian 111 países descritos por 57
variables de petróleo y electricidad de la U.S. Energy Information
Administration (EIA). Primero se eligen cinco variables, una por
dimensión y con poca redundancia entre sí (producción de crudo, consumo
de petróleo, exportación de crudo, generación eléctrica y generación
renovable), se estiman su vector de medias y su matriz de covarianzas y
se describen de forma univariada y bivariada, encontrando distribuciones
muy asimétricas dominadas por unos pocos países y dos bloques casi
independientes, el de la oferta petrolera y el del tamaño energético.
Luego se ejecuta un análisis de componentes principales (ACP) normado y
otro sin normalizar, se eligen cuatro componentes por el criterio del 80
\% de varianza acumulada y se interpretan las variables y los países en
los planos (1,2) y (2,3), donde el primer eje resulta ser el tamaño
energético del país y los siguientes oponen productores de petróleo a
productores de electricidad renovable. Al retirar a Estados Unidos,
China, Arabia Saudita y Rusia el primer eje pasa de explicar el 62 \% al
49 \% y aparecen dos brazos de países que sugieren tres grupos.
Finalmente se eligen tres grupos por la reducción de varianza explicada
y se comparan k-medias y el agrupamiento jerárquico de enlace promedio,
que coinciden en 101 de los 107 países y en separar a India y Japón,
pero difieren en cómo tratan a ocho economías grandes, mientras que las
cajas por grupo muestran que la partición sigue al consumo, generación
eléctrica y no a la producción ni a la exportación de crudo.}
\newcommand{\tallerpalabras}{Análisis de componentes principales;
Círculo de correlaciones; Criterio de Kaiser; K-medias; Agrupamiento
jerárquico; Silueta; EIA}
\input{common/meta-taller-es}

\begin{document}
\input{common/frontmatter-taller}

\chapter{Los datos}\label{datos}

La base \texttt{EIA.csv} reúne, para 111 países, 57 variables de
producción, consumo, comercio y capacidad instalada de petróleo y
electricidad publicadas por la U.S. Energy Information Administration
\parencite{eia-datos}, donde cada unidad estadística (UE) es un país.
Las unidades vienen en el sufijo de cada nombre, es decir, TBPD son
miles de barriles por día, MMTCD millones de toneladas métricas de
CO\textsubscript{2}, BKWH miles de millones de kWh, MK millones de kW y
BB miles de millones de barriles. Todo el análisis se realiza en R
\parencite{r-core} con gráficos de \texttt{ggplot2} \parencite{ggplot2},
\texttt{GGally} \parencite{ggally} y \texttt{factoextra}
\parencite{factoextra}, mientras que la documentación de cada función
utilizada se reúne en el Apéndice \ref{documentacion-r}.

La base se descarga de Kaggle o, en caso de no haber conexión, se busca
una copia local.

\Needspace{2.28in}

\begin{Shaded}
\begin{Highlighting}[]
\NormalTok{url }\OtherTok{\textless{}{-}} \FunctionTok{paste0}\NormalTok{(}\StringTok{"https://www.kaggle.com/api/v1/datasets/download/"}\NormalTok{,}
              \StringTok{"overhcv/eia{-}united{-}states"}\NormalTok{)}
\NormalTok{zip }\OtherTok{\textless{}{-}} \FunctionTok{file.path}\NormalTok{(}\FunctionTok{tempdir}\NormalTok{(), }\StringTok{"kaggle.zip"}\NormalTok{)}
\NormalTok{ruta }\OtherTok{\textless{}{-}} \FunctionTok{tryCatch}\NormalTok{(\{}
  \FunctionTok{download.file}\NormalTok{(url, zip, }\AttributeTok{mode =} \StringTok{"wb"}\NormalTok{, }\AttributeTok{quiet =} \ConstantTok{TRUE}\NormalTok{)}
  \FunctionTok{unzip}\NormalTok{(zip, }\AttributeTok{files =} \StringTok{"EIA.csv"}\NormalTok{, }\AttributeTok{exdir =} \FunctionTok{tempdir}\NormalTok{())}
\NormalTok{\}, }\AttributeTok{error =} \ControlFlowTok{function}\NormalTok{(e) \{}
  \FunctionTok{Filter}\NormalTok{(file.exists, }\FunctionTok{c}\NormalTok{(}\StringTok{"EIA.csv"}\NormalTok{, }\StringTok{"data/EIA.csv"}\NormalTok{,}
                        \StringTok{"../../../data/EIA.csv"}\NormalTok{))[}\DecValTok{1}\NormalTok{]}
\NormalTok{\})}
\NormalTok{datos }\OtherTok{\textless{}{-}} \FunctionTok{read.csv}\NormalTok{(ruta, }\AttributeTok{sep =} \StringTok{";"}\NormalTok{, }\AttributeTok{dec =} \StringTok{"."}\NormalTok{,}
                  \AttributeTok{stringsAsFactors =} \ConstantTok{FALSE}\NormalTok{)}
\end{Highlighting}
\end{Shaded}

\section{Estructura de las variables}\label{estructura}

Los nombres originales llegan a tener 60 caracteres, por lo que se
reemplazan por abreviaturas con la función \texttt{abreviar()} del
archivo \texttt{R/abreviaturas.R}, que se carga al inicio y cuyo código
y tabla de equivalencias decidimos mover al Apéndice \ref{abreviaturas}
al ser tan extensa la tabla. Por ejemplo \texttt{PetLiq\_cons} es el
consumo de petróleo y otros líquidos y \texttt{Renov\_gen} la generación
neta de electricidad renovable. Antes de calcular nada vale la pena ver
cómo están organizadas las 57 variables pues no son 57 mediciones
independientes sino un árbol en el que cada total se vuelve a medir en
sus partes.

\Needspace{4.20in}

\includegraphics[width=5.2in,height=3.62in]{sda-lab-exploracion_files/figure-latex/mermaid-figure-1.png}

Estructura de las 57 variables por sector energético, bloque y unidad:
Cada total se repite en sus partes y cada fuente eléctrica se mide dos
veces, por lo que la base es redundante desde su diseño.

El petróleo se describe en tres niveles (el crudo que se extrae y se
comercia, los totales de petróleo, refinados y emisiones, y los
productos uno por uno), mientras que la electricidad se describe como
sistema (generación, consumo, capacidad, comercio y pérdidas) y por
fuente, con cada una de las 11 fuentes medida en capacidad instalada y
en generación.

\section{Diagnóstico de la base}\label{diagnostico}

La matriz \texttt{num} guarda las 57 variables numéricas con el país
como nombre de fila y sobre ella se cuentan los faltantes y las
variables casi repetidas.

\Needspace{3.23in}

\begin{Shaded}
\begin{Highlighting}[]
\NormalTok{num }\OtherTok{\textless{}{-}} \FunctionTok{numericas}\NormalTok{(datos)}
\NormalTok{r }\OtherTok{\textless{}{-}} \FunctionTok{cor}\NormalTok{(num)}
\NormalTok{knitr}\SpecialCharTok{::}\FunctionTok{kable}\NormalTok{(}\FunctionTok{data.frame}\NormalTok{(}\FunctionTok{nrow}\NormalTok{(num), }\FunctionTok{ncol}\NormalTok{(num), }\FunctionTok{sum}\NormalTok{(}\FunctionTok{is.na}\NormalTok{(num)),}
                        \FunctionTok{anyDuplicated}\NormalTok{(}\FunctionTok{names}\NormalTok{(num)),}
                        \FunctionTok{sum}\NormalTok{(r[}\FunctionTok{upper.tri}\NormalTok{(r)] }\SpecialCharTok{\textgreater{}} \FloatTok{0.99}\NormalTok{),}
                        \FunctionTok{sum}\NormalTok{(}\FunctionTok{duplicated}\NormalTok{(}\FunctionTok{as.list}\NormalTok{(num)))),}
             \AttributeTok{col.names =} \FunctionTok{c}\NormalTok{(}\StringTok{"Países"}\NormalTok{, }\StringTok{"Variables"}\NormalTok{, }\StringTok{"Faltantes"}\NormalTok{,}
                           \StringTok{"Nombres repetidos"}\NormalTok{, }\StringTok{"Pares r \textgreater{} 0.99"}\NormalTok{,}
                           \StringTok{"Columnas repetidas"}\NormalTok{),}
             \AttributeTok{caption =} \StringTok{"Diagnóstico de la base."}\NormalTok{)}
\end{Highlighting}
\end{Shaded}

\Needspace{1.36in}

\begin{longtable}[]{@{}
  >{\raggedleft\arraybackslash}p{(\linewidth - 10\tabcolsep) * \real{0.0886}}
  >{\raggedleft\arraybackslash}p{(\linewidth - 10\tabcolsep) * \real{0.1266}}
  >{\raggedleft\arraybackslash}p{(\linewidth - 10\tabcolsep) * \real{0.1266}}
  >{\raggedleft\arraybackslash}p{(\linewidth - 10\tabcolsep) * \real{0.2278}}
  >{\raggedleft\arraybackslash}p{(\linewidth - 10\tabcolsep) * \real{0.1899}}
  >{\raggedleft\arraybackslash}p{(\linewidth - 10\tabcolsep) * \real{0.2405}}@{}}
\caption{Diagnóstico de la base.}\tabularnewline
\toprule\noalign{}
\begin{minipage}[b]{\linewidth}\raggedleft
Países
\end{minipage} & \begin{minipage}[b]{\linewidth}\raggedleft
Variables
\end{minipage} & \begin{minipage}[b]{\linewidth}\raggedleft
Faltantes
\end{minipage} & \begin{minipage}[b]{\linewidth}\raggedleft
Nombres repetidos
\end{minipage} & \begin{minipage}[b]{\linewidth}\raggedleft
Pares r \textgreater{} 0.99
\end{minipage} & \begin{minipage}[b]{\linewidth}\raggedleft
Columnas repetidas
\end{minipage} \\
\midrule\noalign{}
\endfirsthead
\toprule\noalign{}
\begin{minipage}[b]{\linewidth}\raggedleft
Países
\end{minipage} & \begin{minipage}[b]{\linewidth}\raggedleft
Variables
\end{minipage} & \begin{minipage}[b]{\linewidth}\raggedleft
Faltantes
\end{minipage} & \begin{minipage}[b]{\linewidth}\raggedleft
Nombres repetidos
\end{minipage} & \begin{minipage}[b]{\linewidth}\raggedleft
Pares r \textgreater{} 0.99
\end{minipage} & \begin{minipage}[b]{\linewidth}\raggedleft
Columnas repetidas
\end{minipage} \\
\midrule\noalign{}
\endhead
\bottomrule\noalign{}
\endlastfoot
111 & 57 & 0 & 0 & 18 & 1 \\
\end{longtable}

\Needspace{2.09in}

\begin{Shaded}
\begin{Highlighting}[]
\NormalTok{pares }\OtherTok{\textless{}{-}} \FunctionTok{which}\NormalTok{(}\FunctionTok{upper.tri}\NormalTok{(r) }\SpecialCharTok{\&}\NormalTok{ r }\SpecialCharTok{\textgreater{}} \FloatTok{0.99}\NormalTok{, }\AttributeTok{arr.ind =} \ConstantTok{TRUE}\NormalTok{)}
\FunctionTok{cat}\NormalTok{(}\FunctionTok{strwrap}\NormalTok{(}\FunctionTok{paste}\NormalTok{(}\FunctionTok{names}\NormalTok{(num)[pares[, }\DecValTok{1}\NormalTok{]], }\FunctionTok{names}\NormalTok{(num)[pares[, }\DecValTok{2}\NormalTok{]],}
                  \AttributeTok{sep =} \StringTok{"\textasciitilde{}"}\NormalTok{, }\AttributeTok{collapse =} \StringTok{", "}\NormalTok{), }\DecValTok{64}\NormalTok{), }\AttributeTok{sep =} \StringTok{"}\SpecialCharTok{\textbackslash{}n}\StringTok{"}\NormalTok{)}
\end{Highlighting}
\end{Shaded}

\begin{verbatim}
Gasolina_cons~Gasolina_prod, PetLiq_CO2~PetLiq_cons,
PetLiq_CO2~Refinados_cons, PetLiq_cons~Refinados_cons,
Gasolina_prod~GananciaRef, Elec_cap~Elec_cons,
Elec_cap~Elec_gen, Elec_cons~Elec_gen, Elec_cap~Fosil_cap,
Elec_cons~Fosil_cap, Elec_gen~Fosil_cap, Elec_cons~Fosil_gen,
Elec_gen~Fosil_gen, Fosil_cap~Fosil_gen,
Geotermica_cap~Geotermica_gen, Hidro_cap~Hidro_gen,
Nuclear_cap~Nuclear_gen, Renov_cap~Renov_gen
\end{verbatim}

\Needspace{2.47in}

\begin{codres}

\begin{Shaded}
\begin{Highlighting}[]
\FunctionTok{names}\NormalTok{(num)[}\FunctionTok{duplicated}\NormalTok{(}
  \FunctionTok{as.list}\NormalTok{(num))]}
\end{Highlighting}
\end{Shaded}

\codressep

\begin{verbatim}
[1] "Refinados_cons"
\end{verbatim}

\end{codres}

\begin{codres}

\begin{Shaded}
\begin{Highlighting}[]
\FunctionTok{max}\NormalTok{(}\FunctionTok{abs}\NormalTok{(num}\SpecialCharTok{$}\NormalTok{Renov\_gen }\SpecialCharTok{{-}}
\NormalTok{  num}\SpecialCharTok{$}\NormalTok{Hidro\_gen }\SpecialCharTok{{-}}
\NormalTok{  num}\SpecialCharTok{$}\NormalTok{NoHidro\_gen))}
\end{Highlighting}
\end{Shaded}

\codressep

\begin{verbatim}
[1] 1
\end{verbatim}

\end{codres}

\textbf{¿Está lista la base para el análisis?} Sí pues no hay datos
faltantes y las 57 abreviaturas son distintas entre sí, aunque la base
es muy redundante y esto hay que tenerlo presente desde ya:

\begin{enumerate}
\def\labelenumi{\alph{enumi}.}
\tightlist
\item
  \textbf{Lectura local.} La columna repetida, \texttt{Refinados\_cons},
  es el consumo de productos refinados, que resulta idéntico al consumo
  de petróleo y otros líquidos, mientras que la generación renovable
  coincide con la suma de la hidroeléctrica y la no hidroeléctrica salvo
  por una unidad de redondeo. Es decir que hay variables que no aportan
  nada que no esté ya en otra.
\item
  \textbf{Lectura global.} Los 18 pares con correlación mayor a 0.99 no
  son casualidad sino consecuencia del árbol de la figura anterior ya
  que cinco son la capacidad y la generación de una misma fuente (fósil,
  hidro, nuclear, geotérmica y renovable), ocho son el sistema eléctrico
  consigo mismo y con su parte fósil, tres son el consumo de petróleo,
  el de refinados y las emisiones entre sí, y dos son la producción de
  gasolina con su consumo y con la ganancia de refinación. Esto pesa en
  el ACP porque cada copia de una misma información suma varianza al
  mismo eje y en el agrupamiento porque las distancias entre países
  cuentan varias veces lo mismo.
\end{enumerate}

\chapter{Normal multivariada}\label{normal}

\section{Pregunta 1 --- Vector de medias y matriz de
covarianzas}\label{p1}

\textbf{¿Qué cinco variables y por qué?} Se eligió una variable por cada
dimensión de la base, buscando que entre ellas hubiera poca redundancia,
ya que dos variables casi iguales describen dos veces lo mismo y no
dicen nada nuevo de los países:

\begin{enumerate}
\def\labelenumi{\alph{enumi}.}
\tightlist
\item
  \textbf{Oferta petrolera.} Se mide con la producción de crudo
  (\texttt{Crudo\_prod}, TBPD), que indica cuánto petróleo extrae el
  país.
\item
  \textbf{Demanda petrolera.} Se mide con el consumo de petróleo y otros
  líquidos (\texttt{PetLiq\_cons}, TBPD), que refleja el tamaño de la
  economía que lo usa.
\item
  \textbf{Comercio petrolero.} Se mide con la exportación de crudo
  (\texttt{Crudo\_exp}, TBPD), que distingue al país que vende su
  petróleo del que lo consume.
\item
  \textbf{Escala eléctrica.} Se mide con la generación neta de
  electricidad (\texttt{Elec\_gen}, BKWH), que es el tamaño del sistema
  eléctrico.
\item
  \textbf{Transición energética.} Se mide con la generación neta
  renovable (\texttt{Renov\_gen}, BKWH), que indica cuánto de esa
  electricidad viene de fuentes renovables.
\end{enumerate}

Intentamos respaldar las elecciones con la matriz de correlación de las
candidatas, que incluye las emisiones de CO\textsubscript{2} por
petróleo (\texttt{PetLiq\_CO2}), la variable que a primera vista
representaría la dimensión ambiental.

\Needspace{3.04in}

\begin{Shaded}
\begin{Highlighting}[]
\NormalTok{candidatas }\OtherTok{\textless{}{-}} \FunctionTok{c}\NormalTok{(}\StringTok{"Crudo\_prod"}\NormalTok{, }\StringTok{"PetLiq\_cons"}\NormalTok{, }\StringTok{"PetLiq\_CO2"}\NormalTok{,}
                \StringTok{"Crudo\_exp"}\NormalTok{, }\StringTok{"Elec\_gen"}\NormalTok{, }\StringTok{"Renov\_gen"}\NormalTok{)}
\NormalTok{knitr}\SpecialCharTok{::}\FunctionTok{kable}\NormalTok{(}\FunctionTok{round}\NormalTok{(}\FunctionTok{cor}\NormalTok{(num[candidatas]), }\DecValTok{2}\NormalTok{),}
             \AttributeTok{caption =} \StringTok{"Correlación de Pearson de las candidatas."}\NormalTok{)}
\end{Highlighting}
\end{Shaded}

\Needspace{2.31in}

\begin{longtable}[]{@{}
  >{\raggedright\arraybackslash}p{(\linewidth - 12\tabcolsep) * \real{0.1600}}
  >{\raggedleft\arraybackslash}p{(\linewidth - 12\tabcolsep) * \real{0.1467}}
  >{\raggedleft\arraybackslash}p{(\linewidth - 12\tabcolsep) * \real{0.1600}}
  >{\raggedleft\arraybackslash}p{(\linewidth - 12\tabcolsep) * \real{0.1467}}
  >{\raggedleft\arraybackslash}p{(\linewidth - 12\tabcolsep) * \real{0.1333}}
  >{\raggedleft\arraybackslash}p{(\linewidth - 12\tabcolsep) * \real{0.1200}}
  >{\raggedleft\arraybackslash}p{(\linewidth - 12\tabcolsep) * \real{0.1333}}@{}}
\caption{Correlación de Pearson de las candidatas.}\tabularnewline
\toprule\noalign{}
\begin{minipage}[b]{\linewidth}\raggedright
\end{minipage} & \begin{minipage}[b]{\linewidth}\raggedleft
Crudo\_prod
\end{minipage} & \begin{minipage}[b]{\linewidth}\raggedleft
PetLiq\_cons
\end{minipage} & \begin{minipage}[b]{\linewidth}\raggedleft
PetLiq\_CO2
\end{minipage} & \begin{minipage}[b]{\linewidth}\raggedleft
Crudo\_exp
\end{minipage} & \begin{minipage}[b]{\linewidth}\raggedleft
Elec\_gen
\end{minipage} & \begin{minipage}[b]{\linewidth}\raggedleft
Renov\_gen
\end{minipage} \\
\midrule\noalign{}
\endfirsthead
\toprule\noalign{}
\begin{minipage}[b]{\linewidth}\raggedright
\end{minipage} & \begin{minipage}[b]{\linewidth}\raggedleft
Crudo\_prod
\end{minipage} & \begin{minipage}[b]{\linewidth}\raggedleft
PetLiq\_cons
\end{minipage} & \begin{minipage}[b]{\linewidth}\raggedleft
PetLiq\_CO2
\end{minipage} & \begin{minipage}[b]{\linewidth}\raggedleft
Crudo\_exp
\end{minipage} & \begin{minipage}[b]{\linewidth}\raggedleft
Elec\_gen
\end{minipage} & \begin{minipage}[b]{\linewidth}\raggedleft
Renov\_gen
\end{minipage} \\
\midrule\noalign{}
\endhead
\bottomrule\noalign{}
\endlastfoot
Crudo\_prod & 1.00 & 0.64 & 0.64 & 0.83 & 0.52 & 0.44 \\
PetLiq\_cons & 0.64 & 1.00 & 1.00 & 0.17 & 0.92 & 0.77 \\
PetLiq\_CO2 & 0.64 & 1.00 & 1.00 & 0.17 & 0.92 & 0.77 \\
Crudo\_exp & 0.83 & 0.17 & 0.17 & 1.00 & 0.07 & 0.06 \\
Elec\_gen & 0.52 & 0.92 & 0.92 & 0.07 & 1.00 & 0.92 \\
Renov\_gen & 0.44 & 0.77 & 0.77 & 0.06 & 0.92 & 1.00 \\
\end{longtable}

Las emisiones de CO\textsubscript{2} se descartan porque tienen
correlación de 1.00 con el consumo, lo cual indica que se derivan del
petróleo consumido y que incluirlas sería contar dos veces la demanda.
En su lugar entra la exportación de crudo, que casi no se asocia con el
bloque de tamaño (0.17 con el consumo, 0.07 con la generación y 0.06 con
la renovable) y sí con la producción (0.83), aunque sin repetirla pues
como se ve en la Pregunta 2 hay países que producen sin exportar. Las
otras tres variables tienen entre sí correlaciones altas (de 0.77 a
0.92) pero describen cosas distintas ya que el consumo de petróleo es
demanda de combustibles y la generación es electricidad, mientras que la
renovable es solo una parte de esa electricidad cuyo peso cambia mucho
de un país a otro.

Sobre las \(n = 111\) filas se estiman el vector de medias y la matriz
de varianzas y covarianzas muestrales

\[
\bar{\mathbf{x}} = \frac{1}{n}\sum_{i=1}^{n}\mathbf{x}_i,
\qquad
S = \frac{1}{n-1}\sum_{i=1}^{n}
    (\mathbf{x}_i - \bar{\mathbf{x}})(\mathbf{x}_i - \bar{\mathbf{x}})^\top,
\]

que son los estimadores insesgados de \(\boldsymbol{\mu}\) y \(\Sigma\)
en el modelo normal multivariado explicado en clase (el de máxima
verosimilitud de \(\Sigma\) que se divide por \(n\) en vez de
\(n - 1\)).

\Needspace{3.99in}

\begin{Shaded}
\begin{Highlighting}[]
\NormalTok{cinco }\OtherTok{\textless{}{-}} \FunctionTok{c}\NormalTok{(}\StringTok{"Crudo\_prod"}\NormalTok{, }\StringTok{"PetLiq\_cons"}\NormalTok{, }\StringTok{"Crudo\_exp"}\NormalTok{, }\StringTok{"Elec\_gen"}\NormalTok{,}
           \StringTok{"Renov\_gen"}\NormalTok{)}
\NormalTok{X }\OtherTok{\textless{}{-}}\NormalTok{ num[cinco]}
\NormalTok{knitr}\SpecialCharTok{::}\FunctionTok{kable}\NormalTok{(}\FunctionTok{data.frame}\NormalTok{(}\AttributeTok{Variable =}\NormalTok{ cinco,}
                        \AttributeTok{Unidad =} \FunctionTok{c}\NormalTok{(}\StringTok{"TBPD"}\NormalTok{, }\StringTok{"TBPD"}\NormalTok{, }\StringTok{"TBPD"}\NormalTok{, }\StringTok{"BKWH"}\NormalTok{,}
                                   \StringTok{"BKWH"}\NormalTok{),}
                        \AttributeTok{Media =} \FunctionTok{round}\NormalTok{(}\FunctionTok{colMeans}\NormalTok{(X), }\DecValTok{2}\NormalTok{),}
                        \AttributeTok{Mediana =} \FunctionTok{sapply}\NormalTok{(X, median),}
                        \AttributeTok{row.names =} \ConstantTok{NULL}\NormalTok{),}
             \AttributeTok{caption =} \StringTok{"Vector de medias (y mediana)."}\NormalTok{)}
\end{Highlighting}
\end{Shaded}

\Needspace{2.12in}

\begin{longtable}[]{@{}llrr@{}}
\caption{Vector de medias (y mediana).}\tabularnewline
\toprule\noalign{}
Variable & Unidad & Media & Mediana \\
\midrule\noalign{}
\endfirsthead
\toprule\noalign{}
Variable & Unidad & Media & Mediana \\
\midrule\noalign{}
\endhead
\bottomrule\noalign{}
\endlastfoot
Crudo\_prod & TBPD & 702.06 & 38.0 \\
PetLiq\_cons & TBPD & 811.40 & 202.0 \\
Crudo\_exp & TBPD & 361.56 & 8.5 \\
Elec\_gen & BKWH & 198.29 & 43.0 \\
Renov\_gen & BKWH & 44.77 & 8.3 \\
\end{longtable}

\Needspace{2.47in}

\begin{Shaded}
\begin{Highlighting}[]
\NormalTok{knitr}\SpecialCharTok{::}\FunctionTok{kable}\NormalTok{(}\FunctionTok{round}\NormalTok{(}\FunctionTok{cov}\NormalTok{(X), }\DecValTok{1}\NormalTok{),}
             \AttributeTok{caption =} \StringTok{"Matriz de varianzas y covarianzas S."}\NormalTok{)}
\end{Highlighting}
\end{Shaded}

\Needspace{2.12in}

\begin{longtable}[]{@{}lrrrrr@{}}
\caption{Matriz de varianzas y covarianzas S.}\tabularnewline
\toprule\noalign{}
& Crudo\_prod & PetLiq\_cons & Crudo\_exp & Elec\_gen & Renov\_gen \\
\midrule\noalign{}
\endfirsthead
\toprule\noalign{}
& Crudo\_prod & PetLiq\_cons & Crudo\_exp & Elec\_gen & Renov\_gen \\
\midrule\noalign{}
\endhead
\bottomrule\noalign{}
\endlastfoot
Crudo\_prod & 3032022.8 & 2471112.6 & 1374119.6 & 594768.9 & 110047.2 \\
PetLiq\_cons & 2471112.6 & 4858882.6 & 358231.7 & 1332865.6 &
245051.1 \\
Crudo\_exp & 1374119.6 & 358231.7 & 907732.0 & 41014.9 & 8667.9 \\
Elec\_gen & 594768.9 & 1332865.6 & 41014.9 & 434095.0 & 87623.9 \\
Renov\_gen & 110047.2 & 245051.1 & 8667.9 & 87623.9 & 20864.5 \\
\end{longtable}

\textbf{Vector de medias.} Se estima
\(\bar{\mathbf{x}} = (702.06,\ 811.40,\ 361.56,\ 198.29,\ 44.77)^\top\),
es decir que un país promedio produce 702 mil barriles diarios de crudo,
consume 811 mil barriles diarios de petróleo, exporta 362 mil barriles
diarios de crudo y genera 198 mil millones de kWh, de los cuales 45 son
renovables.

\textbf{Matriz de covarianzas.} Todas las covarianzas son positivas, lo
que significa que las cinco variables crecen juntas y que un país que
consume más petróleo también tiende a generar más electricidad. Sus
magnitudes en cambio no se pueden comparar entre sí puesto que cada
entrada lleva el producto de las unidades de sus dos variables (por
ejemplo, la varianza del consumo está en TBPD² y su covarianza con la
generación en TBPD·BKWH), por lo que la fuerza de cada asociación se lee
en la Pregunta 2 con la correlación.

\begin{enumerate}
\def\labelenumi{\alph{enumi}.}
\tightlist
\item
  \textbf{Lectura local.} Como tal el país promedio no existe pues la
  media de la exportación de crudo (361.56 TBPD) es 42 veces su mediana
  (8.5 TBPD) y la de la producción (702.06 TBPD) más de 18 veces la suya
  (38 TBPD), de modo que cada componente de \(\bar{\mathbf{x}}\) la
  fijan unos pocos países y no el país típico.
\item
  \textbf{Lectura global.} Bajo normalidad multivariada
  \(\bar{\mathbf{x}}\) y \(S\) resumirían toda la distribución, pero
  estos datos están lejos de ser normales y \(S\) mezcla dos bloques, el
  petrolero (producción y exportación) y el de tamaño (consumo,
  generación y renovable), que casi no se mueven juntos. Esto puede
  anticiparnos que una sola dirección no baste para describir a los
  países y que el ACP necesite al menos dos ejes.
\end{enumerate}

Para ver cuánto de las correlaciones de las candidatas se debe al tamaño
se construyen dos razones internas a la base que no dependen de lo
grande que sea el país, la participación renovable en la generación
(\texttt{Renov\_gen\ /\ Elec\_gen}) y la autosuficiencia petrolera
(\texttt{Crudo\_prod\ /\ PetLiq\_cons}), cuyos denominadores no valen
cero en ningún país, para compararlas con la correlación de las mismas
variables en términos absolutos.

\Needspace{4.37in}

\begin{Shaded}
\begin{Highlighting}[]
\FunctionTok{stopifnot}\NormalTok{(}\FunctionTok{all}\NormalTok{(num}\SpecialCharTok{$}\NormalTok{Elec\_gen }\SpecialCharTok{\textgreater{}} \DecValTok{0}\NormalTok{), }\FunctionTok{all}\NormalTok{(num}\SpecialCharTok{$}\NormalTok{PetLiq\_cons }\SpecialCharTok{\textgreater{}} \DecValTok{0}\NormalTok{))}
\NormalTok{renovable }\OtherTok{\textless{}{-}}\NormalTok{ num}\SpecialCharTok{$}\NormalTok{Renov\_gen }\SpecialCharTok{/}\NormalTok{ num}\SpecialCharTok{$}\NormalTok{Elec\_gen}
\NormalTok{autosuficiencia }\OtherTok{\textless{}{-}}\NormalTok{ num}\SpecialCharTok{$}\NormalTok{Crudo\_prod }\SpecialCharTok{/}\NormalTok{ num}\SpecialCharTok{$}\NormalTok{PetLiq\_cons}
\NormalTok{knitr}\SpecialCharTok{::}\FunctionTok{kable}\NormalTok{(}\FunctionTok{data.frame}\NormalTok{(}
  \AttributeTok{par =} \FunctionTok{c}\NormalTok{(}\StringTok{"Renovable y generación"}\NormalTok{, }\StringTok{"Producción y consumo"}\NormalTok{,}
          \StringTok{"Renovable y producción"}\NormalTok{),}
  \AttributeTok{absolutas =} \FunctionTok{c}\NormalTok{(}\FunctionTok{cor}\NormalTok{(num}\SpecialCharTok{$}\NormalTok{Renov\_gen, num}\SpecialCharTok{$}\NormalTok{Elec\_gen),}
                \FunctionTok{cor}\NormalTok{(num}\SpecialCharTok{$}\NormalTok{Crudo\_prod, num}\SpecialCharTok{$}\NormalTok{PetLiq\_cons),}
                \FunctionTok{cor}\NormalTok{(num}\SpecialCharTok{$}\NormalTok{Renov\_gen, num}\SpecialCharTok{$}\NormalTok{Crudo\_prod)),}
  \AttributeTok{con\_razon =} \FunctionTok{c}\NormalTok{(}\FunctionTok{cor}\NormalTok{(renovable, num}\SpecialCharTok{$}\NormalTok{Elec\_gen),}
                \FunctionTok{cor}\NormalTok{(autosuficiencia, num}\SpecialCharTok{$}\NormalTok{PetLiq\_cons),}
                \FunctionTok{cor}\NormalTok{(renovable, autosuficiencia))),}
  \AttributeTok{digits =} \DecValTok{2}\NormalTok{,}
  \AttributeTok{caption =} \StringTok{"Correlación en términos absolutos y con razones."}\NormalTok{)}
\end{Highlighting}
\end{Shaded}

\Needspace{1.74in}

\begin{longtable}[]{@{}lrr@{}}
\caption{Correlación en términos absolutos y con
razones.}\tabularnewline
\toprule\noalign{}
par & absolutas & con\_razon \\
\midrule\noalign{}
\endfirsthead
\toprule\noalign{}
par & absolutas & con\_razon \\
\midrule\noalign{}
\endhead
\bottomrule\noalign{}
\endlastfoot
Renovable y generación & 0.92 & -0.10 \\
Producción y consumo & 0.64 & -0.07 \\
Renovable y producción & 0.44 & -0.06 \\
\end{longtable}

Del mismo modo se divide el CO\textsubscript{2} entre el consumo de
petróleo para ver si las emisiones se miden por separado o se calculan a
partir del consumo.

\Needspace{2.85in}

\begin{Shaded}
\begin{Highlighting}[]
\NormalTok{co2 }\OtherTok{\textless{}{-}}\NormalTok{ num}\SpecialCharTok{$}\NormalTok{PetLiq\_CO2 }\SpecialCharTok{/}\NormalTok{ num}\SpecialCharTok{$}\NormalTok{PetLiq\_cons}
\NormalTok{knitr}\SpecialCharTok{::}\FunctionTok{kable}\NormalTok{(}\FunctionTok{data.frame}\NormalTok{(}
  \AttributeTok{r =} \FunctionTok{cor}\NormalTok{(num}\SpecialCharTok{$}\NormalTok{PetLiq\_CO2, num}\SpecialCharTok{$}\NormalTok{PetLiq\_cons),}
  \AttributeTok{minimo =} \FunctionTok{min}\NormalTok{(co2), }\AttributeTok{mediana =} \FunctionTok{median}\NormalTok{(co2), }\AttributeTok{maximo =} \FunctionTok{max}\NormalTok{(co2),}
  \AttributeTok{cv =} \FunctionTok{sd}\NormalTok{(co2) }\SpecialCharTok{/} \FunctionTok{mean}\NormalTok{(co2),}
  \AttributeTok{t\_por\_barril =} \FunctionTok{median}\NormalTok{(co2) }\SpecialCharTok{*} \FloatTok{1e6} \SpecialCharTok{/} \FloatTok{365e3}\NormalTok{),}
  \AttributeTok{digits =} \DecValTok{3}\NormalTok{,}
  \AttributeTok{caption =} \StringTok{"Emisiones de CO2 por unidad de petróleo consumido."}\NormalTok{)}
\end{Highlighting}
\end{Shaded}

\Needspace{1.36in}

\begin{longtable}[]{@{}rrrrrr@{}}
\caption{Emisiones de CO2 por unidad de petróleo
consumido.}\tabularnewline
\toprule\noalign{}
r & minimo & mediana & maximo & cv & t\_por\_barril \\
\midrule\noalign{}
\endfirsthead
\toprule\noalign{}
r & minimo & mediana & maximo & cv & t\_por\_barril \\
\midrule\noalign{}
\endhead
\bottomrule\noalign{}
\endlastfoot
0.999 & 0.095 & 0.137 & 0.16 & 0.088 & 0.374 \\
\end{longtable}

\begin{notas}

\textbf{Interpretaciones.}

\begin{enumerate}
\def\labelenumi{\alph{enumi}.}
\tightlist
\item
  \textbf{La elección es del grupo.} Con las cinco variables se cubren
  la oferta, la demanda y el comercio de petróleo, la escala eléctrica y
  la transición, y la única redundancia evidente (consumo y
  CO\textsubscript{2}, \(r = 1.00\)) quedó por fuera.
\item
  \textbf{El costo de la elección.} La producción y la exportación de
  crudo comparten el 69 \% de su variación (\(0.83^2\)) así que el
  bloque petrolero pesa doble en las preguntas 2 y 11, del mismo modo
  que el de tamaño pesa triple.
\end{enumerate}

\textbf{Decisiones.}

\begin{itemize}
\item
  \textbf{Se mantienen las variables en términos absolutos.} El
  enunciado pide caracterizar a los países «en términos de importación,
  exportación, consumo, producción» y la base no trae población, de modo
  que pasar a valores por habitante exigiría traer datos de otra fuente.
  Las razones sí muestran cuánto de la correlación es tamaño pues la
  renovable y la generación pasan de 0.92 a \(-0.10\) cuando la
  renovable se toma como participación, la producción y el consumo de
  0.64 a \(-0.07\) cuando la producción se divide por el consumo y la
  renovable y la producción de 0.44 a \(-0.06\) cuando se comparan las
  dos razones, es decir que casi toda la asociación entre las variables
  absolutas la pone el tamaño del país y que, una vez descontado, el
  perfil renovable y el petrolero son casi independientes. Por esto el
  primer eje del ACP puede resultar más tamaño que perfil energético, lo
  que conviene tener presente desde ya.

  \emph{Respaldo.} En la sesión 5 se advirtió que el agrupamiento queda
  limitado a las variables que se tienen, de modo que lo que no se mide
  (como la población) no lo recoge ningún método, mientras que en la
  sesión 2 se vio que las variables de razón admiten cocientes entre sí,
  que es lo que permite construir participaciones sin salir de la base.
\item
  \textbf{La dimensión ambiental queda implícita en el consumo.} La
  correlación entre el CO\textsubscript{2} y el consumo de petróleo es
  0.999 y el cociente entre ambos casi no cambia de un país a otro (de
  0.095 a 0.16 con mediana 0.137 y coeficiente de variación de 0.088),
  lo cual indica que las emisiones no se miden por separado sino que se
  calculan multiplicando el consumo por un factor de emisión que solo
  varía con la mezcla de productos de cada país. Si la cifra es anual la
  mediana equivale a 0.374 toneladas de CO\textsubscript{2} por barril,
  del orden de lo que emite un barril de combustible al quemarse. Por
  esto incluir el CO\textsubscript{2} sería repetir la demanda en otra
  unidad y no agregar una dimensión nueva.

  \emph{Respaldo.} La EIA calcula las emisiones de CO\textsubscript{2}
  del petróleo aplicando coeficientes de emisión al consumo de cada
  producto y sumándolas \parencite{eia-notas}, es decir que no las mide.
\end{itemize}

\end{notas}

\chapter{Análisis de componentes principales}\label{acp}

Las preguntas 2 a 11 siguen la ruta del diagrama, consideramos que
conviene tenerlo a la mano porque el número de países y lo que se busca
cambian por el camino. Primero se describen las cinco variables
(Pregunta 2) y luego se resume la base completa con el ACP, primero con
los 111 países y después sin los cuatro más extremos, para terminar
agrupando a los 107 restantes.

\Needspace{4.30in}

\includegraphics[width=5.2in,height=3.73in]{sda-lab-exploracion_files/figure-latex/mermaid-figure-5.png}

Ruta del análisis de las preguntas 3 a 11: Cada etapa cambia los países
o la pregunta, y las conclusiones de una justifican la siguiente.

\section{Pregunta 2 --- Análisis descriptivo univariado y
bivariado}\label{p2}

\Needspace{3.04in}

\begin{Shaded}
\begin{Highlighting}[]
\NormalTok{resumen }\OtherTok{\textless{}{-}} \FunctionTok{t}\NormalTok{(}\FunctionTok{sapply}\NormalTok{(X, }\ControlFlowTok{function}\NormalTok{(x)}
  \FunctionTok{c}\NormalTok{(}\AttributeTok{media =} \FunctionTok{mean}\NormalTok{(x), }\AttributeTok{mediana =} \FunctionTok{median}\NormalTok{(x), }\AttributeTok{s =} \FunctionTok{sd}\NormalTok{(x), }\AttributeTok{min =} \FunctionTok{min}\NormalTok{(x),}
    \AttributeTok{max =} \FunctionTok{max}\NormalTok{(x), }\AttributeTok{g1 =} \FunctionTok{g1}\NormalTok{(x), }\AttributeTok{ceros =} \FunctionTok{sum}\NormalTok{(x }\SpecialCharTok{==} \DecValTok{0}\NormalTok{))))}
\NormalTok{knitr}\SpecialCharTok{::}\FunctionTok{kable}\NormalTok{(}\FunctionTok{round}\NormalTok{(resumen, }\DecValTok{2}\NormalTok{),}
             \AttributeTok{caption =} \StringTok{"Resumen univariado de las cinco variables."}\NormalTok{)}
\end{Highlighting}
\end{Shaded}

\Needspace{2.12in}

\begin{longtable}[]{@{}lrrrrrrr@{}}
\caption{Resumen univariado de las cinco variables.}\tabularnewline
\toprule\noalign{}
& media & mediana & s & min & max & g1 & ceros \\
\midrule\noalign{}
\endfirsthead
\toprule\noalign{}
& media & mediana & s & min & max & g1 & ceros \\
\midrule\noalign{}
\endhead
\bottomrule\noalign{}
\endlastfoot
Crudo\_prod & 702.06 & 38.0 & 1741.27 & 0.0 & 10107 & 3.91 & 27 \\
PetLiq\_cons & 811.40 & 202.0 & 2204.29 & 15.0 & 19100 & 6.32 & 0 \\
Crudo\_exp & 361.56 & 8.5 & 952.75 & 0.0 & 7120 & 4.48 & 41 \\
Elec\_gen & 198.29 & 43.0 & 658.86 & 0.9 & 5400 & 6.46 & 0 \\
Renov\_gen & 44.77 & 8.3 & 144.45 & 0.0 & 1287 & 6.56 & 11 \\
\end{longtable}

\Needspace{2.09in}

\begin{Shaded}
\begin{Highlighting}[]
\NormalTok{knitr}\SpecialCharTok{::}\FunctionTok{kable}\NormalTok{(}\FunctionTok{sapply}\NormalTok{(X, }\ControlFlowTok{function}\NormalTok{(x) }\FunctionTok{rownames}\NormalTok{(X)[}\FunctionTok{order}\NormalTok{(}\SpecialCharTok{{-}}\NormalTok{x)][}\DecValTok{1}\SpecialCharTok{:}\DecValTok{3}\NormalTok{]),}
             \AttributeTok{caption =} \StringTok{"Los tres países con el valor más alto."}\NormalTok{)}
\end{Highlighting}
\end{Shaded}

\Needspace{1.74in}

\begin{longtable}[]{@{}
  >{\raggedright\arraybackslash}p{(\linewidth - 8\tabcolsep) * \real{0.1831}}
  >{\raggedright\arraybackslash}p{(\linewidth - 8\tabcolsep) * \real{0.1831}}
  >{\raggedright\arraybackslash}p{(\linewidth - 8\tabcolsep) * \real{0.2676}}
  >{\raggedright\arraybackslash}p{(\linewidth - 8\tabcolsep) * \real{0.1831}}
  >{\raggedright\arraybackslash}p{(\linewidth - 8\tabcolsep) * \real{0.1831}}@{}}
\caption{Los tres países con el valor más alto.}\tabularnewline
\toprule\noalign{}
\begin{minipage}[b]{\linewidth}\raggedright
Crudo\_prod
\end{minipage} & \begin{minipage}[b]{\linewidth}\raggedright
PetLiq\_cons
\end{minipage} & \begin{minipage}[b]{\linewidth}\raggedright
Crudo\_exp
\end{minipage} & \begin{minipage}[b]{\linewidth}\raggedright
Elec\_gen
\end{minipage} & \begin{minipage}[b]{\linewidth}\raggedright
Renov\_gen
\end{minipage} \\
\midrule\noalign{}
\endfirsthead
\toprule\noalign{}
\begin{minipage}[b]{\linewidth}\raggedright
Crudo\_prod
\end{minipage} & \begin{minipage}[b]{\linewidth}\raggedright
PetLiq\_cons
\end{minipage} & \begin{minipage}[b]{\linewidth}\raggedright
Crudo\_exp
\end{minipage} & \begin{minipage}[b]{\linewidth}\raggedright
Elec\_gen
\end{minipage} & \begin{minipage}[b]{\linewidth}\raggedright
Renov\_gen
\end{minipage} \\
\midrule\noalign{}
\endhead
\bottomrule\noalign{}
\endlastfoot
Russia & UnitedStates & SaudiArabia & China & China \\
SaudiArabia & China & Russia & UnitedStates & UnitedStates \\
UnitedStates & Japan & UnitedArabEmirates & India & Brazil \\
\end{longtable}

\Needspace{2.28in}

\begin{codres}

\begin{Shaded}
\begin{Highlighting}[]
\FunctionTok{sum}\NormalTok{(X}\SpecialCharTok{$}\NormalTok{Crudo\_prod }\SpecialCharTok{\textgreater{}} \DecValTok{0} \SpecialCharTok{\&}
\NormalTok{    X}\SpecialCharTok{$}\NormalTok{Crudo\_exp }\SpecialCharTok{==} \DecValTok{0}\NormalTok{)}
\end{Highlighting}
\end{Shaded}

\codressep

\begin{verbatim}
[1] 16
\end{verbatim}

\end{codres}

\begin{codres}

\begin{Shaded}
\begin{Highlighting}[]
\FunctionTok{rownames}\NormalTok{(X)[X}\SpecialCharTok{$}\NormalTok{Crudo\_exp }\SpecialCharTok{\textgreater{}} \DecValTok{0} \SpecialCharTok{\&}
\NormalTok{  X}\SpecialCharTok{$}\NormalTok{Crudo\_prod }\SpecialCharTok{==} \DecValTok{0}\NormalTok{]}
\end{Highlighting}
\end{Shaded}

\codressep

\begin{verbatim}
[1] "Senegal"   "Singapore"
\end{verbatim}

\end{codres}

\Needspace{7.06in}

\begin{Shaded}
\begin{Highlighting}[]
\FunctionTok{ggpairs}\NormalTok{(X, }\AttributeTok{upper =} \FunctionTok{list}\NormalTok{(}\AttributeTok{continuous =} \FunctionTok{wrap}\NormalTok{(}\StringTok{"cor"}\NormalTok{, }\AttributeTok{size =} \DecValTok{3}\NormalTok{)),}
        \AttributeTok{lower =} \FunctionTok{list}\NormalTok{(}\AttributeTok{continuous =} \FunctionTok{wrap}\NormalTok{(}\StringTok{"points"}\NormalTok{, }\AttributeTok{size =} \FloatTok{0.6}\NormalTok{,}
                                       \AttributeTok{alpha =} \FloatTok{0.5}\NormalTok{))) }\SpecialCharTok{+}
  \FunctionTok{theme\_minimal}\NormalTok{(}\AttributeTok{base\_size =} \DecValTok{8}\NormalTok{)}
\end{Highlighting}
\end{Shaded}

\begin{figure}[H]

{\centering \pandocbounded{\includegraphics[keepaspectratio]{sda-lab-exploracion_files/figure-latex/p2-pares-1.pdf}}

}

\caption{Matriz de dispersión de las cinco variables con densidades en
la diagonal y correlación de Pearson arriba: Casi toda la nube está
pegada al origen, unos pocos países definen las correlaciones y la
exportación de crudo solo se asocia con la producción.}

\end{figure}%

\Needspace{2.47in}

\begin{Shaded}
\begin{Highlighting}[]
\NormalTok{knitr}\SpecialCharTok{::}\FunctionTok{kable}\NormalTok{(}\FunctionTok{round}\NormalTok{(}\FunctionTok{cor}\NormalTok{(X, }\AttributeTok{method =} \StringTok{"spearman"}\NormalTok{), }\DecValTok{2}\NormalTok{),}
             \AttributeTok{caption =} \StringTok{"Correlación de Spearman."}\NormalTok{)}
\end{Highlighting}
\end{Shaded}

\Needspace{2.12in}

\begin{longtable}[]{@{}lrrrrr@{}}
\caption{Correlación de Spearman.}\tabularnewline
\toprule\noalign{}
& Crudo\_prod & PetLiq\_cons & Crudo\_exp & Elec\_gen & Renov\_gen \\
\midrule\noalign{}
\endfirsthead
\toprule\noalign{}
& Crudo\_prod & PetLiq\_cons & Crudo\_exp & Elec\_gen & Renov\_gen \\
\midrule\noalign{}
\endhead
\bottomrule\noalign{}
\endlastfoot
Crudo\_prod & 1.00 & 0.48 & 0.89 & 0.42 & 0.14 \\
PetLiq\_cons & 0.48 & 1.00 & 0.39 & 0.92 & 0.57 \\
Crudo\_exp & 0.89 & 0.39 & 1.00 & 0.28 & 0.01 \\
Elec\_gen & 0.42 & 0.92 & 0.28 & 1.00 & 0.67 \\
Renov\_gen & 0.14 & 0.57 & 0.01 & 0.67 & 1.00 \\
\end{longtable}

\Needspace{1.33in}

\begin{codres}

\begin{Shaded}
\begin{Highlighting}[]
\NormalTok{sin }\OtherTok{\textless{}{-}}\NormalTok{ X[}\FunctionTok{rownames}\NormalTok{(X) }\SpecialCharTok{!=} \StringTok{"China"}\NormalTok{, ]}
\FunctionTok{round}\NormalTok{(}\FunctionTok{cor}\NormalTok{(sin}\SpecialCharTok{$}\NormalTok{Elec\_gen,}
\NormalTok{  sin}\SpecialCharTok{$}\NormalTok{Renov\_gen), }\DecValTok{3}\NormalTok{)}
\end{Highlighting}
\end{Shaded}

\codressep

\begin{verbatim}
[1] 0.803
\end{verbatim}

\end{codres}

\textbf{Univariado.} Las cinco variables tienen asimetría positiva
extrema, con coeficientes \(g_1\) entre 3.91 (producción de crudo) y
6.56 (generación renovable), así como medias muy por encima de sus
medianas, por ejemplo 811.40 contra 202 TBPD en el consumo de petróleo y
361.56 contra 8.5 TBPD en la exportación de crudo. Las densidades de la
diagonal lo muestran como un pico pegado al cero con una cola larga a la
derecha formada por los países más grandes ya que Estados Unidos y China
ocupan los primeros puestos en consumo y generación, mientras que Rusia,
Arabia Saudita y Estados Unidos encabezan la producción de crudo y
Arabia Saudita, Rusia y Emiratos Árabes Unidos la exportación. Además la
producción vale cero en 27 países y la exportación en 41, de modo que 16
países producen crudo sin exportarlo y dos (Senegal y Singapur) lo
exportan sin producirlo, lo cual podría deberse a reexportaciones que la
base no permite verificar.

\textbf{Bivariado.} Las asociaciones son todas positivas pero se ordenan
en dos bloques:

\begin{enumerate}
\def\labelenumi{\alph{enumi}.}
\tightlist
\item
  \textbf{Bloque petrolero.} La producción y la exportación de crudo
  tienen correlación de Pearson \parencite{pearson1895} de 0.828 ya que
  solo exporta crudo quien lo produce, aunque no todo productor lo
  exporta.
\item
  \textbf{Bloque de tamaño.} El consumo de petróleo y la generación
  eléctrica se asocian con \(r = 0.918\) y la generación con la
  renovable con \(r = 0.921\), es decir que estas tres siguen al tamaño
  de la economía.
\item
  \textbf{Entre bloques.} La exportación casi no se asocia con el bloque
  de tamaño (0.171 con el consumo, 0.065 con la generación y 0.063 con
  la renovable), dado que los tres mayores consumidores (Estados Unidos,
  China y Japón) exportan poco o nada, mientras que la producción de
  crudo hace de puente entre ambos bloques (0.644 con el consumo y 0.518
  con la generación).
\end{enumerate}

Vale la pena matizar que varias de estas correlaciones dependen de los
países extremos pues la de Spearman \parencite{spearman1904}, que usa
rangos y no se deja arrastrar por los valores grandes, baja de 0.921 a
0.672 entre generación y renovable, de 0.770 a 0.569 entre consumo y
renovable y de 0.438 a 0.140 entre producción de crudo y renovable. Solo
el par consumo-generación (0.918) y el par producción-exportación (de
0.828 a 0.891) se mantienen con ambos coeficientes.

\begin{enumerate}
\def\labelenumi{\alph{enumi}.}
\tightlist
\item
  \textbf{Lectura local.} Un solo país mueve una correlación ya que sin
  China la de Pearson entre generación y renovable baja de 0.921 a 0.803
  porque China es a la vez el mayor generador y el mayor productor de
  renovables.
\item
  \textbf{Lectura global.} Con rangos la separación entre bloques es
  todavía más clara (0.010 entre exportación y renovable), de modo que
  las cinco variables viven en al menos dos dimensiones casi
  independientes, la de tamaño y la petrolera, que podrían ser las que
  el ACP de la Pregunta 5 encuentre en sus dos primeros ejes con las 57
  variables.
\end{enumerate}

\section{Pregunta 3 --- Contribuciones en el ACP normado y sin
normalizar}\label{p3}

El análisis de componentes principales (ACP) \parencite{pearson1901,
hotelling1933} busca combinaciones lineales de las variables, no
correlacionadas entre sí, que recojan la mayor varianza posible. Para
leerlo se usan tres medidas:

\begin{enumerate}
\def\labelenumi{\alph{enumi}.}
\tightlist
\item
  \textbf{ACP normado.} Es el ACP de las variables estandarizadas
  (\(z_{ij} = (x_{ij} - \bar{x}_j)/s_j\)), es decir, el de la matriz de
  correlación, de modo que cada variable aporta varianza 1 sin importar
  su unidad. El ACP sin normalizar usa la matriz de covarianzas en las
  unidades originales.
\item
  \textbf{Valor propio.} El valor propio \(\lambda_k\) es la varianza de
  la componente \(k\), de modo que \(\lambda_k / \sum_j \lambda_j\) es
  la proporción de varianza que explica. En el ACP normado los 57
  valores propios suman 57.
\item
  \textbf{Contribución.} La contribución de la variable \(j\) al eje
  \(k\) es \(100\,u_{jk}^2\), donde \(u_{jk}\) es su carga (su peso en
  el vector propio) y mide qué porcentaje del eje construye esa
  variable. La contribución a un plano pondera las de cada eje por su
  valor propio y se compara contra la referencia uniforme
  \(100/57 = 1.75\) \%.
\end{enumerate}

\Needspace{4.56in}

\begin{Shaded}
\begin{Highlighting}[]
\NormalTok{acp\_normado }\OtherTok{\textless{}{-}} \FunctionTok{prcomp}\NormalTok{(num, }\AttributeTok{scale. =} \ConstantTok{TRUE}\NormalTok{)}
\NormalTok{acp\_crudo }\OtherTok{\textless{}{-}} \FunctionTok{prcomp}\NormalTok{(num, }\AttributeTok{scale. =} \ConstantTok{FALSE}\NormalTok{)}
\NormalTok{aporte }\OtherTok{\textless{}{-}} \ControlFlowTok{function}\NormalTok{(acp, }\AttributeTok{n =} \DecValTok{10}\NormalTok{) \{}
\NormalTok{  a }\OtherTok{\textless{}{-}} \FunctionTok{facto\_summarize}\NormalTok{(acp, }\StringTok{"var"}\NormalTok{, }\AttributeTok{axes =} \FunctionTok{c}\NormalTok{(}\DecValTok{1}\NormalTok{, }\DecValTok{2}\NormalTok{))}
\NormalTok{  a }\OtherTok{\textless{}{-}}\NormalTok{ a[}\FunctionTok{order}\NormalTok{(}\SpecialCharTok{{-}}\NormalTok{a}\SpecialCharTok{$}\NormalTok{contrib), ]}
  \FunctionTok{data.frame}\NormalTok{(}\AttributeTok{variable =}\NormalTok{ a}\SpecialCharTok{$}\NormalTok{name,}
             \AttributeTok{contrib =} \FunctionTok{round}\NormalTok{(a}\SpecialCharTok{$}\NormalTok{contrib, }\DecValTok{2}\NormalTok{))[}\DecValTok{1}\SpecialCharTok{:}\NormalTok{n, ]}
\NormalTok{\}}
\NormalTok{knitr}\SpecialCharTok{::}\FunctionTok{kable}\NormalTok{(}\FunctionTok{cbind}\NormalTok{(}\FunctionTok{aporte}\NormalTok{(acp\_normado), }\FunctionTok{aporte}\NormalTok{(acp\_crudo)),}
             \AttributeTok{row.names =} \ConstantTok{FALSE}\NormalTok{,}
             \AttributeTok{col.names =} \FunctionTok{c}\NormalTok{(}\StringTok{"ACP normado"}\NormalTok{, }\StringTok{"\%"}\NormalTok{,}
                           \StringTok{"Sin normalizar"}\NormalTok{, }\StringTok{"\%"}\NormalTok{),}
             \AttributeTok{caption =} \StringTok{"Variables que más aportan al plano 1{-}2."}\NormalTok{)}
\end{Highlighting}
\end{Shaded}

\Needspace{3.07in}

\begin{longtable}[]{@{}lrlr@{}}
\caption{Variables que más aportan al plano 1-2.}\tabularnewline
\toprule\noalign{}
ACP normado & \% & Sin normalizar & \% \\
\midrule\noalign{}
\endfirsthead
\toprule\noalign{}
ACP normado & \% & Sin normalizar & \% \\
\midrule\noalign{}
\endhead
\bottomrule\noalign{}
\endlastfoot
PetLiq\_cons & 2.49 & Refinados\_prod & 18.01 \\
Refinados\_cons & 2.49 & PetLiq\_cons & 17.31 \\
PetLiq\_CO2 & 2.48 & Refinados\_cons & 17.31 \\
Refinados\_prod & 2.47 & PetLiq\_prod & 16.58 \\
Destilado\_prod & 2.45 & Crudo\_prod & 10.79 \\
Destilado\_cons & 2.44 & Crudo\_imp & 3.57 \\
Elec\_cap & 2.43 & Crudo\_exp & 3.01 \\
Elec\_cons & 2.42 & Gasolina\_prod & 2.84 \\
Elec\_gen & 2.42 & Gasolina\_cons & 2.53 \\
NoHidro\_cap & 2.42 & Destilado\_prod & 1.36 \\
\end{longtable}

\Needspace{3.23in}

\begin{Shaded}
\begin{Highlighting}[]
\NormalTok{varianza }\OtherTok{\textless{}{-}} \FunctionTok{sort}\NormalTok{(}\FunctionTok{sapply}\NormalTok{(num, var), }\AttributeTok{decreasing =} \ConstantTok{TRUE}\NormalTok{)}
\NormalTok{knitr}\SpecialCharTok{::}\FunctionTok{kable}\NormalTok{(}\FunctionTok{data.frame}\NormalTok{(}
  \AttributeTok{variable =} \FunctionTok{names}\NormalTok{(varianza)[}\DecValTok{1}\SpecialCharTok{:}\DecValTok{5}\NormalTok{],}
  \AttributeTok{porcentaje =} \FunctionTok{round}\NormalTok{(}\DecValTok{100} \SpecialCharTok{*}\NormalTok{ varianza[}\DecValTok{1}\SpecialCharTok{:}\DecValTok{5}\NormalTok{] }\SpecialCharTok{/} \FunctionTok{sum}\NormalTok{(varianza), }\DecValTok{1}\NormalTok{),}
  \AttributeTok{row.names =} \ConstantTok{NULL}\NormalTok{),}
  \AttributeTok{caption =} \StringTok{"Peso en la varianza total sin normalizar."}\NormalTok{)}
\end{Highlighting}
\end{Shaded}

\Needspace{2.12in}

\begin{longtable}[]{@{}lr@{}}
\caption{Peso en la varianza total sin normalizar.}\tabularnewline
\toprule\noalign{}
variable & porcentaje \\
\midrule\noalign{}
\endfirsthead
\toprule\noalign{}
variable & porcentaje \\
\midrule\noalign{}
\endhead
\bottomrule\noalign{}
\endlastfoot
Refinados\_prod & 17.7 \\
PetLiq\_cons & 16.9 \\
Refinados\_cons & 16.9 \\
PetLiq\_prod & 16.2 \\
Crudo\_prod & 10.5 \\
\end{longtable}

\Needspace{1.52in}

\begin{codres}

\begin{Shaded}
\begin{Highlighting}[]
\NormalTok{normado }\OtherTok{\textless{}{-}} \FunctionTok{facto\_summarize}\NormalTok{(}
\NormalTok{  acp\_normado, }\StringTok{"var"}\NormalTok{,}
  \AttributeTok{axes =} \FunctionTok{c}\NormalTok{(}\DecValTok{1}\NormalTok{, }\DecValTok{2}\NormalTok{))}
\FunctionTok{sum}\NormalTok{(normado}\SpecialCharTok{$}\NormalTok{contrib }\SpecialCharTok{\textgreater{}} \DecValTok{100} \SpecialCharTok{/} \DecValTok{57}\NormalTok{)}
\end{Highlighting}
\end{Shaded}

\codressep

\begin{verbatim}
[1] 37
\end{verbatim}

\end{codres}

\Needspace{4.27in}

\begin{Shaded}
\begin{Highlighting}[]
\FunctionTok{fviz\_contrib}\NormalTok{(acp\_normado, }\AttributeTok{choice =} \StringTok{"var"}\NormalTok{, }\AttributeTok{axes =} \FunctionTok{c}\NormalTok{(}\DecValTok{1}\NormalTok{, }\DecValTok{2}\NormalTok{),}
             \AttributeTok{top =} \DecValTok{20}\NormalTok{) }\SpecialCharTok{+}
  \FunctionTok{labs}\NormalTok{(}\AttributeTok{title =} \ConstantTok{NULL}\NormalTok{)}
\end{Highlighting}
\end{Shaded}

\begin{figure}[H]

{\centering \pandocbounded{\includegraphics[keepaspectratio]{sda-lab-exploracion_files/figure-latex/p3-normado-1.pdf}}

}

\caption{Contribución de las 20 primeras variables al plano 1-2 del ACP
normado: Todas quedan cerca de la referencia uniforme (línea roja) pues
el plano lo construyen muchas variables a la vez.}

\end{figure}%

\Needspace{4.08in}

\begin{Shaded}
\begin{Highlighting}[]
\FunctionTok{fviz\_contrib}\NormalTok{(acp\_crudo, }\AttributeTok{choice =} \StringTok{"var"}\NormalTok{, }\AttributeTok{axes =} \FunctionTok{c}\NormalTok{(}\DecValTok{1}\NormalTok{, }\DecValTok{2}\NormalTok{), }\AttributeTok{top =} \DecValTok{20}\NormalTok{) }\SpecialCharTok{+}
  \FunctionTok{labs}\NormalTok{(}\AttributeTok{title =} \ConstantTok{NULL}\NormalTok{)}
\end{Highlighting}
\end{Shaded}

\begin{figure}[H]

{\centering \pandocbounded{\includegraphics[keepaspectratio]{sda-lab-exploracion_files/figure-latex/p3-crudo-1.pdf}}

}

\caption{Contribución de las 20 primeras variables al plano 1-2 del ACP
sin normalizar: Cinco variables medidas en TBPD construyen el 80 \% del
plano.}

\end{figure}%

\textbf{¿Cómo cambian las contribuciones al normalizar?} Cambian por
completo pues en el ACP sin normalizar el plano lo construyen casi solo
cinco variables de barriles por día (producción de refinados con 18.01
\%, consumo de petróleo y de refinados con 17.31 \% cada una, producción
de petróleo con 16.58 \% y producción de crudo con 10.79 \%), que suman
el 80 \% del plano. Esto sucede porque sin estandarizar cada variable
pesa según su varianza, de manera que esas cinco variables concentran
por sí solas el 78.2 \% de la varianza total (medidas en miles de
barriles diarios alcanzan valores de miles), mientras que por ejemplo la
generación eléctrica en BKWH tiene una varianza once veces menor que la
del consumo de petróleo. Así el primer plano explica un 97.25 \% de la
varianza (82.92 \% el primer eje y 14.34 \% el segundo), cifra que no
indica una buena síntesis sino que la varianza total está en pocas
columnas.

En el ACP normado, en cambio, ninguna variable pasa de 2.49 \% y las
diez primeras apenas suman 24.5 \%, con 37 de las 57 variables por
encima de la referencia de 1.75 \%. Las que más aportan son de petróleo
(consumo, emisiones, refinados y destilados) y de electricidad
(capacidad instalada, consumo y generación), lo que significa que el
primer plano resume el tamaño energético del país en ambos sectores. Por
esto se usa el ACP normado en el resto del taller ya que las variables
tienen unidades distintas y no es la unidad de medida la que debe
decidir qué variable importa.

\Needspace{2.04in}

\includegraphics[width=5.2in,height=1.32in]{sda-lab-exploracion_files/figure-latex/mermaid-figure-4.png}

Por qué escalar cambia todo: La unidad de una sola columna decide, sin
normalizar, quién construye el plano, mientras que normando el plano lo
construye la estructura de toda la base.

\begin{enumerate}
\def\labelenumi{\alph{enumi}.}
\tightlist
\item
  \textbf{Lectura local.} El peso de una variable en el ACP sin
  normalizar lo decide su unidad y no lo que mide pues el consumo de
  petróleo tiene una varianza once veces mayor que la generación
  eléctrica solo porque los barriles diarios alcanzan cifras de miles y
  los BKWH de cientos.
\item
  \textbf{Lectura global.} El ACP sin normalizar no resume la base sino
  que la ordena por unidades, mientras que el normado muestra que el
  tamaño energético es una propiedad de toda la base pues lo construyen
  a la vez 37 de las 57 variables de los dos sectores.
\end{enumerate}

\section{Pregunta 4 --- Número de componentes}\label{p4}

El criterio de Kaiser \parencite{kaiser1960} retiene las componentes con
valor propio mayor que 1, es decir, las que explican más varianza que
una sola variable original estandarizada. Se compara con el porcentaje
acumulado (retener hasta alcanzar el 80 \%) y con el codo del gráfico de
sedimentación.

\Needspace{4.37in}

\begin{Shaded}
\begin{Highlighting}[]
\NormalTok{autovalores }\OtherTok{\textless{}{-}}\NormalTok{ acp\_normado}\SpecialCharTok{$}\NormalTok{sdev}\SpecialCharTok{\^{}}\DecValTok{2}
\NormalTok{proporcion }\OtherTok{\textless{}{-}}\NormalTok{ autovalores }\SpecialCharTok{/} \FunctionTok{sum}\NormalTok{(autovalores)}
\NormalTok{knitr}\SpecialCharTok{::}\FunctionTok{kable}\NormalTok{(}\FunctionTok{data.frame}\NormalTok{(}
  \AttributeTok{componente =} \DecValTok{1}\SpecialCharTok{:}\DecValTok{10}\NormalTok{, }\AttributeTok{autovalor =} \FunctionTok{round}\NormalTok{(autovalores[}\DecValTok{1}\SpecialCharTok{:}\DecValTok{10}\NormalTok{], }\DecValTok{2}\NormalTok{),}
  \AttributeTok{porcentaje =} \FunctionTok{round}\NormalTok{(}\DecValTok{100} \SpecialCharTok{*}\NormalTok{ proporcion[}\DecValTok{1}\SpecialCharTok{:}\DecValTok{10}\NormalTok{], }\DecValTok{2}\NormalTok{),}
  \AttributeTok{acumulado =} \FunctionTok{round}\NormalTok{(}\DecValTok{100} \SpecialCharTok{*} \FunctionTok{cumsum}\NormalTok{(proporcion)[}\DecValTok{1}\SpecialCharTok{:}\DecValTok{10}\NormalTok{], }\DecValTok{2}\NormalTok{)),}
  \AttributeTok{caption =} \StringTok{"Valores propios de las diez primeras componentes."}\NormalTok{)}
\end{Highlighting}
\end{Shaded}

\Needspace{3.07in}

\begin{longtable}[]{@{}rrrr@{}}
\caption{Valores propios de las diez primeras
componentes.}\tabularnewline
\toprule\noalign{}
componente & autovalor & porcentaje & acumulado \\
\midrule\noalign{}
\endfirsthead
\toprule\noalign{}
componente & autovalor & porcentaje & acumulado \\
\midrule\noalign{}
\endhead
\bottomrule\noalign{}
\endlastfoot
1 & 35.32 & 61.96 & 61.96 \\
2 & 4.41 & 7.74 & 69.69 \\
3 & 4.00 & 7.01 & 76.71 \\
4 & 2.71 & 4.75 & 81.46 \\
5 & 2.45 & 4.31 & 85.76 \\
6 & 1.87 & 3.28 & 89.04 \\
7 & 1.31 & 2.29 & 91.33 \\
8 & 1.12 & 1.96 & 93.29 \\
9 & 0.90 & 1.58 & 94.87 \\
10 & 0.76 & 1.33 & 96.20 \\
\end{longtable}

\Needspace{2.09in}

\begin{codres}

\begin{Shaded}
\begin{Highlighting}[]
\FunctionTok{sum}\NormalTok{(autovalores }\SpecialCharTok{\textgreater{}} \DecValTok{1}\NormalTok{)}
\end{Highlighting}
\end{Shaded}

\codressep

\begin{verbatim}
[1] 8
\end{verbatim}

\end{codres}

\begin{codres}

\begin{Shaded}
\begin{Highlighting}[]
\FunctionTok{which}\NormalTok{(}\FunctionTok{cumsum}\NormalTok{(proporcion)}
      \SpecialCharTok{\textgreater{}=} \FloatTok{0.8}\NormalTok{)[}\DecValTok{1}\NormalTok{]}
\end{Highlighting}
\end{Shaded}

\codressep

\begin{verbatim}
[1] 4
\end{verbatim}

\end{codres}

\Needspace{3.88in}

\begin{Shaded}
\begin{Highlighting}[]
\FunctionTok{fviz\_eig}\NormalTok{(acp\_normado, }\AttributeTok{addlabels =} \ConstantTok{TRUE}\NormalTok{, }\AttributeTok{ncp =} \DecValTok{10}\NormalTok{) }\SpecialCharTok{+}
  \FunctionTok{scale\_y\_continuous}\NormalTok{(}\AttributeTok{limits =} \FunctionTok{c}\NormalTok{(}\DecValTok{0}\NormalTok{, }\DecValTok{70}\NormalTok{)) }\SpecialCharTok{+} \FunctionTok{labs}\NormalTok{(}\AttributeTok{title =} \ConstantTok{NULL}\NormalTok{)}
\end{Highlighting}
\end{Shaded}

\begin{figure}[H]

{\centering \pandocbounded{\includegraphics[keepaspectratio]{sda-lab-exploracion_files/figure-latex/p4-scree-1.pdf}}

}

\caption{Gráfico de sedimentación del ACP normado: El primer eje explica
el 62 \% y desde el segundo la curva baja despacio.}

\end{figure}%

\textbf{¿Cuántas componentes se deberían seleccionar?} Cuatro, bajo el
criterio del porcentaje de varianza acumulada, dado que con ellas se
explica el 81.46 \% de la variación de las 57 variables (61.96 \% la
primera, 7.74 \% la segunda, 7.01 \% la tercera y 4.75 \% la cuarta).
Los otros dos criterios no sirven bien en este caso:

\begin{enumerate}
\def\labelenumi{\alph{enumi}.}
\tightlist
\item
  \textbf{Codo.} El gráfico de sedimentación se quiebra en la segunda
  componente, pero quedarse con dos componentes deja por fuera el 30.3
  \% de la varianza y además la tercera explica casi lo mismo que la
  segunda (7.01 \% frente a 7.74 \%), por lo que no hay razón para
  cortar entre ellas.
\item
  \textbf{Kaiser.} Retiene ocho componentes (93.29 \%), lo cual es
  excesivo porque con 57 variables este criterio tiende a retener de más
  y las componentes de la quinta a la octava explican entre 1.96 \% y
  4.31 \% cada una.
\end{enumerate}

Retener cuatro es coherente con lo que pide la Pregunta 5, que analiza
el plano (2,3) porque la tercera componente todavía carga información
relevante.

\begin{enumerate}
\def\labelenumi{\alph{enumi}.}
\tightlist
\item
  \textbf{Lectura local.} La cuarta componente apenas agrega 4.75 \% y
  es la que hace cruzar el umbral del 80 \%, de modo que la frontera
  entre tres y cuatro componentes la marca la regla y no un cambio en la
  curva.
\item
  \textbf{Lectura global.} El primer valor propio vale 35.32, es decir
  que una sola componente equivale a 35 de las 57 variables, lo cual es
  la otra cara de la redundancia del diagnóstico pues las 57 columnas
  describen en realidad unas pocas dimensiones de los países.
\end{enumerate}

Para ver si la redundancia del diagnóstico mueve los criterios se repite
el ACP normado quitando un miembro de cada uno de los 18 pares con
correlación mayor a 0.99, dejando siempre las cinco variables de la
Pregunta 1.

\Needspace{1.33in}

\begin{Shaded}
\begin{Highlighting}[]
\NormalTok{nombres }\OtherTok{\textless{}{-}} \FunctionTok{names}\NormalTok{(num)}
\NormalTok{copia }\OtherTok{\textless{}{-}} \FunctionTok{unique}\NormalTok{(}\FunctionTok{ifelse}\NormalTok{(nombres[pares[, }\DecValTok{2}\NormalTok{]] }\SpecialCharTok{\%in\%}\NormalTok{ cinco,}
\NormalTok{                       nombres[pares[, }\DecValTok{1}\NormalTok{]], nombres[pares[, }\DecValTok{2}\NormalTok{]]))}
\FunctionTok{cat}\NormalTok{(}\FunctionTok{strwrap}\NormalTok{(}\FunctionTok{paste}\NormalTok{(copia, }\AttributeTok{collapse =} \StringTok{", "}\NormalTok{), }\DecValTok{64}\NormalTok{), }\AttributeTok{sep =} \StringTok{"}\SpecialCharTok{\textbackslash{}n}\StringTok{"}\NormalTok{)}
\end{Highlighting}
\end{Shaded}

\begin{verbatim}
Gasolina_prod, PetLiq_CO2, Refinados_cons, GananciaRef,
Elec_cons, Elec_cap, Fosil_cap, Fosil_gen, Geotermica_gen,
Hidro_gen, Nuclear_gen, Renov_cap
\end{verbatim}

\Needspace{4.56in}

\begin{Shaded}
\begin{Highlighting}[]
\NormalTok{sin\_copias }\OtherTok{\textless{}{-}}\NormalTok{ num[}\FunctionTok{setdiff}\NormalTok{(nombres, copia)]}
\NormalTok{criterios }\OtherTok{\textless{}{-}} \ControlFlowTok{function}\NormalTok{(x) \{}
\NormalTok{  l }\OtherTok{\textless{}{-}} \FunctionTok{prcomp}\NormalTok{(x, }\AttributeTok{scale. =} \ConstantTok{TRUE}\NormalTok{)}\SpecialCharTok{$}\NormalTok{sdev}\SpecialCharTok{\^{}}\DecValTok{2}
\NormalTok{  p }\OtherTok{\textless{}{-}} \DecValTok{100} \SpecialCharTok{*}\NormalTok{ l }\SpecialCharTok{/} \FunctionTok{sum}\NormalTok{(l)}
\NormalTok{  rx }\OtherTok{\textless{}{-}} \FunctionTok{cor}\NormalTok{(x)}
  \FunctionTok{c}\NormalTok{(}\AttributeTok{variables =} \FunctionTok{ncol}\NormalTok{(x), }\AttributeTok{pares =} \FunctionTok{sum}\NormalTok{(rx[}\FunctionTok{upper.tri}\NormalTok{(rx)] }\SpecialCharTok{\textgreater{}} \FloatTok{0.99}\NormalTok{),}
    \AttributeTok{cp1 =}\NormalTok{ p[}\DecValTok{1}\NormalTok{], }\AttributeTok{cp2 =}\NormalTok{ p[}\DecValTok{2}\NormalTok{], }\AttributeTok{cp3 =}\NormalTok{ p[}\DecValTok{3}\NormalTok{],}
    \AttributeTok{acumulado\_4 =} \FunctionTok{sum}\NormalTok{(p[}\DecValTok{1}\SpecialCharTok{:}\DecValTok{4}\NormalTok{]), }\AttributeTok{regla\_80 =} \FunctionTok{which}\NormalTok{(}\FunctionTok{cumsum}\NormalTok{(p) }\SpecialCharTok{\textgreater{}=} \DecValTok{80}\NormalTok{)[}\DecValTok{1}\NormalTok{],}
    \AttributeTok{kaiser =} \FunctionTok{sum}\NormalTok{(l }\SpecialCharTok{\textgreater{}} \DecValTok{1}\NormalTok{))}
\NormalTok{\}}
\NormalTok{knitr}\SpecialCharTok{::}\FunctionTok{kable}\NormalTok{(}\FunctionTok{data.frame}\NormalTok{(}\AttributeTok{con\_57 =} \FunctionTok{criterios}\NormalTok{(num),}
                        \AttributeTok{sin\_copias =} \FunctionTok{criterios}\NormalTok{(sin\_copias)),}
             \AttributeTok{digits =} \DecValTok{2}\NormalTok{,}
             \AttributeTok{caption =} \StringTok{"ACP normado con y sin las copias."}\NormalTok{)}
\end{Highlighting}
\end{Shaded}

\Needspace{2.69in}

\begin{longtable}[]{@{}lrr@{}}
\caption{ACP normado con y sin las copias.}\tabularnewline
\toprule\noalign{}
& con\_57 & sin\_copias \\
\midrule\noalign{}
\endfirsthead
\toprule\noalign{}
& con\_57 & sin\_copias \\
\midrule\noalign{}
\endhead
\bottomrule\noalign{}
\endlastfoot
variables & 57.00 & 45.00 \\
pares & 18.00 & 0.00 \\
cp1 & 61.96 & 57.92 \\
cp2 & 7.74 & 9.25 \\
cp3 & 7.01 & 6.55 \\
acumulado\_4 & 81.46 & 78.99 \\
regla\_80 & 4.00 & 5.00 \\
kaiser & 8.00 & 7.00 \\
\end{longtable}

\begin{notas}

\textbf{Interpretaciones.}

\begin{enumerate}
\def\labelenumi{\alph{enumi}.}
\tightlist
\item
  \textbf{Tres criterios, tres respuestas.} El codo da dos componentes
  (69.69 \%), el 80 \% da cuatro (81.46 \%) y Kaiser da ocho (93.29 \%),
  de modo que el número no lo decide el dato sino el criterio.
\item
  \textbf{Lo que sí se interpreta.} Las preguntas 5 y 6 encuentran
  sentido hasta el tercer eje (exportadores frente a importadores de
  electricidad), lo que respalda quedarse con cuatro y no con dos.
\end{enumerate}

\textbf{Decisiones.}

\begin{itemize}
\item
  \textbf{Cuatro componentes por la regla del 80 \%.} Es el criterio que
  el curso toma como referencia y con él cuatro componentes explican el
  81.46 \%, mientras que el codo deja dos (69.69 \%) y Kaiser ocho
  (93.29 \%), de modo que la regla queda entre el criterio gráfico (que
  deja por fuera el 30.3 \% de la varianza) y el más permisivo (que
  retiene componentes de entre 1.96 \% y 4.31 \% cada una). A esto se
  suma que la tercera componente todavía se interpreta en las preguntas
  5 y 6, lo que hace que dos componentes se queden cortas.

  \emph{Respaldo.} En la sesión 4 se presentó el 80 \% como «la regla
  que usualmente usamos», el codo como el método gráfico y Kaiser como
  el criterio que «admite muchas» componentes, mientras que en la sesión
  5 se repitió que el 80 \% separa la señal del ruido.
  \textcite[\S~6.1.1]{jolliffe2002} ubica el corte usual del porcentaje
  acumulado entre el 70 \% y el 90 \%.
\item
  \textbf{Las copias no cambian la estructura pero dejan la cuarta
  componente en el límite.} Al quitar las 12 columnas redundantes de los
  18 pares quedan 45 variables sin ningún par mayor a 0.99, donde el
  primer eje baja de 61.96 \% a 57.92 \%, el codo sigue en la segunda
  componente (9.25 \% frente a 6.55 \% de la tercera) y Kaiser pasa de
  ocho a siete. La regla del 80 \% en cambio pide cinco pues las cuatro
  primeras explican el 78.99 \%, es decir que las copias inflan el
  primer eje en unos cuatro puntos y que con ellas cuatro componentes
  pasan el umbral por 1.46 puntos mientras que sin ellas les falta 1.01.
  Se mantienen las 57 variables y las cuatro componentes porque la
  Pregunta 3 pide el ACP sobre la base completa y lo que se interpreta
  (un eje de tamaño dominante y el codo en dos) no cambia, pero la
  cuarta componente queda en el límite pues sin las copias no alcanzaría
  el 80 \% y su inclusión la decide el umbral más que la curva.

  \emph{Respaldo.} En la sesión 4 se aclaró que el 80 \% «es usualmente
  un punto de referencia» y no un corte exacto, de modo que una
  diferencia de un punto alrededor de él no obliga por sí sola a cambiar
  el número de componentes que se interpretan aunque sí deja a la última
  en duda.
\end{itemize}

\end{notas}

\section{Pregunta 5 --- Representación de las variables}\label{p5}

El círculo de correlaciones ubica cada variable con coordenadas iguales
a su correlación con cada eje del plano, de modo que la longitud de la
flecha es \(\sqrt{\cos^2}\), donde el \(\cos^2\) (la correlación al
cuadrado con el eje) mide qué tan bien representa el plano a esa
variable. Una flecha que llega al círculo está bien representada,
mientras que una corta está mayormente en otras componentes. Se dibujan
solo las 30 variables que más contribuyen al plano.

\Needspace{6.87in}

\begin{Shaded}
\begin{Highlighting}[]
\FunctionTok{fviz\_pca\_var}\NormalTok{(acp\_normado, }\AttributeTok{axes =} \FunctionTok{c}\NormalTok{(}\DecValTok{1}\NormalTok{, }\DecValTok{2}\NormalTok{),}
             \AttributeTok{select.var =} \FunctionTok{list}\NormalTok{(}\AttributeTok{contrib =} \DecValTok{30}\NormalTok{), }\AttributeTok{col.var =} \StringTok{"contrib"}\NormalTok{,}
             \AttributeTok{labelsize =} \FloatTok{2.2}\NormalTok{, }\AttributeTok{repel =} \ConstantTok{TRUE}\NormalTok{) }\SpecialCharTok{+} \FunctionTok{labs}\NormalTok{(}\AttributeTok{title =} \ConstantTok{NULL}\NormalTok{)}
\end{Highlighting}
\end{Shaded}

\begin{figure}[H]

{\centering \pandocbounded{\includegraphics[keepaspectratio]{sda-lab-exploracion_files/figure-latex/p5-12-1.pdf}}

}

\caption{Círculo de correlaciones del ACP normado en el plano (1,2), con
las 30 variables que más contribuyen: Casi todas apuntan hacia la
derecha, de modo que el primer eje es el tamaño energético del país.}

\end{figure}%

\Needspace{7.41in}

\begin{Shaded}
\begin{Highlighting}[]
\CommentTok{\# Las mismas 30 sin la de bombeo, que apunta al otro lado.}
\NormalTok{plano }\OtherTok{\textless{}{-}} \FunctionTok{facto\_summarize}\NormalTok{(acp\_normado, }\StringTok{"var"}\NormalTok{, }\AttributeTok{axes =} \FunctionTok{c}\NormalTok{(}\DecValTok{1}\NormalTok{, }\DecValTok{2}\NormalTok{))}
\NormalTok{haz }\OtherTok{\textless{}{-}} \FunctionTok{setdiff}\NormalTok{(plano}\SpecialCharTok{$}\NormalTok{name[}\FunctionTok{order}\NormalTok{(}\SpecialCharTok{{-}}\NormalTok{plano}\SpecialCharTok{$}\NormalTok{contrib)][}\DecValTok{1}\SpecialCharTok{:}\DecValTok{30}\NormalTok{],}
               \StringTok{"Bombeo\_gen"}\NormalTok{)}
\FunctionTok{fviz\_pca\_var}\NormalTok{(acp\_normado, }\AttributeTok{axes =} \FunctionTok{c}\NormalTok{(}\DecValTok{1}\NormalTok{, }\DecValTok{2}\NormalTok{),}
             \AttributeTok{select.var =} \FunctionTok{list}\NormalTok{(}\AttributeTok{name =}\NormalTok{ haz), }\AttributeTok{col.var =} \StringTok{"contrib"}\NormalTok{,}
             \AttributeTok{labelsize =} \FloatTok{2.6}\NormalTok{, }\AttributeTok{repel =} \ConstantTok{TRUE}\NormalTok{) }\SpecialCharTok{+}
  \FunctionTok{coord\_cartesian}\NormalTok{(}\AttributeTok{xlim =} \FunctionTok{c}\NormalTok{(}\FloatTok{0.65}\NormalTok{, }\FloatTok{1.02}\NormalTok{), }\AttributeTok{ylim =} \FunctionTok{c}\NormalTok{(}\SpecialCharTok{{-}}\FloatTok{0.42}\NormalTok{, }\FloatTok{0.68}\NormalTok{)) }\SpecialCharTok{+}
  \FunctionTok{labs}\NormalTok{(}\AttributeTok{title =} \ConstantTok{NULL}\NormalTok{)}
\end{Highlighting}
\end{Shaded}

\begin{figure}[H]

{\centering \pandocbounded{\includegraphics[keepaspectratio]{sda-lab-exploracion_files/figure-latex/p5-zoom-1.pdf}}

}

\caption{Acercamiento al haz de flechas del plano (1,2): La mayoría de
las de petróleo quedan arriba del eje horizontal y las de electricidad
abajo, todas con correlación de al menos 0.70 con el primer eje.}

\end{figure}%

\Needspace{2.09in}

\begin{Shaded}
\begin{Highlighting}[]
\NormalTok{coord }\OtherTok{\textless{}{-}} \FunctionTok{get\_pca\_var}\NormalTok{(acp\_normado)}\SpecialCharTok{$}\NormalTok{coord[, }\DecValTok{1}\SpecialCharTok{:}\DecValTok{3}\NormalTok{]}
\FunctionTok{round}\NormalTok{(coord[}\FunctionTok{order}\NormalTok{(}\SpecialCharTok{{-}}\FunctionTok{abs}\NormalTok{(coord[, }\DecValTok{2}\NormalTok{]))[}\DecValTok{1}\SpecialCharTok{:}\DecValTok{8}\NormalTok{], ], }\DecValTok{2}\NormalTok{)}
\end{Highlighting}
\end{Shaded}

\begin{verbatim}
               Dim.1 Dim.2 Dim.3
LGN_prod        0.70  0.65  0.06
Crudo_exp       0.13  0.59 -0.68
PetLiq_prod     0.69  0.58 -0.35
Crudo_prod      0.60  0.57 -0.49
Crudo_reservas  0.15  0.46 -0.60
Hidro_gen       0.74 -0.42 -0.35
Hidro_cap       0.78 -0.41 -0.33
Solar_cap       0.77 -0.40  0.05
\end{verbatim}

\Needspace{1.90in}

\begin{Shaded}
\begin{Highlighting}[]
\FunctionTok{round}\NormalTok{(coord[}\FunctionTok{order}\NormalTok{(}\SpecialCharTok{{-}}\FunctionTok{abs}\NormalTok{(coord[, }\DecValTok{3}\NormalTok{]))[}\DecValTok{1}\SpecialCharTok{:}\DecValTok{8}\NormalTok{], ], }\DecValTok{2}\NormalTok{)}
\end{Highlighting}
\end{Shaded}

\begin{verbatim}
               Dim.1 Dim.2 Dim.3
Crudo_exp       0.13  0.59 -0.68
Crudo_reservas  0.15  0.46 -0.60
Residual_prod   0.49  0.35 -0.53
Crudo_prod      0.60  0.57 -0.49
Residual_cons   0.53  0.09 -0.48
Geotermica_cap  0.51  0.33  0.46
Geotermica_gen  0.51  0.34  0.46
Elec_imp        0.54  0.13  0.41
\end{verbatim}

\Needspace{0.95in}

\begin{codres}

\begin{Shaded}
\begin{Highlighting}[]
\FunctionTok{sum}\NormalTok{(}\FunctionTok{abs}\NormalTok{(coord[, }\DecValTok{1}\NormalTok{]) }\SpecialCharTok{\textgreater{}} \FloatTok{0.7}\NormalTok{)}
\end{Highlighting}
\end{Shaded}

\codressep

\begin{verbatim}
[1] 39
\end{verbatim}

\end{codres}

\textbf{Plano (1,2).} Las 30 flechas más largas llegan casi al borde del
círculo, lo cual muestra que el plano las representa muy bien (el
consumo de petróleo tiene \(\cos^2 = 0.988\) y la generación eléctrica
0.961). Se concluye que:

\begin{enumerate}
\def\labelenumi{\alph{enumi}.}
\tightlist
\item
  \textbf{El primer eje es un factor de tamaño.} Prácticamente todas las
  variables tienen correlación alta y positiva con él, tanto las de
  petróleo como las de electricidad, por lo que un país está a la
  derecha si produce, consume, genera e importa mucho en términos
  absolutos. La única flecha que apunta hacia la izquierda es la
  generación neta por bombeo hidroeléctrico (\texttt{Bombeo\_gen},
  \(r = -0.90\)), lo cual tiene sentido porque esa generación neta es
  negativa (el bombeo consume más energía de la que devuelve) en 27
  países y lo es más en los grandes, como Estados Unidos (\(-6.2\) BKWH)
  y China (\(-4.7\) BKWH).
\item
  \textbf{El segundo eje separa petróleo de renovables.} Según las
  correlaciones impresas arriba, del lado positivo quedan los líquidos
  de gas natural (\texttt{LGN\_prod}, 0.65), las exportaciones de crudo
  (0.59) y la producción de crudo y de petróleo (0.57 y 0.59), mientras
  que del lado negativo quedan la generación y la capacidad
  hidroeléctrica (−0.42 y −0.41) y la capacidad solar y renovable (−0.40
  y −0.39). Es decir que entre países de igual tamaño el segundo eje
  distingue a los productores de hidrocarburos de los que generan
  electricidad renovable.
\end{enumerate}

\Needspace{5.87in}

\begin{Shaded}
\begin{Highlighting}[]
\FunctionTok{fviz\_pca\_var}\NormalTok{(acp\_normado, }\AttributeTok{axes =} \FunctionTok{c}\NormalTok{(}\DecValTok{2}\NormalTok{, }\DecValTok{3}\NormalTok{),}
             \AttributeTok{select.var =} \FunctionTok{list}\NormalTok{(}\AttributeTok{contrib =} \DecValTok{30}\NormalTok{), }\AttributeTok{col.var =} \StringTok{"contrib"}\NormalTok{,}
             \AttributeTok{labelsize =} \FloatTok{2.5}\NormalTok{, }\AttributeTok{repel =} \ConstantTok{TRUE}\NormalTok{) }\SpecialCharTok{+} \FunctionTok{labs}\NormalTok{(}\AttributeTok{title =} \ConstantTok{NULL}\NormalTok{)}
\end{Highlighting}
\end{Shaded}

\begin{figure}[H]

{\centering \pandocbounded{\includegraphics[keepaspectratio]{sda-lab-exploracion_files/figure-latex/p5-23-1.pdf}}

}

\caption{Círculo de correlaciones del ACP normado en el plano (2,3), con
las 30 variables que más contribuyen: Las exportaciones y reservas de
crudo bajan, las renovables van a la izquierda y la geotermia y las
importaciones eléctricas suben.}

\end{figure}%

\textbf{Plano (2,3).} Al quitar el efecto de tamaño las flechas se
acortan y se abren en varias direcciones, lo que indica que este plano
representa peor a cada variable (el mejor \(\cos^2\) es el de las
exportaciones de crudo con 0.81) pero muestra perfiles que el primero
escondía:

\begin{enumerate}
\def\labelenumi{\alph{enumi}.}
\tightlist
\item
  \textbf{Perfil exportador de crudo.} Abajo a la derecha quedan las
  exportaciones de crudo (\(-0.68\) con el tercer eje), las reservas
  (\(-0.60\)), la producción de combustóleo residual (\(-0.53\)) y la
  producción de crudo (\(-0.49\)), es decir, los países con yacimientos
  que venden su petróleo.
\item
  \textbf{Perfil renovable.} A la izquierda quedan la capacidad y
  generación hidroeléctrica, solar y renovable, con correlación negativa
  con el segundo eje.
\item
  \textbf{Perfil importador de electricidad.} Arriba quedan la capacidad
  y la generación geotérmica (0.46), las importaciones de electricidad
  (0.41) y las importaciones netas (0.37), así como la ganancia de
  refinación y la electricidad nuclear con correlaciones menores.
\end{enumerate}

Se concluye que el tercer eje separa a los exportadores de crudo de los
países que compran electricidad o la generan con fuentes como la
geotérmica y la nuclear, lo cual justifica haber retenido más de dos
componentes.

El diagrama resume lo que significa cada uno de los tres primeros ejes
según las variables que más se asocian con él.

\Needspace{4.30in}

\includegraphics[width=5.2in,height=3.46in]{sda-lab-exploracion_files/figure-latex/mermaid-figure-3.png}

Qué mide cada eje del ACP normado con los 111 países y qué variables lo
construyen: El primero ordena por tamaño y solo los siguientes
distinguen perfiles energéticos.

\begin{enumerate}
\def\labelenumi{\alph{enumi}.}
\tightlist
\item
  \textbf{Lectura local.} La exportación de crudo, que en el primer
  plano casi no aparece (0.13 con el primer eje), es la variable mejor
  representada del plano (2,3), lo cual respalda que la dimensión
  comercial que se buscó en la Pregunta 1 existe y está fuera del eje de
  tamaño.
\item
  \textbf{Lectura global.} El primer eje es el tamaño porque 39 de las
  57 variables tienen con él correlación mayor a 0.7 en valor absoluto,
  de modo que lo que distingue a los países más allá de su tamaño
  (petróleo, renovables, importaciones) solo aparece en los ejes
  siguientes, que explican mucho menos.
\end{enumerate}

\begin{quote}
\textbf{Nota.} El signo de cada componente es arbitrario pues
\texttt{prcomp()} puede devolver \(u_k\) o \(-u_k\) y ambos son igual de
válidos. Lo que se interpreta es la oposición entre grupos de variables
y no si quedan arriba o abajo.
\end{quote}

\section{Pregunta 6 --- Mapa de individuos}\label{p6}

Para que el mapa se pueda leer se dibujan todos los países como puntos y
se rotulan solo los 20 más alejados del origen en cada plano, que son
los que más contribuyen a él.

\Needspace{2.47in}

\begin{Shaded}
\begin{Highlighting}[]
\NormalTok{mapa\_paises }\OtherTok{\textless{}{-}} \ControlFlowTok{function}\NormalTok{(acp, ejes, }\AttributeTok{n =} \DecValTok{20}\NormalTok{, }\AttributeTok{grupos =} \StringTok{"none"}\NormalTok{) \{}
\NormalTok{  coord }\OtherTok{\textless{}{-}} \FunctionTok{data.frame}\NormalTok{(acp}\SpecialCharTok{$}\NormalTok{x[, ejes])}
  \FunctionTok{names}\NormalTok{(coord) }\OtherTok{\textless{}{-}} \FunctionTok{c}\NormalTok{(}\StringTok{"x"}\NormalTok{, }\StringTok{"y"}\NormalTok{)}
\NormalTok{  coord}\SpecialCharTok{$}\NormalTok{pais }\OtherTok{\textless{}{-}} \FunctionTok{rownames}\NormalTok{(coord)}
\NormalTok{  lejanos }\OtherTok{\textless{}{-}}\NormalTok{ coord[}\FunctionTok{order}\NormalTok{(}\SpecialCharTok{{-}}\NormalTok{(coord}\SpecialCharTok{$}\NormalTok{x}\SpecialCharTok{\^{}}\DecValTok{2} \SpecialCharTok{+}\NormalTok{ coord}\SpecialCharTok{$}\NormalTok{y}\SpecialCharTok{\^{}}\DecValTok{2}\NormalTok{))[}\FunctionTok{seq\_len}\NormalTok{(n)], ]}
  \FunctionTok{fviz\_pca\_ind}\NormalTok{(acp, }\AttributeTok{axes =}\NormalTok{ ejes, }\AttributeTok{geom =} \StringTok{"point"}\NormalTok{, }\AttributeTok{habillage =}\NormalTok{ grupos,}
               \AttributeTok{col.ind =} \StringTok{"grey45"}\NormalTok{, }\AttributeTok{mean.point =} \ConstantTok{FALSE}\NormalTok{) }\SpecialCharTok{+}
\NormalTok{    ggrepel}\SpecialCharTok{::}\FunctionTok{geom\_text\_repel}\NormalTok{(}\AttributeTok{data =}\NormalTok{ lejanos, }\AttributeTok{inherit.aes =} \ConstantTok{FALSE}\NormalTok{,}
                             \FunctionTok{aes}\NormalTok{(x, y, }\AttributeTok{label =}\NormalTok{ pais), }\AttributeTok{size =} \FloatTok{2.4}\NormalTok{,}
                             \AttributeTok{max.overlaps =} \ConstantTok{Inf}\NormalTok{,}
                             \AttributeTok{min.segment.length =} \DecValTok{0}\NormalTok{) }\SpecialCharTok{+}
    \FunctionTok{labs}\NormalTok{(}\AttributeTok{title =} \ConstantTok{NULL}\NormalTok{, }\AttributeTok{colour =} \StringTok{"Grupo"}\NormalTok{, }\AttributeTok{shape =} \StringTok{"Grupo"}\NormalTok{)}
\NormalTok{\}}
\end{Highlighting}
\end{Shaded}

\Needspace{4.69in}

\begin{Shaded}
\begin{Highlighting}[]
\FunctionTok{mapa\_paises}\NormalTok{(acp\_normado, }\FunctionTok{c}\NormalTok{(}\DecValTok{1}\NormalTok{, }\DecValTok{2}\NormalTok{))}
\end{Highlighting}
\end{Shaded}

\begin{figure}[H]

{\centering \pandocbounded{\includegraphics[keepaspectratio]{sda-lab-exploracion_files/figure-latex/p6-12-1.pdf}}

}

\caption{Países en el plano (1,2) del ACP normado: Estados Unidos y
China quedan aislados a la derecha y el resto se amontona cerca del
origen.}

\end{figure}%

\Needspace{4.69in}

\begin{Shaded}
\begin{Highlighting}[]
\FunctionTok{mapa\_paises}\NormalTok{(acp\_normado, }\FunctionTok{c}\NormalTok{(}\DecValTok{2}\NormalTok{, }\DecValTok{3}\NormalTok{))}
\end{Highlighting}
\end{Shaded}

\begin{figure}[H]

{\centering \pandocbounded{\includegraphics[keepaspectratio]{sda-lab-exploracion_files/figure-latex/p6-23-1.pdf}}

}

\caption{Países en el plano (2,3) del ACP normado: Arabia Saudita y
Rusia quedan juntos abajo a la derecha, China a la izquierda y Estados
Unidos arriba a la derecha.}

\end{figure}%

\Needspace{0.76in}

\begin{Shaded}
\begin{Highlighting}[]
\NormalTok{contribucion }\OtherTok{\textless{}{-}} \FunctionTok{get\_pca\_ind}\NormalTok{(acp\_normado)}\SpecialCharTok{$}\NormalTok{contrib}
\FunctionTok{round}\NormalTok{(}\FunctionTok{head}\NormalTok{(}\FunctionTok{sort}\NormalTok{(contribucion[, }\DecValTok{1}\NormalTok{], }\AttributeTok{decreasing =} \ConstantTok{TRUE}\NormalTok{), }\DecValTok{4}\NormalTok{), }\DecValTok{2}\NormalTok{)}
\end{Highlighting}
\end{Shaded}

\begin{verbatim}
UnitedStates        China        India       Russia 
       52.02        30.96         2.01         1.98 
\end{verbatim}

\Needspace{2.47in}

\begin{Shaded}
\begin{Highlighting}[]
\FunctionTok{round}\NormalTok{(acp\_normado}\SpecialCharTok{$}\NormalTok{x[}\FunctionTok{c}\NormalTok{(}\StringTok{"SaudiArabia"}\NormalTok{, }\StringTok{"Russia"}\NormalTok{, }\StringTok{"Kuwait"}\NormalTok{, }\StringTok{"Iraq"}\NormalTok{,}
                      \StringTok{"Japan"}\NormalTok{, }\StringTok{"India"}\NormalTok{, }\StringTok{"Italy"}\NormalTok{, }\StringTok{"France"}\NormalTok{,}
                      \StringTok{"Germany"}\NormalTok{), }\DecValTok{1}\SpecialCharTok{:}\DecValTok{3}\NormalTok{], }\DecValTok{2}\NormalTok{)}
\end{Highlighting}
\end{Shaded}

\begin{verbatim}
              PC1   PC2   PC3
SaudiArabia  3.82  8.48 -8.63
Russia       8.77  6.65 -8.26
Kuwait      -0.82  2.03 -2.68
Iraq        -0.52  2.49 -2.73
Japan        8.65 -3.65 -2.15
India        8.83 -2.88 -2.25
Italy        4.31 -2.25  4.12
France       4.01 -2.56  1.83
Germany      1.65 -2.72  1.69
\end{verbatim}

\Needspace{1.14in}

\begin{codres}

\begin{Shaded}
\begin{Highlighting}[]
\FunctionTok{mean}\NormalTok{(}\FunctionTok{abs}\NormalTok{(acp\_normado}\SpecialCharTok{$}\NormalTok{x[, }\DecValTok{1}\NormalTok{])}
     \SpecialCharTok{\textless{}} \DecValTok{2}\NormalTok{)}
\end{Highlighting}
\end{Shaded}

\codressep

\begin{verbatim}
[1] 0.5495495
\end{verbatim}

\end{codres}

\textbf{¿Qué se concluye de la cercanía entre países?} Que el ACP con
los 111 países está dominado por cuatro de ellos y que la cercanía
revela perfiles energéticos comunes:

\begin{enumerate}
\def\labelenumi{\alph{enumi}.}
\tightlist
\item
  \textbf{Dos gigantes.} Estados Unidos (44.95 en el primer eje) y China
  (34.68) están tan lejos del resto que construyen el 52.02 \% y el
  30.96 \% del primer eje respectivamente, es decir, entre los dos el 83
  \%. Ningún otro país pasa de 9 en esa coordenada y el 55 \% de los
  países queda a menos de 2 unidades del origen.
\item
  \textbf{Petroleros.} En el plano (2,3) Arabia Saudita (8.48, −8.63) y
  Rusia (6.65, −8.26) están juntos y lejos del resto, en la dirección de
  las exportaciones y reservas de crudo de la Pregunta 5, mientras que
  Kuwait, Irak, Emiratos Árabes, Venezuela e Irán aparecen en la misma
  dirección aunque mucho más cerca del origen.
\item
  \textbf{Grandes consumidores sin petróleo.} Japón (8.65, −3.65) e
  India (8.83, −2.88) quedan prácticamente superpuestos en el plano
  (1,2) pues ambos son economías grandes que consumen mucho petróleo y
  electricidad sin producir crudo en cantidad.
\item
  \textbf{China frente a Estados Unidos.} A pesar de estar ambos a la
  derecha del primer plano, se separan en el segundo eje (−12.99 contra
  9.48) ya que China pesa por su generación hidroeléctrica y renovable
  mientras que Estados Unidos lo hace también por su producción de
  petróleo.
\item
  \textbf{Importadores de electricidad.} En el plano (2,3) Italia (4.12
  en el tercer eje), Francia (1.83) y Alemania (1.69) quedan arriba a la
  izquierda, que es la dirección de las importaciones de electricidad y
  la generación nuclear y geotérmica de la Pregunta 5.
\end{enumerate}

Estas cercanías se leen en dos niveles:

\begin{enumerate}
\def\labelenumi{\alph{enumi}.}
\tightlist
\item
  \textbf{Lectura local.} La cercanía de India y Japón no dice que sean
  países parecidos en todo sino que tienen el mismo perfil energético,
  es decir, mucho consumo y poca producción de crudo, de modo que dos
  economías que en otros aspectos pueden ser muy distintas quedan
  superpuestas porque las 57 variables solo miden energía.
\item
  \textbf{Lectura global.} El plano lo organizan pocas economías pues
  Estados Unidos y China construyen el 83 \% del primer eje y fijan la
  escala del mapa, de modo que la mayoría de países no se distingue en
  estos planos y se justifica la Pregunta 7.
\end{enumerate}

\section{Pregunta 7 --- ACP sin los cuatro países atípicos}\label{p7}

Desde aquí se excluyen Estados Unidos, China, Arabia Saudita y Rusia,
tanto para el ACP como para el agrupamiento de las preguntas 8 a 11.

\Needspace{4.56in}

\begin{Shaded}
\begin{Highlighting}[]
\NormalTok{atipicos }\OtherTok{\textless{}{-}} \FunctionTok{c}\NormalTok{(}\StringTok{"UnitedStates"}\NormalTok{, }\StringTok{"China"}\NormalTok{, }\StringTok{"SaudiArabia"}\NormalTok{, }\StringTok{"Russia"}\NormalTok{)}
\NormalTok{datos }\OtherTok{\textless{}{-}}\NormalTok{ datos[}\SpecialCharTok{!}\NormalTok{datos}\SpecialCharTok{$}\NormalTok{Country }\SpecialCharTok{\%in\%}\NormalTok{ atipicos, ]}
\NormalTok{num }\OtherTok{\textless{}{-}} \FunctionTok{numericas}\NormalTok{(datos)}
\NormalTok{acp\_sin4 }\OtherTok{\textless{}{-}} \FunctionTok{prcomp}\NormalTok{(num, }\AttributeTok{scale. =} \ConstantTok{TRUE}\NormalTok{)}
\NormalTok{valores }\OtherTok{\textless{}{-}} \ControlFlowTok{function}\NormalTok{(acp) \{}
\NormalTok{  p }\OtherTok{\textless{}{-}}\NormalTok{ acp}\SpecialCharTok{$}\NormalTok{sdev}\SpecialCharTok{\^{}}\DecValTok{2} \SpecialCharTok{/} \FunctionTok{sum}\NormalTok{(acp}\SpecialCharTok{$}\NormalTok{sdev}\SpecialCharTok{\^{}}\DecValTok{2}\NormalTok{)}
  \FunctionTok{round}\NormalTok{(}\DecValTok{100} \SpecialCharTok{*} \FunctionTok{c}\NormalTok{(p[}\DecValTok{1}\SpecialCharTok{:}\DecValTok{3}\NormalTok{], }\FunctionTok{cumsum}\NormalTok{(p)[}\DecValTok{2}\SpecialCharTok{:}\DecValTok{3}\NormalTok{]), }\DecValTok{2}\NormalTok{)}
\NormalTok{\}}
\NormalTok{knitr}\SpecialCharTok{::}\FunctionTok{kable}\NormalTok{(}\FunctionTok{data.frame}\NormalTok{(}
  \AttributeTok{medida =} \FunctionTok{c}\NormalTok{(}\StringTok{"\% CP1"}\NormalTok{, }\StringTok{"\% CP2"}\NormalTok{, }\StringTok{"\% CP3"}\NormalTok{, }\StringTok{"\% plano 1{-}2"}\NormalTok{,}
             \StringTok{"\% acumulado 3"}\NormalTok{),}
  \AttributeTok{con\_111 =} \FunctionTok{valores}\NormalTok{(acp\_normado), }\AttributeTok{sin\_4 =} \FunctionTok{valores}\NormalTok{(acp\_sin4)),}
  \AttributeTok{caption =} \StringTok{"Varianza explicada con y sin los cuatro países."}\NormalTok{)}
\end{Highlighting}
\end{Shaded}

\Needspace{2.12in}

\begin{longtable}[]{@{}lrr@{}}
\caption{Varianza explicada con y sin los cuatro países.}\tabularnewline
\toprule\noalign{}
medida & con\_111 & sin\_4 \\
\midrule\noalign{}
\endfirsthead
\toprule\noalign{}
medida & con\_111 & sin\_4 \\
\midrule\noalign{}
\endhead
\bottomrule\noalign{}
\endlastfoot
\% CP1 & 61.96 & 48.95 \\
\% CP2 & 7.74 & 10.24 \\
\% CP3 & 7.01 & 8.11 \\
\% plano 1-2 & 69.69 & 59.19 \\
\% acumulado 3 & 76.71 & 67.30 \\
\end{longtable}

\Needspace{1.90in}

\begin{Shaded}
\begin{Highlighting}[]
\NormalTok{autovalores }\OtherTok{\textless{}{-}}\NormalTok{ acp\_sin4}\SpecialCharTok{$}\NormalTok{sdev}\SpecialCharTok{\^{}}\DecValTok{2}
\NormalTok{coord }\OtherTok{\textless{}{-}} \FunctionTok{get\_pca\_var}\NormalTok{(acp\_sin4)}\SpecialCharTok{$}\NormalTok{coord}
\FunctionTok{round}\NormalTok{(coord[}\FunctionTok{order}\NormalTok{(}\SpecialCharTok{{-}}\FunctionTok{abs}\NormalTok{(coord[, }\DecValTok{2}\NormalTok{]))[}\DecValTok{1}\SpecialCharTok{:}\DecValTok{6}\NormalTok{], }\DecValTok{1}\SpecialCharTok{:}\DecValTok{2}\NormalTok{], }\DecValTok{2}\NormalTok{)}
\end{Highlighting}
\end{Shaded}

\begin{verbatim}
               Dim.1 Dim.2
PetLiq_prod     0.35 -0.84
Crudo_prod      0.30 -0.83
Crudo_exp       0.16 -0.75
LGN_prod        0.31 -0.69
Crudo_reservas  0.18 -0.66
Bombeo_gen     -0.49 -0.53
\end{verbatim}

\Needspace{0.76in}

\begin{Shaded}
\begin{Highlighting}[]
\FunctionTok{round}\NormalTok{(}\FunctionTok{head}\NormalTok{(}\FunctionTok{sort}\NormalTok{(}\FunctionTok{get\_pca\_ind}\NormalTok{(acp\_sin4)}\SpecialCharTok{$}\NormalTok{contrib[, }\DecValTok{1}\NormalTok{],}
                \AttributeTok{decreasing =} \ConstantTok{TRUE}\NormalTok{), }\DecValTok{4}\NormalTok{), }\DecValTok{2}\NormalTok{)}
\end{Highlighting}
\end{Shaded}

\begin{verbatim}
 India  Japan Brazil Canada 
 23.37  20.04  11.49   6.38 
\end{verbatim}

\Needspace{3.23in}

\begin{codres}

\begin{Shaded}
\begin{Highlighting}[]
\FunctionTok{c}\NormalTok{(}\AttributeTok{paises =} \FunctionTok{nrow}\NormalTok{(num),}
  \AttributeTok{kaiser =} \FunctionTok{sum}\NormalTok{(autovalores }\SpecialCharTok{\textgreater{}} \DecValTok{1}\NormalTok{))}
\end{Highlighting}
\end{Shaded}

\codressep

\begin{verbatim}
paises kaiser 
   107      9 
\end{verbatim}

\end{codres}

\begin{codres}

\begin{Shaded}
\begin{Highlighting}[]
\FunctionTok{c}\NormalTok{(}\AttributeTok{acumulado\_80 =} \FunctionTok{which}\NormalTok{(}
  \FunctionTok{cumsum}\NormalTok{(autovalores) }\SpecialCharTok{/} \DecValTok{57}
    \SpecialCharTok{\textgreater{}=} \FloatTok{0.8}\NormalTok{)[}\DecValTok{1}\NormalTok{],}
  \AttributeTok{tamano =} \FunctionTok{sum}\NormalTok{(}\FunctionTok{abs}\NormalTok{(}
\NormalTok{    coord[, }\DecValTok{1}\NormalTok{]) }\SpecialCharTok{\textgreater{}} \FloatTok{0.7}\NormalTok{))}
\end{Highlighting}
\end{Shaded}

\codressep

\begin{verbatim}
acumulado_80       tamano 
           6           27 
\end{verbatim}

\end{codres}

\Needspace{4.89in}

\begin{Shaded}
\begin{Highlighting}[]
\FunctionTok{mapa\_paises}\NormalTok{(acp\_sin4, }\FunctionTok{c}\NormalTok{(}\DecValTok{1}\NormalTok{, }\DecValTok{2}\NormalTok{), }\AttributeTok{n =} \DecValTok{24}\NormalTok{)}
\end{Highlighting}
\end{Shaded}

\begin{figure}[H]

{\centering \pandocbounded{\includegraphics[keepaspectratio]{sda-lab-exploracion_files/figure-latex/p7-mapa-1.pdf}}

}

\caption{Países en el plano (1,2) del ACP normado sin Estados Unidos,
China, Arabia Saudita y Rusia: La masa central se abre y aparecen dos
brazos, uno de grandes consumidores hacia arriba y otro de productores
de petróleo hacia abajo.}

\end{figure}%

\textbf{¿En qué cambia el ACP al excluir estos países?} En tres
aspectos:

\begin{enumerate}
\def\labelenumi{\alph{enumi}.}
\tightlist
\item
  \textbf{La varianza se reparte más.} La primera componente baja de
  61.96 \% a 48.95 \% mientras que la segunda sube de 7.74 \% a 10.24 \%
  y la tercera de 7.01 \% a 8.11 \%, por lo que ahora se necesitan seis
  componentes para el 80 \% (antes cuatro) y Kaiser retiene nueve (antes
  ocho). Esto muestra que buena parte del 62 \% del primer eje se debía
  a que Estados Unidos y China son enormes en todas las variables a la
  vez.
\item
  \textbf{El segundo eje se vuelve un eje petrolero más nítido.} Las
  correlaciones de la producción de petróleo (\(-0.84\)), la producción
  de crudo (\(-0.83\)), las exportaciones de crudo (\(-0.75\)) y los
  líquidos de gas (\(-0.69\)) con la segunda componente son mucho más
  fuertes que con los 111 países (alrededor de 0.6), mientras que en el
  lado opuesto queda la electricidad solar (0.48). El signo cambió
  respecto al ACP anterior, lo cual no tiene importancia según la nota
  de la Pregunta 5.
\item
  \textbf{Aparecen nuevos países dominantes.} India (23.37 \%) y Japón
  (20.04 \%) pasan a construir el primer eje, de modo que el efecto de
  tamaño no desaparece sino que lo heredan los siguientes países más
  grandes.
\end{enumerate}

\textbf{¿Se perciben clusters con mayor claridad?} Sí, aunque no como
grupos compactos. La masa de países pequeños se abre hacia la derecha y
de ella salen dos brazos, uno hacia arriba con grandes consumidores que
no producen petróleo (Francia, Italia, Corea del Sur, Japón, Alemania y
España) y otro hacia abajo con productores y exportadores de petróleo
(Irak, Emiratos Árabes, Venezuela, Irán y México), al extremo del cual
quedan Brasil y Canadá, que combinan producción de crudo (2255 y 3613
TBPD) con la mayor generación renovable de los 107 países (429 y 412
BKWH). India queda sola en el extremo derecho.

\begin{enumerate}
\def\labelenumi{\alph{enumi}.}
\tightlist
\item
  \textbf{Lectura local.} Quitar cuatro países no borra el tamaño sino
  que lo traslada pues India y Japón construyen ahora el 43 \% del
  primer eje, que es la mitad de lo que construían Estados Unidos y
  China pero sigue siendo mucho para dos países de 107.
\item
  \textbf{Lectura global.} El tamaño es una propiedad de la base y no de
  cuatro países ya que todavía 27 de las 57 variables tienen correlación
  mayor a 0.7 con el primer eje (antes 39). Lo que sí cambia es que, al
  bajar la escala, la dimensión petrolera gana espacio y deja ver la
  estructura en brazos que la Pregunta 8 lee como grupos.
\end{enumerate}

El enunciado fija los cuatro países, pero se puede comprobar si son los
atípicos multivariados de la base con la distancia de Mahalanobis. Sobre
las 57 variables no se puede calcular porque su matriz de covarianzas es
singular así que se calcula sobre las \(k\) primeras componentes del ACP
normado de los 111 países y se compara con el cuantil 0.975 de una
\(\chi^2_k\).

\Needspace{0.57in}

\begin{Shaded}
\begin{Highlighting}[]
\NormalTok{num111 }\OtherTok{\textless{}{-}} \FunctionTok{numericas}\NormalTok{(}\FunctionTok{read.csv}\NormalTok{(ruta, }\AttributeTok{sep =} \StringTok{";"}\NormalTok{, }\AttributeTok{dec =} \StringTok{"."}\NormalTok{))}
\FunctionTok{qr}\NormalTok{(}\FunctionTok{cov}\NormalTok{(num111))}\SpecialCharTok{$}\NormalTok{rank}
\end{Highlighting}
\end{Shaded}

\begin{verbatim}
[1] 50
\end{verbatim}

\Needspace{4.37in}

\begin{Shaded}
\begin{Highlighting}[]
\NormalTok{mahal }\OtherTok{\textless{}{-}} \ControlFlowTok{function}\NormalTok{(k) \{}
\NormalTok{  s }\OtherTok{\textless{}{-}}\NormalTok{ acp\_normado}\SpecialCharTok{$}\NormalTok{x[, }\DecValTok{1}\SpecialCharTok{:}\NormalTok{k]}
\NormalTok{  d2 }\OtherTok{\textless{}{-}} \FunctionTok{sort}\NormalTok{(}\FunctionTok{mahalanobis}\NormalTok{(s, }\FunctionTok{colMeans}\NormalTok{(s), }\FunctionTok{cov}\NormalTok{(s)), }\ConstantTok{TRUE}\NormalTok{)}
\NormalTok{  corte }\OtherTok{\textless{}{-}} \FunctionTok{qchisq}\NormalTok{(}\FloatTok{0.975}\NormalTok{, }\AttributeTok{df =}\NormalTok{ k)}
  \FunctionTok{list}\NormalTok{(}\AttributeTok{fila =} \FunctionTok{c}\NormalTok{(}\AttributeTok{k =}\NormalTok{ k, }\AttributeTok{corte =}\NormalTok{ corte, }\AttributeTok{sobre =} \FunctionTok{sum}\NormalTok{(d2 }\SpecialCharTok{\textgreater{}}\NormalTok{ corte),}
                \AttributeTok{d2\_India =}\NormalTok{ d2[[}\StringTok{"India"}\NormalTok{]], }\AttributeTok{d2\_Japon =}\NormalTok{ d2[[}\StringTok{"Japan"}\NormalTok{]],}
                \AttributeTok{puesto\_India =} \FunctionTok{match}\NormalTok{(}\StringTok{"India"}\NormalTok{, }\FunctionTok{names}\NormalTok{(d2)),}
                \AttributeTok{puesto\_Japon =} \FunctionTok{match}\NormalTok{(}\StringTok{"Japan"}\NormalTok{, }\FunctionTok{names}\NormalTok{(d2))),}
       \AttributeTok{paises =} \FunctionTok{names}\NormalTok{(d2)[d2 }\SpecialCharTok{\textgreater{}}\NormalTok{ corte])}
\NormalTok{\}}
\NormalTok{casos }\OtherTok{\textless{}{-}} \FunctionTok{lapply}\NormalTok{(}\FunctionTok{c}\NormalTok{(}\DecValTok{2}\NormalTok{, }\DecValTok{3}\NormalTok{, }\DecValTok{4}\NormalTok{, }\DecValTok{8}\NormalTok{), mahal)}
\NormalTok{knitr}\SpecialCharTok{::}\FunctionTok{kable}\NormalTok{(}\FunctionTok{t}\NormalTok{(}\FunctionTok{sapply}\NormalTok{(casos, }\StringTok{\textasciigrave{}}\AttributeTok{[[}\StringTok{\textasciigrave{}}\NormalTok{, }\StringTok{"fila"}\NormalTok{)), }\AttributeTok{digits =} \DecValTok{2}\NormalTok{,}
             \AttributeTok{caption =} \StringTok{"Corte chi{-}cuadrado(k, 0.975) en k CP."}\NormalTok{)}
\end{Highlighting}
\end{Shaded}

\Needspace{1.93in}

\begin{longtable}[]{@{}rrrrrrr@{}}
\caption{Corte chi-cuadrado(k, 0.975) en k CP.}\tabularnewline
\toprule\noalign{}
k & corte & sobre & d2\_India & d2\_Japon & puesto\_India &
puesto\_Japon \\
\midrule\noalign{}
\endfirsthead
\toprule\noalign{}
k & corte & sobre & d2\_India & d2\_Japon & puesto\_India &
puesto\_Japon \\
\midrule\noalign{}
\endhead
\bottomrule\noalign{}
\endlastfoot
2 & 7.38 & 4 & 4.09 & 5.14 & 6 & 5 \\
3 & 9.35 & 4 & 5.36 & 6.30 & 7 & 5 \\
4 & 11.14 & 8 & 5.41 & 15.95 & 12 & 7 \\
8 & 17.53 & 11 & 15.43 & 78.60 & 13 & 4 \\
\end{longtable}

\Needspace{0.76in}

\begin{Shaded}
\begin{Highlighting}[]
\FunctionTok{cat}\NormalTok{(}\FunctionTok{strwrap}\NormalTok{(}\FunctionTok{paste}\NormalTok{(casos[[}\DecValTok{3}\NormalTok{]]}\SpecialCharTok{$}\NormalTok{paises, }\AttributeTok{collapse =} \StringTok{", "}\NormalTok{), }\DecValTok{64}\NormalTok{),}
    \AttributeTok{sep =} \StringTok{"}\SpecialCharTok{\textbackslash{}n}\StringTok{"}\NormalTok{)}
\end{Highlighting}
\end{Shaded}

\begin{verbatim}
UnitedStates, China, SaudiArabia, France, SouthKorea, Russia,
Japan, Italy
\end{verbatim}

\Needspace{0.76in}

\begin{Shaded}
\begin{Highlighting}[]
\NormalTok{cp4 }\OtherTok{\textless{}{-}} \FunctionTok{abs}\NormalTok{(acp\_normado}\SpecialCharTok{$}\NormalTok{rotation[, }\DecValTok{4}\NormalTok{])}
\FunctionTok{round}\NormalTok{(}\FunctionTok{head}\NormalTok{(}\FunctionTok{sort}\NormalTok{(cp4, }\AttributeTok{decreasing =} \ConstantTok{TRUE}\NormalTok{), }\DecValTok{3}\NormalTok{), }\DecValTok{2}\NormalTok{)}
\end{Highlighting}
\end{Shaded}

\begin{verbatim}
Mareas_gen Mareas_cap   Elec_exp 
      0.50       0.47       0.27 
\end{verbatim}

\Needspace{0.57in}

\begin{Shaded}
\begin{Highlighting}[]
\FunctionTok{rownames}\NormalTok{(num111)[num111}\SpecialCharTok{$}\NormalTok{Mareas\_gen }\SpecialCharTok{\textgreater{}} \DecValTok{0}\NormalTok{]}
\end{Highlighting}
\end{Shaded}

\begin{verbatim}
[1] "France"     "Germany"    "Italy"      "Slovakia"  
[5] "SouthKorea" "Turkey"    
\end{verbatim}

Como alternativa a excluir países se repite el ACP normado de los 111
países sobre el logaritmo de las variables. Cuatro columnas
(\texttt{PetLiq\_prod}, \texttt{GananciaRef}, \texttt{Elec\_impNeta} y
\texttt{Bombeo\_gen}) traen 80 valores negativos con los que
\texttt{log1p()} daría \texttt{NaN}, por lo que se usa el logaritmo con
signo \(\operatorname{sign}(x)\log(1 + |x|)\), que deja el cero en cero.

\Needspace{3.61in}

\begin{Shaded}
\begin{Highlighting}[]
\NormalTok{log111 }\OtherTok{\textless{}{-}} \FunctionTok{sign}\NormalTok{(num111) }\SpecialCharTok{*} \FunctionTok{log1p}\NormalTok{(}\FunctionTok{abs}\NormalTok{(num111))}
\NormalTok{acp\_log }\OtherTok{\textless{}{-}} \FunctionTok{prcomp}\NormalTok{(log111, }\AttributeTok{scale. =} \ConstantTok{TRUE}\NormalTok{)}
\NormalTok{knitr}\SpecialCharTok{::}\FunctionTok{kable}\NormalTok{(}\FunctionTok{data.frame}\NormalTok{(}
  \AttributeTok{medida =} \FunctionTok{c}\NormalTok{(}\StringTok{"\% CP1"}\NormalTok{, }\StringTok{"\% CP2"}\NormalTok{, }\StringTok{"\% CP3"}\NormalTok{, }\StringTok{"\% plano 1{-}2"}\NormalTok{,}
             \StringTok{"\% acumulado 3"}\NormalTok{),}
  \AttributeTok{con\_111 =} \FunctionTok{valores}\NormalTok{(acp\_normado), }\AttributeTok{sin\_4 =} \FunctionTok{valores}\NormalTok{(acp\_sin4),}
  \AttributeTok{log\_111 =} \FunctionTok{valores}\NormalTok{(acp\_log)),}
  \AttributeTok{caption =} \StringTok{"Varianza explicada sin quitar países, con logaritmo."}\NormalTok{)}
\end{Highlighting}
\end{Shaded}

\Needspace{2.12in}

\begin{longtable}[]{@{}lrrr@{}}
\caption{Varianza explicada sin quitar países, con
logaritmo.}\tabularnewline
\toprule\noalign{}
medida & con\_111 & sin\_4 & log\_111 \\
\midrule\noalign{}
\endfirsthead
\toprule\noalign{}
medida & con\_111 & sin\_4 & log\_111 \\
\midrule\noalign{}
\endhead
\bottomrule\noalign{}
\endlastfoot
\% CP1 & 61.96 & 48.95 & 55.81 \\
\% CP2 & 7.74 & 10.24 & 11.83 \\
\% CP3 & 7.01 & 8.11 & 4.95 \\
\% plano 1-2 & 69.69 & 59.19 & 67.63 \\
\% acumulado 3 & 76.71 & 67.30 & 72.58 \\
\end{longtable}

\Needspace{1.14in}

\begin{Shaded}
\begin{Highlighting}[]
\NormalTok{autovalores\_log }\OtherTok{\textless{}{-}}\NormalTok{ acp\_log}\SpecialCharTok{$}\NormalTok{sdev}\SpecialCharTok{\^{}}\DecValTok{2}
\FunctionTok{c}\NormalTok{(}\AttributeTok{negativos =} \FunctionTok{sum}\NormalTok{(num111 }\SpecialCharTok{\textless{}} \DecValTok{0}\NormalTok{),}
  \AttributeTok{acumulado\_80 =} \FunctionTok{which}\NormalTok{(}\FunctionTok{cumsum}\NormalTok{(autovalores\_log) }\SpecialCharTok{/} \DecValTok{57} \SpecialCharTok{\textgreater{}=} \FloatTok{0.8}\NormalTok{)[}\DecValTok{1}\NormalTok{],}
  \AttributeTok{kaiser =} \FunctionTok{sum}\NormalTok{(autovalores\_log }\SpecialCharTok{\textgreater{}} \DecValTok{1}\NormalTok{))}
\end{Highlighting}
\end{Shaded}

\begin{verbatim}
   negativos acumulado_80       kaiser 
          80            5            9 
\end{verbatim}

\Needspace{0.38in}

\begin{Shaded}
\begin{Highlighting}[]
\FunctionTok{names}\NormalTok{(}\FunctionTok{which}\NormalTok{(}\FunctionTok{colSums}\NormalTok{(num111 }\SpecialCharTok{\textless{}} \DecValTok{0}\NormalTok{) }\SpecialCharTok{\textgreater{}} \DecValTok{0}\NormalTok{))}
\end{Highlighting}
\end{Shaded}

\begin{verbatim}
[1] "PetLiq_prod"  "GananciaRef"  "Elec_impNeta" "Bombeo_gen"  
\end{verbatim}

\Needspace{0.76in}

\begin{Shaded}
\begin{Highlighting}[]
\NormalTok{aporte\_log }\OtherTok{\textless{}{-}} \FunctionTok{get\_pca\_ind}\NormalTok{(acp\_log)}\SpecialCharTok{$}\NormalTok{contrib[, }\DecValTok{1}\NormalTok{]}
\FunctionTok{round}\NormalTok{(}\FunctionTok{head}\NormalTok{(}\FunctionTok{sort}\NormalTok{(aporte\_log, }\AttributeTok{decreasing =} \ConstantTok{TRUE}\NormalTok{), }\DecValTok{5}\NormalTok{), }\DecValTok{2}\NormalTok{)}
\end{Highlighting}
\end{Shaded}

\begin{verbatim}
UnitedStates        China        India        Japan       Brazil 
       12.48        10.27         4.60         3.89         3.15 
\end{verbatim}

\Needspace{0.38in}

\begin{Shaded}
\begin{Highlighting}[]
\FunctionTok{round}\NormalTok{(}\FunctionTok{sum}\NormalTok{(aporte\_log[}\FunctionTok{c}\NormalTok{(}\StringTok{"UnitedStates"}\NormalTok{, }\StringTok{"China"}\NormalTok{)]), }\DecValTok{2}\NormalTok{)}
\end{Highlighting}
\end{Shaded}

\begin{verbatim}
[1] 22.75
\end{verbatim}

\Needspace{1.33in}

\begin{Shaded}
\begin{Highlighting}[]
\FunctionTok{round}\NormalTok{(}\FunctionTok{get\_pca\_var}\NormalTok{(acp\_log)}\SpecialCharTok{$}\NormalTok{coord[cinco, }\DecValTok{1}\SpecialCharTok{:}\DecValTok{2}\NormalTok{], }\DecValTok{2}\NormalTok{)}
\end{Highlighting}
\end{Shaded}

\begin{verbatim}
            Dim.1 Dim.2
Crudo_prod   0.45 -0.69
PetLiq_cons  0.95 -0.21
Crudo_exp    0.27 -0.79
Elec_gen     0.95 -0.05
Renov_gen    0.75  0.33
\end{verbatim}

\begin{notas}

\textbf{Interpretaciones.}

\begin{enumerate}
\def\labelenumi{\alph{enumi}.}
\tightlist
\item
  \textbf{Qué cambió.} El primer eje pasa de 61.96 \% a 48.95 \%, las
  variables ligadas a él de 39 a 27 y los países que lo construyen de
  Estados Unidos y China (83 \%) a India y Japón (43 \%).
\item
  \textbf{Qué no cambió.} El primer eje sigue siendo tamaño y el segundo
  sigue siendo petróleo, solo que más nítido.
\item
  \textbf{Quitar no es corregir.} Los cuatro países no son errores de
  medición sino datos reales, de modo que excluirlos cambia la población
  que se describe: Los resultados de las preguntas 7 a 11 valen para 107
  países sin los cuatro mayores.
\end{enumerate}

\textbf{Decisiones.}

\begin{itemize}
\item
  \textbf{Se quitan los cuatro países que pide el enunciado y no seis.}
  El enunciado nombra a Estados Unidos, China, Arabia Saudita y Rusia y
  la distancia de Mahalanobis lo respalda pues sobre las dos y las tres
  primeras componentes de los 111 países (la covarianza de las 57
  variables es singular, de rango 50) son exactamente esos cuatro los
  que superan el corte de la \(\chi^2\) (7.38 y 9.35), mientras que
  Japón queda quinto con 5.14 y 6.30 e India sexta y séptima con 4.09 y
  5.36, ambos por debajo. India no supera el corte con ningún número de
  componentes (5.41 con cuatro y 15.43 con ocho frente a 11.14 y 17.53)
  y Japón solo lo supera desde la cuarta componente (15.95), que
  construyen sobre todo las mareas (0.50 y 0.47) y con la que también
  salen Francia, Corea del Sur e Italia, tres de los seis países con
  generación por mareas. Es decir que quitar seis sería tan arbitrario
  como quitar ocho pues la lejanía de Japón con cuatro componentes no es
  la del tamaño que motivó la exclusión, de modo que India y Japón
  heredan el primer eje (43 \%) sin ser atípicos por este criterio.

  \emph{Respaldo.} En la sesión 3 se detectaron los atípicos
  multivariados con la distancia de Mahalanobis frente al cuantil 0.975
  de una \(\chi^2\) con tantos grados de libertad como variables (el
  ejemplo de clase, con cuatro variables, daba el mismo corte de 11.14),
  que es el procedimiento de \textcite[\S~4.7]{johnson2007}, mientras
  que en la sesión 4 se recomendó calcularla antes del ACP porque un
  punto lejano «va a jalonar» la dirección de máxima varianza.
\item
  \textbf{El logaritmo entra como prueba de robustez y no reemplaza la
  exclusión.} Con el logaritmo con signo, necesario porque cuatro
  variables (producción de petróleo, ganancia de refinación, importación
  neta de electricidad y almacenamiento por bombeo) traen 80 valores
  negativos, el ACP normado de los 111 países explica el 55.81 \% en el
  primer eje y el 11.83 \% en el segundo, mientras que Estados Unidos y
  China pasan de construir el 83 \% del primer eje a construir el 22.75
  \% (12.48 \% y 10.27 \%), que es además menos que el 43 \% que India y
  Japón construyen sin los cuatro. La estructura en cambio se conserva
  ya que el primer eje sigue siendo tamaño (consumo y generación con
  0.95) y el segundo petróleo (exportación de crudo con \(-0.79\) y
  producción con \(-0.69\)), con lo cual la regla del 80 \% pide cinco
  componentes, entre las cuatro de los 111 y las seis de los 107. Se
  queda la exclusión porque la Pregunta 7 la pide y las preguntas 8 a 11
  se plantean sobre los 107 países, pero el logaritmo sugiere que el
  dominio de dos países es un efecto de la escala de las variables más
  que de la estructura de la base y que con esta transformación los
  resultados podrían extenderse a los 111 países.

  \emph{Respaldo.} En la sesión 4 se advirtió que el ACP está «altamente
  influenciado por valores atípicos», mientras que reexpresar con el
  logaritmo las variables positivas muy asimétricas es el remedio
  clásico del análisis exploratorio \parencite{tukey1977} y una
  transformación usual antes del ACP
  \parencite[\S~14.2.3]{jolliffe2002}.
\end{itemize}

\end{notas}

\chapter{Agrupamiento}\label{agrupamiento}

Como pide el enunciado, las preguntas 8 a 11 se resuelven sin Estados
Unidos, China, Arabia Saudita y Rusia, es decir, sobre los 107 países
que quedaron en \texttt{datos} y \texttt{num} al final de la Pregunta 7.

\section{Pregunta 8 --- Número de grupos en el primer plano}\label{p8}

\textbf{¿Cuántos clusters se espera que existan?} Tres, leídos sobre el
plano (1,2) de la Pregunta 7:

\begin{enumerate}
\def\labelenumi{\alph{enumi}.}
\tightlist
\item
  \textbf{Masa central.} La mayoría de los 107 países, de consumo y
  producción pequeños, que se amontonan cerca del origen.
\item
  \textbf{Brazo de consumidores.} Las economías grandes sin producción
  importante de petróleo hacia arriba y a la derecha, con Japón e India
  en el extremo.
\item
  \textbf{Brazo de productores.} Los exportadores de petróleo hacia
  abajo, con Brasil y Canadá en el extremo.
\end{enumerate}

Cabe aclarar que esta es una estimación a ojo sobre un plano que solo
explica el 59.19 \% de la varianza, que los brazos salen de la masa
central sin un corte claro y que India, Japón, Brasil y Canadá están tan
alejados que un método de agrupamiento podría separarlos como grupos
propios. Por esto el número se contrasta en la Pregunta 9.

\begin{enumerate}
\def\labelenumi{\alph{enumi}.}
\tightlist
\item
  \textbf{Lectura local.} Los cuatro países de los extremos (India,
  Japón, Brasil y Canadá) son los que más construyen el primer eje (de
  23.37 \% a 6.38 \%), de modo que pueden terminar como grupos propios
  aunque a ojo parezcan la punta de un brazo.
\item
  \textbf{Lectura global.} El 40.81 \% de la varianza no está en este
  plano, por lo que dos países que se ven cerca pueden estar lejos en
  otras componentes y por eso el número de grupos se lee aquí como una
  hipótesis y no como un resultado.
\end{enumerate}

\section{Pregunta 9 --- Número de grupos por reducción de varianza
explicada}\label{p9}

La inercia intra-grupo \(W(k)\) es la suma de las distancias al cuadrado
de cada país al centroide de su grupo, que mide qué tan compactos son
los grupos y la varianza explicada por la partición es

\[
1 - \frac{W(k)}{T}, \qquad
T = W(1) = \sum_{i=1}^{n} \lVert \mathbf{z}_i - \bar{\mathbf{z}} \rVert^2,
\]

es decir, la fracción de la inercia total \(T\) que queda entre los
grupos. Con las 57 variables estandarizadas
\(T = (n - 1)\,p = 106 \times 57 = 6042\). Se corre k-medias
\parencite{macqueen1967} para \(k = 1, \dots, 10\) y se busca el punto
desde el que agregar un grupo ya no reduce mucho la inercia. Como
verificación adicional se calcula la silueta media, que se define en la
Pregunta 10.

\Needspace{4.56in}

\begin{Shaded}
\begin{Highlighting}[]
\NormalTok{z }\OtherTok{\textless{}{-}} \FunctionTok{scale}\NormalTok{(num)}
\FunctionTok{set.seed}\NormalTok{(}\DecValTok{42}\NormalTok{)}
\NormalTok{ajustes }\OtherTok{\textless{}{-}} \FunctionTok{lapply}\NormalTok{(}\DecValTok{1}\SpecialCharTok{:}\DecValTok{10}\NormalTok{, }\ControlFlowTok{function}\NormalTok{(k)}
  \FunctionTok{kmeans}\NormalTok{(z, }\AttributeTok{centers =}\NormalTok{ k, }\AttributeTok{nstart =} \DecValTok{100}\NormalTok{))}
\NormalTok{inercia }\OtherTok{\textless{}{-}} \FunctionTok{sapply}\NormalTok{(ajustes, }\StringTok{\textasciigrave{}}\AttributeTok{[[}\StringTok{\textasciigrave{}}\NormalTok{, }\StringTok{"tot.withinss"}\NormalTok{)}
\NormalTok{explicada }\OtherTok{\textless{}{-}} \DecValTok{100} \SpecialCharTok{*}\NormalTok{ (}\DecValTok{1} \SpecialCharTok{{-}}\NormalTok{ inercia }\SpecialCharTok{/}\NormalTok{ inercia[}\DecValTok{1}\NormalTok{])}
\NormalTok{silueta }\OtherTok{\textless{}{-}} \FunctionTok{c}\NormalTok{(}\ConstantTok{NA}\NormalTok{, }\FunctionTok{sapply}\NormalTok{(ajustes[}\SpecialCharTok{{-}}\DecValTok{1}\NormalTok{], }\ControlFlowTok{function}\NormalTok{(a)}
  \FunctionTok{mean}\NormalTok{(}\FunctionTok{silhouette}\NormalTok{(a}\SpecialCharTok{$}\NormalTok{cluster, }\FunctionTok{dist}\NormalTok{(z))[, }\DecValTok{3}\NormalTok{])))}
\NormalTok{tamanos }\OtherTok{\textless{}{-}} \FunctionTok{sapply}\NormalTok{(ajustes, }\ControlFlowTok{function}\NormalTok{(a)}
  \FunctionTok{paste}\NormalTok{(}\FunctionTok{sort}\NormalTok{(a}\SpecialCharTok{$}\NormalTok{size, }\AttributeTok{decreasing =} \ConstantTok{TRUE}\NormalTok{), }\AttributeTok{collapse =} \StringTok{"/"}\NormalTok{))}
\NormalTok{knitr}\SpecialCharTok{::}\FunctionTok{kable}\NormalTok{(}\FunctionTok{data.frame}\NormalTok{(}\AttributeTok{k =} \DecValTok{1}\SpecialCharTok{:}\DecValTok{10}\NormalTok{, }\AttributeTok{W =} \FunctionTok{round}\NormalTok{(inercia, }\DecValTok{1}\NormalTok{),}
                        \AttributeTok{explicada =} \FunctionTok{round}\NormalTok{(explicada, }\DecValTok{1}\NormalTok{),}
                        \AttributeTok{ganancia =} \FunctionTok{round}\NormalTok{(}\FunctionTok{c}\NormalTok{(}\ConstantTok{NA}\NormalTok{, }\FunctionTok{diff}\NormalTok{(explicada)), }\DecValTok{1}\NormalTok{),}
                        \AttributeTok{silueta =} \FunctionTok{round}\NormalTok{(silueta, }\DecValTok{3}\NormalTok{), tamanos),}
             \AttributeTok{col.names =} \FunctionTok{c}\NormalTok{(}\StringTok{"k"}\NormalTok{, }\StringTok{"W"}\NormalTok{, }\StringTok{"explicada"}\NormalTok{, }\StringTok{"ganancia"}\NormalTok{,}
                           \StringTok{"silueta"}\NormalTok{, }\StringTok{"tamaños"}\NormalTok{),}
             \AttributeTok{caption =} \StringTok{"Inercia intra{-}grupo y varianza explicada."}\NormalTok{)}
\end{Highlighting}
\end{Shaded}

\Needspace{3.07in}

\begin{longtable}[]{@{}rrrrrl@{}}
\caption{Inercia intra-grupo y varianza explicada.}\tabularnewline
\toprule\noalign{}
k & W & explicada & ganancia & silueta & tamaños \\
\midrule\noalign{}
\endfirsthead
\toprule\noalign{}
k & W & explicada & ganancia & silueta & tamaños \\
\midrule\noalign{}
\endhead
\bottomrule\noalign{}
\endlastfoot
1 & 6042.0 & 0.0 & NA & NA & 107 \\
2 & 3838.6 & 36.5 & 36.5 & 0.698 & 98/9 \\
3 & 3226.7 & 46.6 & 10.1 & 0.666 & 97/8/2 \\
4 & 2767.7 & 54.2 & 7.6 & 0.538 & 87/10/8/2 \\
5 & 2303.1 & 61.9 & 7.7 & 0.546 & 87/10/6/2/2 \\
6 & 2023.6 & 66.5 & 4.6 & 0.552 & 87/10/4/2/2/2 \\
7 & 1828.8 & 69.7 & 3.2 & 0.444 & 76/15/6/4/2/2/2 \\
8 & 1615.5 & 73.3 & 3.5 & 0.429 & 73/17/6/4/2/2/2/1 \\
9 & 1468.9 & 75.7 & 2.4 & 0.412 & 71/17/5/4/3/2/2/2/1 \\
10 & 1376.4 & 77.2 & 1.5 & 0.391 & 65/16/7/5/4/3/2/2/2/1 \\
\end{longtable}

\Needspace{5.02in}

\begin{Shaded}
\begin{Highlighting}[]
\FunctionTok{ggplot}\NormalTok{(}\FunctionTok{data.frame}\NormalTok{(}\AttributeTok{k =} \DecValTok{1}\SpecialCharTok{:}\DecValTok{10}\NormalTok{, explicada), }\FunctionTok{aes}\NormalTok{(k, explicada)) }\SpecialCharTok{+}
  \FunctionTok{geom\_line}\NormalTok{(}\AttributeTok{colour =} \StringTok{"grey50"}\NormalTok{) }\SpecialCharTok{+}
  \FunctionTok{geom\_point}\NormalTok{(}\AttributeTok{colour =} \StringTok{"\#a6192e"}\NormalTok{, }\AttributeTok{size =} \DecValTok{2}\NormalTok{) }\SpecialCharTok{+}
  \FunctionTok{geom\_vline}\NormalTok{(}\AttributeTok{xintercept =} \DecValTok{3}\NormalTok{, }\AttributeTok{linetype =} \StringTok{"dashed"}\NormalTok{,}
             \AttributeTok{colour =} \StringTok{"grey50"}\NormalTok{) }\SpecialCharTok{+}
  \FunctionTok{scale\_x\_continuous}\NormalTok{(}\AttributeTok{breaks =} \DecValTok{1}\SpecialCharTok{:}\DecValTok{10}\NormalTok{) }\SpecialCharTok{+}
  \FunctionTok{labs}\NormalTok{(}\AttributeTok{x =} \StringTok{"Número de grupos k"}\NormalTok{, }\AttributeTok{y =} \StringTok{"Varianza explicada (\%)"}\NormalTok{) }\SpecialCharTok{+}
  \FunctionTok{theme\_minimal}\NormalTok{()}
\end{Highlighting}
\end{Shaded}

\begin{figure}[H]

{\centering \pandocbounded{\includegraphics[keepaspectratio]{sda-lab-exploracion_files/figure-latex/p9-curva-1.pdf}}

}

\caption{Varianza explicada por la partición de k-medias según el número
de grupos: La ganancia cae de 36.5 a 10.1 puntos al pasar al tercer
grupo y desde ahí se mantiene por debajo de 8.}

\end{figure}%

\textbf{¿Cuántos grupos se fijan?} Tres, que explican el 46.6 \% de la
inercia total. La reducción más grande ocurre al pasar de uno a dos
grupos (36.5 puntos), el tercer grupo todavía agrega 10.1 puntos y a
partir de ahí cada grupo nuevo agrega 7.7 puntos o menos, cayendo a 4.6
en \(k = 6\) y a 3.2 en \(k = 7\). Vale la pena anotar con honestidad
que la curva no tiene un codo marcado después de \(k = 2\), sino una
caída gradual, por lo que la elección se apoya además en dos hechos:

\begin{enumerate}
\def\labelenumi{\alph{enumi}.}
\tightlist
\item
  \textbf{Silueta.} La silueta media pasa de 0.698 con dos grupos a
  0.666 con tres, pero cae a 0.538 con cuatro, de modo que el cuarto
  grupo empeora claramente la separación.
\item
  \textbf{Coherencia con el plano.} Tres grupos es lo que se esperaba a
  ojo en la Pregunta 8, mientras que dos grupos solo separarían a nueve
  países grandes del resto (partición 98/9).
\end{enumerate}

La elección también se lee en dos niveles:

\begin{enumerate}
\def\labelenumi{\alph{enumi}.}
\tightlist
\item
  \textbf{Lectura local.} El tercer grupo cuesta poco y aporta mucho
  pues con 10.1 puntos más de varianza explicada la silueta solo baja de
  0.698 a 0.666, mientras que el cuarto parte la masa central
  (87/10/8/2) y la silueta cae a 0.538.
\item
  \textbf{Lectura global.} La curva sin codo después de \(k = 2\) indica
  que la base no tiene grupos naturales bien separados sino una masa
  continua con unos pocos países extremos, de modo que cualquier \(k\)
  es en parte una decisión y no un descubrimiento.
\end{enumerate}

\section{Pregunta 10 --- K-medias frente a jerárquico de enlace
promedio}\label{p10}

Se comparan dos formas de partir los 107 países en tres grupos con las
57 variables estandarizadas:

\begin{enumerate}
\def\labelenumi{\alph{enumi}.}
\tightlist
\item
  \textbf{K-medias.} Asigna cada país al centroide más cercano y mueve
  los centroides hasta minimizar \(W(k)\). Se usan 100 arranques
  aleatorios (\texttt{nstart\ =\ 100}) con semilla fija, por la razón
  que explica la nota del final.
\item
  \textbf{Enlace promedio.} Es un agrupamiento jerárquico aglomerativo
  que parte de un grupo por país y une en cada paso los dos grupos con
  menor distancia promedio entre todos sus pares de países
  \parencite{sokal1958}. Se corta el árbol en tres grupos.
\item
  \textbf{Correlación cofenética.} Es la correlación entre las
  distancias originales y la altura del árbol a la que se unen cada par
  de países, y mide qué tan fielmente el dendrograma conserva las
  distancias.
\item
  \textbf{Silueta.} Para cada país es
  \(s(i) = (b_i - a_i)/\max(a_i, b_i)\), con \(a_i\) su distancia
  promedio a su propio grupo y \(b_i\) al grupo vecino más cercano
  \parencite{rousseeuw1987}. Vale 1 si el país está bien ubicado, 0 si
  está en la frontera y es negativa si estaría mejor en otro grupo.
\end{enumerate}

\Needspace{2.04in}

\includegraphics[width=5.2in,height=1.29in]{sda-lab-exploracion_files/figure-latex/mermaid-figure-2.png}

Qué optimiza cada método sobre los 107 países: K-medias busca grupos
compactos y arma uno intermedio, mientras que el enlace promedio deja
para el final a los países lejanos y los aísla en parejas.

\Needspace{3.61in}

\begin{Shaded}
\begin{Highlighting}[]
\FunctionTok{set.seed}\NormalTok{(}\DecValTok{42}\NormalTok{)}
\NormalTok{km3 }\OtherTok{\textless{}{-}} \FunctionTok{kmeans}\NormalTok{(z, }\AttributeTok{centers =} \DecValTok{3}\NormalTok{, }\AttributeTok{nstart =} \DecValTok{100}\NormalTok{)}
\NormalTok{distancias }\OtherTok{\textless{}{-}} \FunctionTok{dist}\NormalTok{(z)}
\NormalTok{hc\_average }\OtherTok{\textless{}{-}} \FunctionTok{hclust}\NormalTok{(distancias, }\AttributeTok{method =} \StringTok{"average"}\NormalTok{)}
\NormalTok{jerarquico }\OtherTok{\textless{}{-}} \FunctionTok{cutree}\NormalTok{(hc\_average, }\AttributeTok{k =} \DecValTok{3}\NormalTok{)}
\CommentTok{\# Grupos rotulados por tamaño: A el más grande. Los números que}
\CommentTok{\# devuelven kmeans() y cutree() son arbitrarios.}
\NormalTok{rotular }\OtherTok{\textless{}{-}} \ControlFlowTok{function}\NormalTok{(g) \{}
\NormalTok{  orden }\OtherTok{\textless{}{-}} \FunctionTok{rank}\NormalTok{(}\SpecialCharTok{{-}}\FunctionTok{tabulate}\NormalTok{(g), }\AttributeTok{ties.method =} \StringTok{"first"}\NormalTok{)}
  \FunctionTok{factor}\NormalTok{(LETTERS[orden[g]])}
\NormalTok{\}}
\NormalTok{kmedias }\OtherTok{\textless{}{-}} \FunctionTok{rotular}\NormalTok{(km3}\SpecialCharTok{$}\NormalTok{cluster)}
\NormalTok{promedio }\OtherTok{\textless{}{-}} \FunctionTok{rotular}\NormalTok{(jerarquico)}
\FunctionTok{table}\NormalTok{(kmedias, promedio)}
\end{Highlighting}
\end{Shaded}

\begin{verbatim}
       promedio
kmedias  A  B  C
      A 97  0  0
      B  6  2  0
      C  0  0  2
\end{verbatim}

\Needspace{1.33in}

\begin{Shaded}
\begin{Highlighting}[]
\FunctionTok{split}\NormalTok{(}\FunctionTok{rownames}\NormalTok{(num), kmedias)[}\FunctionTok{c}\NormalTok{(}\StringTok{"B"}\NormalTok{, }\StringTok{"C"}\NormalTok{)]}
\end{Highlighting}
\end{Shaded}

\begin{verbatim}
$B
[1] "Brazil"        "Canada"        "France"       
[4] "Germany"       "Italy"         "SouthKorea"   
[7] "Spain"         "UnitedKingdom"

$C
[1] "India" "Japan"
\end{verbatim}

\Needspace{0.95in}

\begin{Shaded}
\begin{Highlighting}[]
\FunctionTok{split}\NormalTok{(}\FunctionTok{rownames}\NormalTok{(num), promedio)[}\FunctionTok{c}\NormalTok{(}\StringTok{"B"}\NormalTok{, }\StringTok{"C"}\NormalTok{)]}
\end{Highlighting}
\end{Shaded}

\begin{verbatim}
$B
[1] "Brazil" "Canada"

$C
[1] "India" "Japan"
\end{verbatim}

\Needspace{3.42in}

\begin{Shaded}
\begin{Highlighting}[]
\NormalTok{explicada }\OtherTok{\textless{}{-}} \ControlFlowTok{function}\NormalTok{(g) }\DecValTok{1} \SpecialCharTok{{-}} \FunctionTok{sum}\NormalTok{(}\FunctionTok{sapply}\NormalTok{(}\FunctionTok{split}\NormalTok{(}\FunctionTok{as.data.frame}\NormalTok{(z), g),}
  \ControlFlowTok{function}\NormalTok{(b) }\FunctionTok{sum}\NormalTok{(}\FunctionTok{scale}\NormalTok{(b, }\AttributeTok{scale =} \ConstantTok{FALSE}\NormalTok{)}\SpecialCharTok{\^{}}\DecValTok{2}\NormalTok{))) }\SpecialCharTok{/} \FunctionTok{sum}\NormalTok{(z}\SpecialCharTok{\^{}}\DecValTok{2}\NormalTok{)}
\NormalTok{sil }\OtherTok{\textless{}{-}} \ControlFlowTok{function}\NormalTok{(g) }\FunctionTok{summary}\NormalTok{(}\FunctionTok{silhouette}\NormalTok{(}\FunctionTok{as.integer}\NormalTok{(g), distancias))}
\NormalTok{medidas }\OtherTok{\textless{}{-}} \ControlFlowTok{function}\NormalTok{(g) }\FunctionTok{c}\NormalTok{(}\FunctionTok{explicada}\NormalTok{(g), }\FunctionTok{sil}\NormalTok{(g)}\SpecialCharTok{$}\NormalTok{avg.width,}
                         \FunctionTok{sil}\NormalTok{(g)}\SpecialCharTok{$}\NormalTok{clus.avg.widths)}
\NormalTok{knitr}\SpecialCharTok{::}\FunctionTok{kable}\NormalTok{(}\FunctionTok{round}\NormalTok{(}\FunctionTok{rbind}\NormalTok{(}\StringTok{"k{-}medias"} \OtherTok{=} \FunctionTok{medidas}\NormalTok{(kmedias),}
                         \StringTok{"enlace promedio"} \OtherTok{=} \FunctionTok{medidas}\NormalTok{(promedio)), }\DecValTok{3}\NormalTok{),}
             \AttributeTok{col.names =} \FunctionTok{c}\NormalTok{(}\StringTok{"Explicada"}\NormalTok{, }\StringTok{"Silueta media"}\NormalTok{,}
                           \StringTok{"Silueta A"}\NormalTok{, }\StringTok{"Silueta B"}\NormalTok{, }\StringTok{"Silueta C"}\NormalTok{),}
             \AttributeTok{caption =} \StringTok{"Varianza explicada y silueta."}\NormalTok{)}
\end{Highlighting}
\end{Shaded}

\Needspace{1.55in}

\begin{longtable}[]{@{}
  >{\raggedright\arraybackslash}p{(\linewidth - 10\tabcolsep) * \real{0.2286}}
  >{\raggedleft\arraybackslash}p{(\linewidth - 10\tabcolsep) * \real{0.1429}}
  >{\raggedleft\arraybackslash}p{(\linewidth - 10\tabcolsep) * \real{0.2000}}
  >{\raggedleft\arraybackslash}p{(\linewidth - 10\tabcolsep) * \real{0.1429}}
  >{\raggedleft\arraybackslash}p{(\linewidth - 10\tabcolsep) * \real{0.1429}}
  >{\raggedleft\arraybackslash}p{(\linewidth - 10\tabcolsep) * \real{0.1429}}@{}}
\caption{Varianza explicada y silueta.}\tabularnewline
\toprule\noalign{}
\begin{minipage}[b]{\linewidth}\raggedright
\end{minipage} & \begin{minipage}[b]{\linewidth}\raggedleft
Explicada
\end{minipage} & \begin{minipage}[b]{\linewidth}\raggedleft
Silueta media
\end{minipage} & \begin{minipage}[b]{\linewidth}\raggedleft
Silueta A
\end{minipage} & \begin{minipage}[b]{\linewidth}\raggedleft
Silueta B
\end{minipage} & \begin{minipage}[b]{\linewidth}\raggedleft
Silueta C
\end{minipage} \\
\midrule\noalign{}
\endfirsthead
\toprule\noalign{}
\begin{minipage}[b]{\linewidth}\raggedright
\end{minipage} & \begin{minipage}[b]{\linewidth}\raggedleft
Explicada
\end{minipage} & \begin{minipage}[b]{\linewidth}\raggedleft
Silueta media
\end{minipage} & \begin{minipage}[b]{\linewidth}\raggedleft
Silueta A
\end{minipage} & \begin{minipage}[b]{\linewidth}\raggedleft
Silueta B
\end{minipage} & \begin{minipage}[b]{\linewidth}\raggedleft
Silueta C
\end{minipage} \\
\midrule\noalign{}
\endhead
\bottomrule\noalign{}
\endlastfoot
k-medias & 0.466 & 0.666 & 0.731 & -0.022 & 0.254 \\
enlace promedio & 0.390 & 0.713 & 0.732 & 0.174 & 0.266 \\
\end{longtable}

\Needspace{1.14in}

\begin{codres}

\begin{Shaded}
\begin{Highlighting}[]
\FunctionTok{round}\NormalTok{(}\FunctionTok{cor}\NormalTok{(}\FunctionTok{cophenetic}\NormalTok{(}
\NormalTok{  hc\_average), distancias), }\DecValTok{3}\NormalTok{)}
\end{Highlighting}
\end{Shaded}

\codressep

\begin{verbatim}
[1] 0.986
\end{verbatim}

\end{codres}

\Needspace{7.72in}

\begin{Shaded}
\begin{Highlighting}[]
\CommentTok{\# Con 107 hojas los nombres no se leen: se rotulan los extremos.}
\FunctionTok{par}\NormalTok{(}\AttributeTok{mar =} \FunctionTok{c}\NormalTok{(}\DecValTok{4}\NormalTok{, }\DecValTok{4}\NormalTok{, }\DecValTok{1}\NormalTok{, }\DecValTok{1}\NormalTok{))}
\FunctionTok{plot}\NormalTok{(hc\_average, }\AttributeTok{labels =} \ConstantTok{FALSE}\NormalTok{, }\AttributeTok{hang =} \SpecialCharTok{{-}}\DecValTok{1}\NormalTok{, }\AttributeTok{main =} \StringTok{""}\NormalTok{, }\AttributeTok{xlab =} \StringTok{""}\NormalTok{,}
     \AttributeTok{sub =} \StringTok{""}\NormalTok{, }\AttributeTok{ylab =} \StringTok{"Altura (distancia promedio)"}\NormalTok{)}
\FunctionTok{rect.hclust}\NormalTok{(hc\_average, }\AttributeTok{k =} \DecValTok{3}\NormalTok{, }\AttributeTok{border =} \FunctionTok{c}\NormalTok{(}\StringTok{"\#a6192e"}\NormalTok{, }\StringTok{"\#f08a3c"}\NormalTok{,}
                                          \StringTok{"\#3572a8"}\NormalTok{))}
\CommentTok{\# Cada par se rotula con una línea a la altura a la que se une.}
\NormalTok{hojas }\OtherTok{\textless{}{-}}\NormalTok{ hc\_average}\SpecialCharTok{$}\NormalTok{labels[hc\_average}\SpecialCharTok{$}\NormalTok{order]}
\NormalTok{pares }\OtherTok{\textless{}{-}} \FunctionTok{list}\NormalTok{(}\StringTok{"India y Japón"} \OtherTok{=} \FunctionTok{c}\NormalTok{(}\StringTok{"India"}\NormalTok{, }\StringTok{"Japan"}\NormalTok{),}
              \StringTok{"Brasil y Canadá"} \OtherTok{=} \FunctionTok{c}\NormalTok{(}\StringTok{"Brazil"}\NormalTok{, }\StringTok{"Canada"}\NormalTok{))}
\NormalTok{rotulo\_y }\OtherTok{\textless{}{-}} \FunctionTok{c}\NormalTok{(}\FloatTok{27.5}\NormalTok{, }\FloatTok{25.5}\NormalTok{)}
\ControlFlowTok{for}\NormalTok{ (i }\ControlFlowTok{in} \DecValTok{1}\SpecialCharTok{:}\DecValTok{2}\NormalTok{) \{}
\NormalTok{  par\_i }\OtherTok{\textless{}{-}}\NormalTok{ pares[[i]]}
\NormalTok{  x }\OtherTok{\textless{}{-}} \FunctionTok{mean}\NormalTok{(}\FunctionTok{match}\NormalTok{(par\_i, hojas))}
\NormalTok{  y }\OtherTok{\textless{}{-}} \FunctionTok{as.matrix}\NormalTok{(}\FunctionTok{cophenetic}\NormalTok{(hc\_average))[par\_i[}\DecValTok{1}\NormalTok{], par\_i[}\DecValTok{2}\NormalTok{]]}
  \FunctionTok{segments}\NormalTok{(x, y, }\DecValTok{12}\NormalTok{, rotulo\_y[i], }\AttributeTok{col =} \StringTok{"grey40"}\NormalTok{)}
  \FunctionTok{text}\NormalTok{(}\DecValTok{12}\NormalTok{, rotulo\_y[i], }\FunctionTok{names}\NormalTok{(pares)[i], }\AttributeTok{pos =} \DecValTok{4}\NormalTok{, }\AttributeTok{cex =} \FloatTok{0.7}\NormalTok{)}
\NormalTok{\}}
\end{Highlighting}
\end{Shaded}

\begin{figure}[H]

{\centering \pandocbounded{\includegraphics[keepaspectratio]{sda-lab-exploracion_files/figure-latex/p10-dendrograma-1.pdf}}

}

\caption{Dendrograma del enlace promedio con el corte en tres grupos,
rotulando solo los cuatro países que quedan fuera de la rama grande:
India y Japón, y Brasil y Canadá, se separan arriba del árbol mientras
que los otros 103 países se unen en una sola rama.}

\end{figure}%

\Needspace{3.42in}

\begin{Shaded}
\begin{Highlighting}[]
\NormalTok{knitr}\SpecialCharTok{::}\FunctionTok{kable}\NormalTok{(}\FunctionTok{data.frame}\NormalTok{(}
  \AttributeTok{k =} \DecValTok{2}\SpecialCharTok{:}\DecValTok{5}\NormalTok{,}
  \AttributeTok{tamanos =} \FunctionTok{sapply}\NormalTok{(}\DecValTok{2}\SpecialCharTok{:}\DecValTok{5}\NormalTok{, }\ControlFlowTok{function}\NormalTok{(k) }\FunctionTok{paste}\NormalTok{(}\FunctionTok{sort}\NormalTok{(}\FunctionTok{table}\NormalTok{(}
    \FunctionTok{cutree}\NormalTok{(hc\_average, k)), }\AttributeTok{decreasing =} \ConstantTok{TRUE}\NormalTok{), }\AttributeTok{collapse =} \StringTok{"/"}\NormalTok{)),}
  \AttributeTok{silueta =} \FunctionTok{sapply}\NormalTok{(}\DecValTok{2}\SpecialCharTok{:}\DecValTok{5}\NormalTok{, }\ControlFlowTok{function}\NormalTok{(k)}
    \FunctionTok{round}\NormalTok{(}\FunctionTok{sil}\NormalTok{(}\FunctionTok{cutree}\NormalTok{(hc\_average, k))}\SpecialCharTok{$}\NormalTok{avg.width, }\DecValTok{3}\NormalTok{))),}
  \AttributeTok{col.names =} \FunctionTok{c}\NormalTok{(}\StringTok{"k"}\NormalTok{, }\StringTok{"tamaños"}\NormalTok{, }\StringTok{"silueta"}\NormalTok{),}
  \AttributeTok{caption =} \StringTok{"Cortes del árbol de enlace promedio."}\NormalTok{)}
\end{Highlighting}
\end{Shaded}

\Needspace{1.93in}

\begin{longtable}[]{@{}rlr@{}}
\caption{Cortes del árbol de enlace promedio.}\tabularnewline
\toprule\noalign{}
k & tamaños & silueta \\
\midrule\noalign{}
\endfirsthead
\toprule\noalign{}
k & tamaños & silueta \\
\midrule\noalign{}
\endhead
\bottomrule\noalign{}
\endlastfoot
2 & 105/2 & 0.752 \\
3 & 103/2/2 & 0.713 \\
4 & 101/2/2/2 & 0.686 \\
5 & 101/2/2/1/1 & 0.678 \\
\end{longtable}

\Needspace{4.69in}

\begin{Shaded}
\begin{Highlighting}[]
\FunctionTok{mapa\_paises}\NormalTok{(acp\_sin4, }\FunctionTok{c}\NormalTok{(}\DecValTok{1}\NormalTok{, }\DecValTok{2}\NormalTok{), }\AttributeTok{n =} \DecValTok{14}\NormalTok{, }\AttributeTok{grupos =}\NormalTok{ kmedias)}
\end{Highlighting}
\end{Shaded}

\begin{figure}[H]

{\centering \pandocbounded{\includegraphics[keepaspectratio]{sda-lab-exploracion_files/figure-latex/p10-mapa-1.pdf}}

}

\caption{Grupos de k-medias sobre el plano (1,2) del ACP sin los cuatro
países: El grupo B reúne las economías grandes de los dos brazos y el C
a India y Japón.}

\end{figure}%

\textbf{¿Son diferentes los resultados de los dos análisis?} Sí, aunque
coinciden en lo principal pues 101 de los 107 países quedan en el grupo
equivalente con ambos métodos:

\begin{enumerate}
\def\labelenumi{\alph{enumi}.}
\tightlist
\item
  \textbf{En qué coinciden.} Los dos aíslan a India y Japón como un
  grupo propio (grupo C) y dejan juntos a los mismos 97 países de la
  masa central.
\item
  \textbf{En qué difieren.} K-medias forma un grupo B de ocho economías
  grandes (Brasil, Canadá, Francia, Alemania, Italia, Corea del Sur,
  España y Reino Unido), mientras que el enlace promedio solo separa a
  Brasil y Canadá y devuelve los otros seis a la masa central, que queda
  de 103 países.
\end{enumerate}

Esta diferencia viene de lo que optimiza cada método. K-medias minimiza
la inercia intra-grupo y por eso arma un grupo intermedio que le da más
varianza explicada (46.6 \% frente a 39.0 \%), mientras que el enlace
promedio une primero lo que está cerca en promedio, de modo que solo
deja por fuera a los pares más alejados, tanto así que cortar el árbol
en cuatro o cinco grupos tampoco rescata a las seis economías
intermedias sino que separa otro par o países sueltos (particiones de
101/2/2/2 y 101/2/2/1/1). La correlación cofenética de 0.986 indica que
el árbol conserva muy bien las distancias, aunque cabe matizar que con
una masa tan compacta y cuatro países tan lejanos cualquier árbol tiende
a lograrlo. Por su parte la silueta media favorece al jerárquico (0.713
frente a 0.666), aunque la ventaja no se explica solo por el grupo B de
k-medias, cuya silueta promedio de \(-0.022\) indica que esos ocho
países están en promedio tan cerca de la masa central como entre sí,
sino sobre todo por la masa central pues las seis economías intermedias,
al volver a ella, pasan de una silueta media de \(-0.051\) a 0.262 y los
97 países que los dos métodos dejan juntos de 0.731 a 0.761, mientras
que las dos parejas aisladas por el enlace promedio solo promedian 0.220
(el cálculo está al final de la pregunta). En resumen ambos métodos ven
la misma estructura de una masa grande y pocos países extremos y lo que
no está bien definido es si las economías grandes intermedias forman un
grupo aparte.

\Needspace{1.52in}

\begin{Shaded}
\begin{Highlighting}[]
\NormalTok{inicios }\OtherTok{\textless{}{-}} \ControlFlowTok{function}\NormalTok{(n) }\FunctionTok{sapply}\NormalTok{(}\FunctionTok{c}\NormalTok{(}\DecValTok{1}\NormalTok{, }\DecValTok{42}\NormalTok{, }\DecValTok{2026}\NormalTok{), }\ControlFlowTok{function}\NormalTok{(s) \{}
  \FunctionTok{set.seed}\NormalTok{(s)}
  \FunctionTok{round}\NormalTok{(}\FunctionTok{kmeans}\NormalTok{(z, }\AttributeTok{centers =} \DecValTok{3}\NormalTok{, }\AttributeTok{nstart =}\NormalTok{ n)}\SpecialCharTok{$}\NormalTok{tot.withinss, }\DecValTok{1}\NormalTok{)}
\NormalTok{\})}
\FunctionTok{rbind}\NormalTok{(}\AttributeTok{nstart\_25 =} \FunctionTok{inicios}\NormalTok{(}\DecValTok{25}\NormalTok{), }\AttributeTok{nstart\_100 =} \FunctionTok{inicios}\NormalTok{(}\DecValTok{100}\NormalTok{))}
\end{Highlighting}
\end{Shaded}

\begin{verbatim}
             [,1]   [,2]   [,3]
nstart_25  3274.5 3231.9 3274.5
nstart_100 3226.7 3226.7 3226.7
\end{verbatim}

\begin{quote}
\textbf{Nota.} Con \texttt{nstart\ =\ 25} k-medias cae en óptimos
locales distintos según la semilla (inercias de 3274.5 o 3231.9, con
grupos de 90/14/3 y 95/10/2 países), mientras que con 100 arranques las
tres semillas llegan a la misma partición de inercia 3226.7, que es
también la que encuentra SDA Lab con 25 reinicios de k-medias++. Por eso
se usan 100 arranques.
\end{quote}

\begin{enumerate}
\def\labelenumi{\alph{enumi}.}
\tightlist
\item
  \textbf{Lectura local.} El grupo B es el punto débil de k-medias pues
  su silueta promedio es negativa (\(-0.022\)) y la de Alemania llega a
  \(-0.222\), es decir que Alemania está más cerca, en promedio, de la
  masa central que de los otros siete países de su grupo.
\item
  \textbf{Lectura global.} Los dos métodos responden a la misma
  estructura de la base, una masa compacta con unos pocos extremos y
  difieren solo en la franja intermedia, que es justo donde los datos no
  tienen una frontera clara y donde cada criterio de optimización decide
  por su cuenta.
\end{enumerate}

\Needspace{4.18in}

\begin{Shaded}
\begin{Highlighting}[]
\NormalTok{s }\OtherTok{\textless{}{-}} \FunctionTok{silhouette}\NormalTok{(}\FunctionTok{as.integer}\NormalTok{(kmedias), distancias)}
\NormalTok{b }\OtherTok{\textless{}{-}}\NormalTok{ kmedias }\SpecialCharTok{==} \StringTok{"B"}
\CommentTok{\# vecino: El grupo más cercano a cada país después del suyo.}
\NormalTok{silueta\_b }\OtherTok{\textless{}{-}} \FunctionTok{data.frame}\NormalTok{(}\AttributeTok{silueta =} \FunctionTok{round}\NormalTok{(s[b, }\StringTok{"sil\_width"}\NormalTok{], }\DecValTok{3}\NormalTok{),}
                        \AttributeTok{vecino =}\NormalTok{ LETTERS[s[b, }\StringTok{"neighbor"}\NormalTok{]],}
                        \AttributeTok{row.names =} \FunctionTok{rownames}\NormalTok{(num)[b])}
\NormalTok{knitr}\SpecialCharTok{::}\FunctionTok{kable}\NormalTok{(silueta\_b[}\FunctionTok{order}\NormalTok{(silueta\_b}\SpecialCharTok{$}\NormalTok{silueta), ],}
             \AttributeTok{caption =} \StringTok{"Silueta y grupo vecino del grupo B."}\NormalTok{)}
\end{Highlighting}
\end{Shaded}

\Needspace{2.69in}

\begin{longtable}[]{@{}lrl@{}}
\caption{Silueta y grupo vecino del grupo B.}\tabularnewline
\toprule\noalign{}
& silueta & vecino \\
\midrule\noalign{}
\endfirsthead
\toprule\noalign{}
& silueta & vecino \\
\midrule\noalign{}
\endhead
\bottomrule\noalign{}
\endlastfoot
Germany & -0.222 & A \\
UnitedKingdom & -0.092 & A \\
Spain & -0.054 & A \\
SouthKorea & -0.021 & A \\
Italy & 0.020 & A \\
Canada & 0.046 & A \\
France & 0.062 & A \\
Brazil & 0.082 & C \\
\end{longtable}

\Needspace{1.14in}

\begin{Shaded}
\begin{Highlighting}[]
\CommentTok{\# Si B se devuelve a la masa central quedan dos grupos: A+B y C.}
\NormalTok{sin\_b }\OtherTok{\textless{}{-}} \FunctionTok{droplevels}\NormalTok{(}\FunctionTok{replace}\NormalTok{(kmedias, b, }\StringTok{"A"}\NormalTok{))}
\FunctionTok{round}\NormalTok{(}\FunctionTok{c}\NormalTok{(}\AttributeTok{explicada =} \FunctionTok{explicada}\NormalTok{(sin\_b),}
        \AttributeTok{silueta =} \FunctionTok{sil}\NormalTok{(sin\_b)}\SpecialCharTok{$}\NormalTok{avg.width), }\DecValTok{3}\NormalTok{)}
\end{Highlighting}
\end{Shaded}

\begin{verbatim}
explicada   silueta 
    0.245     0.752 
\end{verbatim}

Para decidir qué partición defender se comparan con k-medias otras
cuatro particiones en tres grupos, dos de otros enlaces jerárquicos y
dos de k-medias sobre las primeras componentes del ACP de la Pregunta 7,
todas medidas sobre las 57 variables estandarizadas. Luego se descompone
la silueta de los dos métodos por bloques de países.

\Needspace{4.56in}

\begin{Shaded}
\begin{Highlighting}[]
\CommentTok{\# Índice de Rand ajustado: 1 si dos particiones coinciden y cerca}
\CommentTok{\# de 0 si coinciden lo que se esperaría al azar.}
\NormalTok{ari }\OtherTok{\textless{}{-}} \ControlFlowTok{function}\NormalTok{(a, b) \{}
\NormalTok{  t }\OtherTok{\textless{}{-}} \FunctionTok{table}\NormalTok{(a, b)}
\NormalTok{  pares }\OtherTok{\textless{}{-}} \ControlFlowTok{function}\NormalTok{(x) }\FunctionTok{sum}\NormalTok{(}\FunctionTok{choose}\NormalTok{(x, }\DecValTok{2}\NormalTok{))}
\NormalTok{  f }\OtherTok{\textless{}{-}} \FunctionTok{pares}\NormalTok{(}\FunctionTok{rowSums}\NormalTok{(t))}
\NormalTok{  g }\OtherTok{\textless{}{-}} \FunctionTok{pares}\NormalTok{(}\FunctionTok{colSums}\NormalTok{(t))}
\NormalTok{  e }\OtherTok{\textless{}{-}}\NormalTok{ f }\SpecialCharTok{*}\NormalTok{ g }\SpecialCharTok{/} \FunctionTok{choose}\NormalTok{(}\FunctionTok{sum}\NormalTok{(t), }\DecValTok{2}\NormalTok{)}
\NormalTok{  (}\FunctionTok{pares}\NormalTok{(t) }\SpecialCharTok{{-}}\NormalTok{ e) }\SpecialCharTok{/}\NormalTok{ ((f }\SpecialCharTok{+}\NormalTok{ g) }\SpecialCharTok{/} \DecValTok{2} \SpecialCharTok{{-}}\NormalTok{ e)}
\NormalTok{\}}
\CommentTok{\# K{-}medias sobre las primeras CP sin re{-}escalarlas, de modo que}
\CommentTok{\# cada componente pesa según su varianza.}
\NormalTok{en\_cp }\OtherTok{\textless{}{-}} \ControlFlowTok{function}\NormalTok{(q) \{}
  \FunctionTok{set.seed}\NormalTok{(}\DecValTok{42}\NormalTok{)}
  \FunctionTok{rotular}\NormalTok{(}\FunctionTok{kmeans}\NormalTok{(acp\_sin4}\SpecialCharTok{$}\NormalTok{x[, }\DecValTok{1}\SpecialCharTok{:}\NormalTok{q], }\AttributeTok{centers =} \DecValTok{3}\NormalTok{,}
                 \AttributeTok{nstart =} \DecValTok{100}\NormalTok{)}\SpecialCharTok{$}\NormalTok{cluster)}
\NormalTok{\}}
\NormalTok{arbol }\OtherTok{\textless{}{-}} \ControlFlowTok{function}\NormalTok{(m) }\FunctionTok{rotular}\NormalTok{(}\FunctionTok{cutree}\NormalTok{(}\FunctionTok{hclust}\NormalTok{(distancias, m), }\DecValTok{3}\NormalTok{))}
\NormalTok{alternativas }\OtherTok{\textless{}{-}} \FunctionTok{list}\NormalTok{(}\StringTok{"k{-}medias"} \OtherTok{=}\NormalTok{ kmedias,}
                     \StringTok{"promedio"} \OtherTok{=}\NormalTok{ promedio,}
                     \StringTok{"completo"} \OtherTok{=} \FunctionTok{arbol}\NormalTok{(}\StringTok{"complete"}\NormalTok{),}
                     \StringTok{"Ward"} \OtherTok{=} \FunctionTok{arbol}\NormalTok{(}\StringTok{"ward.D2"}\NormalTok{),}
                     \StringTok{"k{-}medias CP1{-}CP4"} \OtherTok{=} \FunctionTok{en\_cp}\NormalTok{(}\DecValTok{4}\NormalTok{),}
                     \StringTok{"k{-}medias CP1{-}CP6"} \OtherTok{=} \FunctionTok{en\_cp}\NormalTok{(}\DecValTok{6}\NormalTok{))}
\NormalTok{knitr}\SpecialCharTok{::}\FunctionTok{kable}\NormalTok{(}\FunctionTok{t}\NormalTok{(}\FunctionTok{sapply}\NormalTok{(alternativas, }\ControlFlowTok{function}\NormalTok{(g) }\FunctionTok{c}\NormalTok{(}
  \AttributeTok{tamanos =} \FunctionTok{paste}\NormalTok{(}\FunctionTok{table}\NormalTok{(g), }\AttributeTok{collapse =} \StringTok{"/"}\NormalTok{),}
  \AttributeTok{explicada =} \FunctionTok{sprintf}\NormalTok{(}\StringTok{"\%.1f"}\NormalTok{, }\DecValTok{100} \SpecialCharTok{*} \FunctionTok{explicada}\NormalTok{(g)),}
  \AttributeTok{silueta =} \FunctionTok{sprintf}\NormalTok{(}\StringTok{"\%.3f"}\NormalTok{, }\FunctionTok{sil}\NormalTok{(g)}\SpecialCharTok{$}\NormalTok{avg.width),}
  \AttributeTok{ARI =} \FunctionTok{sprintf}\NormalTok{(}\StringTok{"\%.3f"}\NormalTok{, }\FunctionTok{ari}\NormalTok{(kmedias, g))))),}
  \AttributeTok{col.names =} \FunctionTok{c}\NormalTok{(}\StringTok{"tamaños"}\NormalTok{, }\StringTok{"explicada (\%)"}\NormalTok{, }\StringTok{"silueta"}\NormalTok{, }\StringTok{"ARI"}\NormalTok{),}
  \AttributeTok{align =} \StringTok{"lrrr"}\NormalTok{,}
  \AttributeTok{caption =} \StringTok{"Particiones en tres grupos frente a k{-}medias."}\NormalTok{)}
\end{Highlighting}
\end{Shaded}

\Needspace{2.31in}

\begin{longtable}[]{@{}llrrr@{}}
\caption{Particiones en tres grupos frente a k-medias.}\tabularnewline
\toprule\noalign{}
& tamaños & explicada (\%) & silueta & ARI \\
\midrule\noalign{}
\endfirsthead
\toprule\noalign{}
& tamaños & explicada (\%) & silueta & ARI \\
\midrule\noalign{}
\endhead
\bottomrule\noalign{}
\endlastfoot
k-medias & 97/8/2 & 46.6 & 0.666 & 1.000 \\
promedio & 103/2/2 & 39.0 & 0.713 & 0.528 \\
completo & 103/2/2 & 39.0 & 0.713 & 0.528 \\
Ward & 97/8/2 & 46.6 & 0.666 & 1.000 \\
k-medias CP1-CP4 & 97/7/3 & 46.5 & 0.650 & 0.994 \\
k-medias CP1-CP6 & 97/8/2 & 46.6 & 0.666 & 1.000 \\
\end{longtable}

\Needspace{0.95in}

\begin{Shaded}
\begin{Highlighting}[]
\FunctionTok{round}\NormalTok{(}\FunctionTok{c}\NormalTok{(}\AttributeTok{promedio\_completo =} \FunctionTok{ari}\NormalTok{(promedio, alternativas}\SpecialCharTok{$}\NormalTok{completo),}
        \AttributeTok{cp4 =} \DecValTok{100} \SpecialCharTok{*} \FunctionTok{sum}\NormalTok{(acp\_sin4}\SpecialCharTok{$}\NormalTok{sdev[}\DecValTok{1}\SpecialCharTok{:}\DecValTok{4}\NormalTok{]}\SpecialCharTok{\^{}}\DecValTok{2}\NormalTok{) }\SpecialCharTok{/} \DecValTok{57}\NormalTok{,}
        \AttributeTok{cp6 =} \DecValTok{100} \SpecialCharTok{*} \FunctionTok{sum}\NormalTok{(acp\_sin4}\SpecialCharTok{$}\NormalTok{sdev[}\DecValTok{1}\SpecialCharTok{:}\DecValTok{6}\NormalTok{]}\SpecialCharTok{\^{}}\DecValTok{2}\NormalTok{) }\SpecialCharTok{/} \DecValTok{57}\NormalTok{), }\DecValTok{2}\NormalTok{)}
\end{Highlighting}
\end{Shaded}

\begin{verbatim}
promedio_completo               cp4               cp6 
             1.00             73.92             82.78 
\end{verbatim}

\Needspace{4.37in}

\begin{Shaded}
\begin{Highlighting}[]
\NormalTok{sp }\OtherTok{\textless{}{-}} \FunctionTok{silhouette}\NormalTok{(}\FunctionTok{as.integer}\NormalTok{(promedio), distancias)}
\NormalTok{comunes }\OtherTok{\textless{}{-}}\NormalTok{ kmedias }\SpecialCharTok{==} \StringTok{"A"} \SpecialCharTok{\&}\NormalTok{ promedio }\SpecialCharTok{==} \StringTok{"A"}
\NormalTok{seis }\OtherTok{\textless{}{-}}\NormalTok{ kmedias }\SpecialCharTok{==} \StringTok{"B"} \SpecialCharTok{\&}\NormalTok{ promedio }\SpecialCharTok{==} \StringTok{"A"}
\NormalTok{media }\OtherTok{\textless{}{-}} \ControlFlowTok{function}\NormalTok{(x, i) }\FunctionTok{round}\NormalTok{(}\FunctionTok{mean}\NormalTok{(x[i, }\StringTok{"sil\_width"}\NormalTok{]), }\DecValTok{3}\NormalTok{)}
\NormalTok{knitr}\SpecialCharTok{::}\FunctionTok{kable}\NormalTok{(}\FunctionTok{data.frame}\NormalTok{(}
  \AttributeTok{paises =} \FunctionTok{c}\NormalTok{(}\FunctionTok{sum}\NormalTok{(comunes), }\FunctionTok{sum}\NormalTok{(seis), }\FunctionTok{sum}\NormalTok{(promedio }\SpecialCharTok{!=} \StringTok{"A"}\NormalTok{)),}
  \AttributeTok{k\_medias =} \FunctionTok{c}\NormalTok{(}\FunctionTok{media}\NormalTok{(s, comunes), }\FunctionTok{media}\NormalTok{(s, seis),}
               \FunctionTok{media}\NormalTok{(s, promedio }\SpecialCharTok{!=} \StringTok{"A"}\NormalTok{)),}
  \AttributeTok{promedio =} \FunctionTok{c}\NormalTok{(}\FunctionTok{media}\NormalTok{(sp, comunes), }\FunctionTok{media}\NormalTok{(sp, seis),}
               \FunctionTok{media}\NormalTok{(sp, promedio }\SpecialCharTok{!=} \StringTok{"A"}\NormalTok{)),}
  \AttributeTok{row.names =} \FunctionTok{c}\NormalTok{(}\StringTok{"masa común"}\NormalTok{, }\StringTok{"seis intermedios"}\NormalTok{,}
                \StringTok{"cuatro extremos"}\NormalTok{)),}
  \AttributeTok{col.names =} \FunctionTok{c}\NormalTok{(}\StringTok{"países"}\NormalTok{, }\StringTok{"k{-}medias"}\NormalTok{, }\StringTok{"promedio"}\NormalTok{),}
  \AttributeTok{caption =} \StringTok{"Silueta media de cada bloque de países."}\NormalTok{)}
\end{Highlighting}
\end{Shaded}

\Needspace{1.74in}

\begin{longtable}[]{@{}lrrr@{}}
\caption{Silueta media de cada bloque de países.}\tabularnewline
\toprule\noalign{}
& países & k-medias & promedio \\
\midrule\noalign{}
\endfirsthead
\toprule\noalign{}
& países & k-medias & promedio \\
\midrule\noalign{}
\endhead
\bottomrule\noalign{}
\endlastfoot
masa común & 97 & 0.731 & 0.761 \\
seis intermedios & 6 & -0.051 & 0.262 \\
cuatro extremos & 4 & 0.159 & 0.220 \\
\end{longtable}

\begin{notas}

\textbf{Interpretaciones.}

\begin{enumerate}
\def\labelenumi{\alph{enumi}.}
\tightlist
\item
  \textbf{Cada método gana en su propio criterio.} K-medias explica más
  varianza (46.6 \% frente a 39.0 \%) porque eso es lo que minimiza y el
  enlace promedio tiene mejor silueta (0.713 frente a 0.666) porque no
  fuerza un grupo intermedio, de modo que «mejor» depende de qué se le
  pida a la partición.
\item
  \textbf{El grupo B de k-medias.} Su silueta promedio es \(-0.022\),
  con cuatro de sus ocho países en negativo (Alemania \(-0.222\), Reino
  Unido \(-0.092\), España \(-0.054\) y Corea del Sur \(-0.021\)) y los
  otros cuatro apenas por encima de cero (el mejor es Brasil con 0.082),
  por lo que B se comporta más como una franja de transición que como un
  grupo.
\item
  \textbf{Por qué el enlace promedio aísla parejas.} India y Japón, y
  Brasil y Canadá, están lejos en promedio de todos los demás, de modo
  que el árbol los une entre sí muy arriba y los deja para el final. Al
  cortar en cuatro o cinco grupos se siguen separando parejas o países
  sueltos (101/2/2/2 y 101/2/2/1/1) y nunca las seis economías
  intermedias.
\item
  \textbf{El arranque importa.} Con \texttt{nstart\ =\ 25} la inercia
  depende de la semilla (3274.5 o 3231.9) y solo con 100 arranques se
  llega siempre a 3226.7 así que una partición de k-medias sin varios
  arranques puede ser un óptimo local.
\item
  \textbf{Qué k defendería cada uno.} Por silueta los dos preferirían
  \(k = 2\) pues k-medias llega a 0.698 (98/9) y el enlace promedio a
  0.752 (105/2, solo India y Japón aparte). El \(k = 3\) se eligió por
  la varianza explicada y por la lectura del plano, no por la silueta.
\end{enumerate}

\textbf{Decisiones.}

\begin{itemize}
\item
  \textbf{Las ocho economías del grupo B son el borde de la masa central
  y no un grupo propio.} Su silueta promedio es de \(-0.022\) y la más
  alta es la de Brasil (0.082), de modo que las ocho quedan por debajo
  de 0.25, que es el valor desde el cual no se reconoce una estructura
  sustancial, mientras que siete de ellas tienen como grupo vecino a la
  masa central A y solo Brasil tiene como vecino al grupo de India y
  Japón. Si B se devuelve a la masa central la silueta sube de 0.666 a
  0.752 pero la varianza explicada cae de 46.6 \% a 24.5 \% así que B
  se mantiene en la partición como una franja de transición que recoge
  la inercia de las economías grandes y no como un perfil energético
  distinto.

  \emph{Respaldo.} \textcite[cap.~2]{kaufman1990} leen una silueta media
  de 0.25 o menos como ausencia de estructura sustancial y en la sesión
  5 se mostró que un punto que queda entre dos grupos es el caso en que
  un método suave le asignaría una probabilidad cercana a 0.5 de
  pertenecer a cada uno.
\item
  \textbf{Se defiende la partición de k-medias y el enlace promedio
  queda como chequeo.} K-medias explica el 46.6 \% de la inercia frente
  al 39.0 \% y es la partición que el enunciado usa en la Pregunta 11,
  mientras que la ventaja del enlace promedio en silueta (0.713 frente a
  0.666) no viene de sus dos parejas, cuyos cuatro países promedian
  0.220, sino de que al devolver a los seis países intermedios a la masa
  central estos pasan de \(-0.051\) a 0.262 y los 97 países comunes de
  0.731 a 0.761 porque su grupo vecino queda más lejos. Es decir que la
  silueta premia no formar el grupo intermedio, que es justo lo que se
  discute, por lo que no sirve para decidir entre los dos métodos. El
  índice de Rand ajustado (ARI), que vale 1 si dos particiones coinciden
  y cerca de 0 si coinciden como al azar, es de 0.528 entre ambos y esto
  se debe a que la diferencia se concentra en esos seis países pues los
  otros 101 quedan en el grupo equivalente.

  \emph{Respaldo.} En la sesión 5 el objetivo del agrupamiento se
  planteó como grupos homogéneos por dentro y heterogéneos entre sí, que
  es el criterio que k-medias minimiza, se recomendó el argumento
  \texttt{nstart} para no quedarse en un óptimo local y los métodos se
  compararon por la concordancia de sus asignaciones, que aquí se mide
  con el índice de \textcite{hubert1985}. El enlace promedio es el que
  se usa en la mayoría de los casos por ser robusto ante distancias
  particulares, por lo que entra como chequeo de la masa central y de
  los extremos.
\item
  \textbf{Se agrupa sobre las 57 variables estandarizadas y las
  componentes quedan como diagnóstico.} K-medias sobre las cuatro
  primeras componentes sin re-escalar, que en los 107 países reúnen el
  73.92 \% de la varianza, da grupos de 97, 7 y 3 países con 46.5 \%
  explicada, silueta de 0.650 y ARI de 0.994 frente a la partición de
  las 57 variables, es decir que la única diferencia es Brasil, que pasa
  del grupo B al de India y Japón, mientras que con seis componentes
  (82.78 \%, las que pide el 80 \% en la Pregunta 7) la partición es
  idéntica (ARI de 1.000). Esto tiene sentido pues las coordenadas sin
  re-escalar son una rotación de las variables estandarizadas que con
  todas las componentes conserva las mismas distancias, de modo que
  agrupar sobre las componentes no quita el peso repetido de las
  variables casi iguales (que ya está dentro de las componentes) y lo
  único que agrega quedarse con unas pocas es el riesgo de descartar
  información.

  \emph{Respaldo.} El profesor dijo sobre agrupar en componentes que «en
  general no lo recomiendo» (sesión 4), usó el plano factorial como «una
  herramienta de diagnóstico, no como una parte del proceso» y advirtió
  que las componentes tardías pueden caracterizar un grupo que se pierde
  al quedarse con las primeras (sesión 5), que es lo que mostró
  \textcite{chang1983} y lo que se critica del llamado análisis en
  tándem \parencite{arabie1994}.
\item
  \textbf{El grupo intermedio es un rasgo del criterio de inercia y no
  solo de k-medias.} El enlace de Ward reproduce exactamente la
  partición de k-medias (97/8/2 y ARI de 1.000) mientras que el enlace
  completo reproduce la del promedio (103/2/2 y ARI de 1.00 entre
  ambos), de modo que los métodos se separan por lo que optimizan y no
  por ser jerárquicos o de partición pues Ward une en cada paso los dos
  grupos que menos aumentan la inercia intra-grupo, que es la misma
  \(W(k)\) que minimiza k-medias, mientras que el promedio y el completo
  miden distancias entre pares de países y por eso dejan para el final a
  India, Japón, Brasil y Canadá.

  \emph{Respaldo.} \textcite{ward1963} define su enlace como el que
  minimiza el aumento de la suma de cuadrados dentro de los grupos y
  \texttt{ward.D2} es la versión de \texttt{hclust()} que aplica ese
  criterio sobre distancias euclidianas \parencite{murtagh2014}. En la
  sesión 5 se recomendó el enlace promedio en la mayoría de los casos y
  el completo o el simple para caracterizar valores inusuales, que es lo
  que hacen aquí al aislar las dos parejas extremas.
\end{itemize}

\end{notas}

\Needspace{4in}

\section{Pregunta 11 --- Distribución de las cinco variables por
grupo}\label{p11}

\Needspace{8.33in}

\begin{Shaded}
\begin{Highlighting}[]
\NormalTok{largo }\OtherTok{\textless{}{-}} \FunctionTok{stack}\NormalTok{(num[cinco])}
\NormalTok{largo}\SpecialCharTok{$}\NormalTok{grupo }\OtherTok{\textless{}{-}} \FunctionTok{rep}\NormalTok{(kmedias, }\AttributeTok{times =} \FunctionTok{length}\NormalTok{(cinco))}
\NormalTok{etiqueta }\OtherTok{\textless{}{-}} \ControlFlowTok{function}\NormalTok{(g) }\FunctionTok{sprintf}\NormalTok{(}\StringTok{"\%s (\%d)"}\NormalTok{, g, }\FunctionTok{table}\NormalTok{(kmedias)[g])}
\FunctionTok{ggplot}\NormalTok{(largo, }\FunctionTok{aes}\NormalTok{(}\AttributeTok{x =} \FunctionTok{etiqueta}\NormalTok{(grupo), }\AttributeTok{y =}\NormalTok{ values, }\AttributeTok{fill =}\NormalTok{ grupo)) }\SpecialCharTok{+}
  \FunctionTok{geom\_boxplot}\NormalTok{(}\AttributeTok{alpha =} \FloatTok{0.6}\NormalTok{, }\AttributeTok{show.legend =} \ConstantTok{FALSE}\NormalTok{,}
               \AttributeTok{outlier.size =} \FloatTok{0.8}\NormalTok{) }\SpecialCharTok{+}
  \FunctionTok{stat\_summary}\NormalTok{(}\AttributeTok{fun =}\NormalTok{ median, }\AttributeTok{geom =} \StringTok{"text"}\NormalTok{, }\AttributeTok{size =} \FloatTok{2.4}\NormalTok{,}
               \AttributeTok{vjust =} \SpecialCharTok{{-}}\FloatTok{0.5}\NormalTok{,}
               \FunctionTok{aes}\NormalTok{(}\AttributeTok{label =} \FunctionTok{round}\NormalTok{(}\FunctionTok{after\_stat}\NormalTok{(y), }\DecValTok{1}\NormalTok{))) }\SpecialCharTok{+}
  \FunctionTok{facet\_wrap}\NormalTok{(}\SpecialCharTok{\textasciitilde{}}\NormalTok{ ind, }\AttributeTok{scales =} \StringTok{"free\_y"}\NormalTok{, }\AttributeTok{nrow =} \DecValTok{2}\NormalTok{) }\SpecialCharTok{+}
  \CommentTok{\# La escala se transforma DESPUÉS de calcular cajas y medianas,}
  \CommentTok{\# para que la mediana anotada sea la de los datos originales.}
  \FunctionTok{coord\_transform}\NormalTok{(}\AttributeTok{y =} \StringTok{"pseudo\_log"}\NormalTok{) }\SpecialCharTok{+}
  \FunctionTok{scale\_y\_continuous}\NormalTok{(}\AttributeTok{breaks =} \FunctionTok{c}\NormalTok{(}\DecValTok{0}\NormalTok{, }\DecValTok{10}\NormalTok{, }\DecValTok{100}\NormalTok{, }\DecValTok{1000}\NormalTok{, }\DecValTok{10000}\NormalTok{)) }\SpecialCharTok{+}
  \FunctionTok{scale\_fill\_manual}\NormalTok{(}\AttributeTok{values =} \FunctionTok{c}\NormalTok{(}\StringTok{"\#3572a8"}\NormalTok{, }\StringTok{"\#f08a3c"}\NormalTok{, }\StringTok{"\#a6192e"}\NormalTok{)) }\SpecialCharTok{+}
  \FunctionTok{labs}\NormalTok{(}\AttributeTok{x =} \StringTok{"Grupo de k{-}medias (países)"}\NormalTok{, }\AttributeTok{y =} \StringTok{"Valor"}\NormalTok{) }\SpecialCharTok{+}
  \FunctionTok{theme\_minimal}\NormalTok{(}\AttributeTok{base\_size =} \DecValTok{9}\NormalTok{)}
\end{Highlighting}
\end{Shaded}

\begin{figure}[H]

{\centering \pandocbounded{\includegraphics[keepaspectratio]{sda-lab-exploracion_files/figure-latex/p11-1.pdf}}

}

\caption{Distribución de las cinco variables de la Pregunta 1 por grupo
de k-medias, en escala pseudo-logarítmica y con la mediana anotada: El
consumo y la generación eléctrica y renovable crecen de A a C, mientras
que la producción y la exportación de crudo no separan los grupos.}

\end{figure}%

\Needspace{2.47in}

\begin{Shaded}
\begin{Highlighting}[]
\NormalTok{knitr}\SpecialCharTok{::}\FunctionTok{kable}\NormalTok{(}\FunctionTok{sapply}\NormalTok{(cinco, }\ControlFlowTok{function}\NormalTok{(v) }\FunctionTok{tapply}\NormalTok{(num[[v]], kmedias,}
  \ControlFlowTok{function}\NormalTok{(x) }\FunctionTok{sprintf}\NormalTok{(}\StringTok{"\%g a \%g"}\NormalTok{, }\FunctionTok{round}\NormalTok{(}\FunctionTok{min}\NormalTok{(x), }\DecValTok{1}\NormalTok{),}
                      \FunctionTok{round}\NormalTok{(}\FunctionTok{max}\NormalTok{(x), }\DecValTok{1}\NormalTok{)))),}
  \AttributeTok{caption =} \StringTok{"Rango (mínimo a máximo) de cada variable por grupo."}\NormalTok{)}
\end{Highlighting}
\end{Shaded}

\Needspace{1.74in}

\begin{longtable}[]{@{}llllll@{}}
\caption{Rango (mínimo a máximo) de cada variable por
grupo.}\tabularnewline
\toprule\noalign{}
& Crudo\_prod & PetLiq\_cons & Crudo\_exp & Elec\_gen & Renov\_gen \\
\midrule\noalign{}
\endfirsthead
\toprule\noalign{}
& Crudo\_prod & PetLiq\_cons & Crudo\_exp & Elec\_gen & Renov\_gen \\
\midrule\noalign{}
\endhead
\bottomrule\noalign{}
\endlastfoot
A & 0 a 3370 & 15 a 2040 & 0 a 2651 & 0.9 a 290 & 0 a 138 \\
B & 0 a 3613 & 791.3 a 3279 & 0 a 2382 & 198 a 636 & 10 a 429 \\
C & 4.8 a 767 & 3550 a 4288 & 0 a 0.7 & 1027 a 1258 & 129 a 198 \\
\end{longtable}

\Needspace{2.66in}

\begin{Shaded}
\begin{Highlighting}[]
\NormalTok{exportadores }\OtherTok{\textless{}{-}} \FunctionTok{order}\NormalTok{(}\SpecialCharTok{{-}}\NormalTok{num}\SpecialCharTok{$}\NormalTok{Crudo\_exp)[}\DecValTok{1}\SpecialCharTok{:}\DecValTok{5}\NormalTok{]}
\FunctionTok{data.frame}\NormalTok{(}\AttributeTok{exportacion =}\NormalTok{ num}\SpecialCharTok{$}\NormalTok{Crudo\_exp[exportadores],}
           \AttributeTok{grupo =}\NormalTok{ kmedias[exportadores],}
           \AttributeTok{row.names =} \FunctionTok{rownames}\NormalTok{(num)[exportadores])}
\end{Highlighting}
\end{Shaded}

\Needspace{1.93in}

{\def\LTcaptype{none} % do not increment counter
\begin{longtable}[]{@{}lrl@{}}
\toprule\noalign{}
& exportacion & grupo \\
\midrule\noalign{}
\endhead
\bottomrule\noalign{}
\endlastfoot
UnitedArabEmirates & 2651 & A \\
Iraq & 2518 & A \\
Canada & 2382 & B \\
Nigeria & 2120 & A \\
Kuwait & 1995 & A \\
\end{longtable}
}

\Needspace{0.57in}

\begin{Shaded}
\begin{Highlighting}[]
\FunctionTok{c}\NormalTok{(}\AttributeTok{sin\_exportar\_A =} \FunctionTok{sum}\NormalTok{(num}\SpecialCharTok{$}\NormalTok{Crudo\_exp[kmedias }\SpecialCharTok{==} \StringTok{"A"}\NormalTok{] }\SpecialCharTok{==} \DecValTok{0}\NormalTok{))}
\end{Highlighting}
\end{Shaded}

\begin{verbatim}
sin_exportar_A 
            38 
\end{verbatim}

\textbf{¿Hay diferencias importantes entre grupos en una misma
variable?} Sí en tres de las cinco variables (consumo de petróleo,
generación eléctrica y generación renovable) y no en las dos petroleras
(producción y exportación de crudo). Antes de leerlas vale la pena
anotar que el grupo C tiene solo dos países (India y Japón) y el B ocho,
de modo que sus cajas son apenas un resumen de pocos puntos y no una
distribución y que el eje es pseudo-logarítmico porque los valores de
los tres grupos difieren en órdenes de magnitud.

\begin{enumerate}
\def\labelenumi{\alph{enumi}.}
\tightlist
\item
  \textbf{Consumo de petróleo y generación eléctrica.} Las medianas
  crecen con claridad de A a C pues el consumo pasa de 155 a 1632 y a
  3919 TBPD y la generación de 27 a 422.5 y a 1142.5 BKWH. India y Japón
  superan a todos los países de los otros grupos en las dos variables,
  aunque A y B se solapan en sus extremos, por ejemplo México consume
  2040 TBPD y genera 290 BKWH siendo del grupo A, mientras que Alemania
  consume 791.3 TBPD y genera 198 BKWH siendo del grupo B.
\item
  \textbf{Generación renovable.} También crece de A a C (medianas de 6,
  102.5 y 163.5 BKWH) pero con mucho solapamiento ya que el grupo B va
  de 10 BKWH (Corea del Sur) a 429 BKWH (Brasil) y el máximo de A
  (Noruega, 138 BKWH) supera incluso a Japón (129 BKWH) del grupo C.
\item
  \textbf{Producción de crudo.} No distingue los grupos, dado que los
  tres tienen países que casi no producen y el rango de A (de 0 a 3370
  TBPD, este último Irak) cubre casi por completo los de B (de 0 a 3613)
  y C (de 4.8 a 767).
\item
  \textbf{Exportación de crudo.} Tampoco los distingue y además no crece
  con el tamaño pues las medianas son 3.3, 6.2 y 0.35 TBPD, India y
  Japón casi no exportan (0.7 y 0 TBPD) y cuatro de los cinco mayores
  exportadores de los 107 países (Emiratos Árabes Unidos con 2651 TBPD,
  Irak con 2518, Nigeria con 2120 y Kuwait con 1995) están en el grupo
  A, junto a 38 países que no exportan nada, mientras que el quinto,
  Canadá (2382 TBPD), está en B por su tamaño y no por su petróleo.
\end{enumerate}

Se concluye que k-medias agrupó a los países por el tamaño de su consumo
y de su sistema eléctrico, que es la información del primer eje del ACP
y no por ser o no productores y exportadores de petróleo, que es la
información del segundo eje.

\begin{enumerate}
\def\labelenumi{\alph{enumi}.}
\tightlist
\item
  \textbf{Lectura local.} Emiratos Árabes Unidos, el mayor exportador de
  crudo sin los cuatro atípicos, queda en el mismo grupo que países que
  no exportan nada porque en todo lo demás (consumo, generación) es un
  país mediano.
\item
  \textbf{Lectura global.} La partición sigue al tamaño porque la base
  mide sobre todo tamaño ya que 27 de las 57 variables están ligadas al
  primer eje y la dimensión petrolera la cargan unas pocas, de modo que
  en la distancia euclidiana entre países el petróleo pesa poco. Haber
  elegido la exportación en la Pregunta 1 es lo que permite verlo pues
  con las emisiones de CO\textsubscript{2} las cinco cajas habrían
  repetido el consumo.
\end{enumerate}

Para ver si otra partición separa a los exportadores se repite k-medias
con tres grupos sobre dos conjuntos que dan más peso al petróleo, las
cinco variables de la Pregunta 1 estandarizadas y las componentes 2 y 3
del ACP de la Pregunta 7.

\Needspace{4.37in}

\begin{Shaded}
\begin{Highlighting}[]
\CommentTok{\# Dos particiones que dan más peso al petróleo: Las cinco}
\CommentTok{\# variables de la Pregunta 1 y las componentes 2 y 3 del ACP.}
\FunctionTok{set.seed}\NormalTok{(}\DecValTok{42}\NormalTok{)}
\NormalTok{en\_cinco }\OtherTok{\textless{}{-}} \FunctionTok{rotular}\NormalTok{(}\FunctionTok{kmeans}\NormalTok{(}\FunctionTok{scale}\NormalTok{(num[cinco]), }\AttributeTok{centers =} \DecValTok{3}\NormalTok{,}
                           \AttributeTok{nstart =} \DecValTok{100}\NormalTok{)}\SpecialCharTok{$}\NormalTok{cluster)}
\FunctionTok{set.seed}\NormalTok{(}\DecValTok{42}\NormalTok{)}
\NormalTok{en\_cp23 }\OtherTok{\textless{}{-}} \FunctionTok{rotular}\NormalTok{(}\FunctionTok{kmeans}\NormalTok{(acp\_sin4}\SpecialCharTok{$}\NormalTok{x[, }\DecValTok{2}\SpecialCharTok{:}\DecValTok{3}\NormalTok{], }\AttributeTok{centers =} \DecValTok{3}\NormalTok{,}
                          \AttributeTok{nstart =} \DecValTok{100}\NormalTok{)}\SpecialCharTok{$}\NormalTok{cluster)}
\FunctionTok{data.frame}\NormalTok{(}\AttributeTok{exportacion =}\NormalTok{ num}\SpecialCharTok{$}\NormalTok{Crudo\_exp[exportadores],}
           \AttributeTok{k\_medias =}\NormalTok{ kmedias[exportadores],}
           \AttributeTok{cinco =}\NormalTok{ en\_cinco[exportadores],}
           \AttributeTok{cp2\_cp3 =}\NormalTok{ en\_cp23[exportadores],}
           \AttributeTok{row.names =} \FunctionTok{rownames}\NormalTok{(num)[exportadores])}
\end{Highlighting}
\end{Shaded}

\Needspace{1.93in}

{\def\LTcaptype{none} % do not increment counter
\begin{longtable}[]{@{}lrlll@{}}
\toprule\noalign{}
& exportacion & k\_medias & cinco & cp2\_cp3 \\
\midrule\noalign{}
\endhead
\bottomrule\noalign{}
\endlastfoot
UnitedArabEmirates & 2651 & A & B & B \\
Iraq & 2518 & A & B & B \\
Canada & 2382 & B & C & B \\
Nigeria & 2120 & A & B & B \\
Kuwait & 1995 & A & B & B \\
\end{longtable}
}

\Needspace{2.66in}

\begin{Shaded}
\begin{Highlighting}[]
\NormalTok{knitr}\SpecialCharTok{::}\FunctionTok{kable}\NormalTok{(}\FunctionTok{t}\NormalTok{(}\FunctionTok{sapply}\NormalTok{(}\FunctionTok{list}\NormalTok{(}\AttributeTok{cinco =}\NormalTok{ en\_cinco, }\AttributeTok{cp2\_cp3 =}\NormalTok{ en\_cp23),}
  \ControlFlowTok{function}\NormalTok{(g) }\FunctionTok{c}\NormalTok{(}\AttributeTok{tamanos =} \FunctionTok{paste}\NormalTok{(}\FunctionTok{table}\NormalTok{(g), }\AttributeTok{collapse =} \StringTok{"/"}\NormalTok{),}
    \AttributeTok{explicada =} \FunctionTok{sprintf}\NormalTok{(}\StringTok{"\%.1f"}\NormalTok{, }\DecValTok{100} \SpecialCharTok{*} \FunctionTok{explicada}\NormalTok{(g)),}
    \AttributeTok{ARI =} \FunctionTok{sprintf}\NormalTok{(}\StringTok{"\%.3f"}\NormalTok{, }\FunctionTok{ari}\NormalTok{(kmedias, g))))),}
  \AttributeTok{col.names =} \FunctionTok{c}\NormalTok{(}\StringTok{"tamaños"}\NormalTok{, }\StringTok{"explicada (\%)"}\NormalTok{, }\StringTok{"ARI"}\NormalTok{), }\AttributeTok{align =} \StringTok{"lrr"}\NormalTok{,}
  \AttributeTok{caption =} \StringTok{"Particiones con más peso petrolero."}\NormalTok{)}
\end{Highlighting}
\end{Shaded}

\Needspace{1.55in}

\begin{longtable}[]{@{}llrr@{}}
\caption{Particiones con más peso petrolero.}\tabularnewline
\toprule\noalign{}
& tamaños & explicada (\%) & ARI \\
\midrule\noalign{}
\endfirsthead
\toprule\noalign{}
& tamaños & explicada (\%) & ARI \\
\midrule\noalign{}
\endhead
\bottomrule\noalign{}
\endlastfoot
cinco & 93/10/4 & 36.9 & 0.244 \\
cp2\_cp3 & 93/10/4 & 21.1 & 0.405 \\
\end{longtable}

\Needspace{1.71in}

\begin{Shaded}
\begin{Highlighting}[]
\FunctionTok{split}\NormalTok{(}\FunctionTok{rownames}\NormalTok{(num), en\_cinco)[}\FunctionTok{c}\NormalTok{(}\StringTok{"B"}\NormalTok{, }\StringTok{"C"}\NormalTok{)]}
\end{Highlighting}
\end{Shaded}

\begin{verbatim}
$B
 [1] "Angola"             "Iran"              
 [3] "Iraq"               "Kazakhstan"        
 [5] "Kuwait"             "Mexico"            
 [7] "Nigeria"            "Norway"            
 [9] "UnitedArabEmirates" "Venezuela"         

$C
[1] "Brazil" "Canada" "India"  "Japan" 
\end{verbatim}

\Needspace{1.71in}

\begin{Shaded}
\begin{Highlighting}[]
\FunctionTok{split}\NormalTok{(}\FunctionTok{rownames}\NormalTok{(num), en\_cp23)[}\FunctionTok{c}\NormalTok{(}\StringTok{"B"}\NormalTok{, }\StringTok{"C"}\NormalTok{)]}
\end{Highlighting}
\end{Shaded}

\begin{verbatim}
$B
 [1] "Brazil"             "Canada"            
 [3] "Iran"               "Iraq"              
 [5] "Kuwait"             "Mexico"            
 [7] "Nigeria"            "Norway"            
 [9] "UnitedArabEmirates" "Venezuela"         

$C
[1] "France"  "Germany" "Italy"   "Spain"  
\end{verbatim}

\begin{notas}

\textbf{Interpretaciones.}

\begin{enumerate}
\def\labelenumi{\alph{enumi}.}
\tightlist
\item
  \textbf{Lo que separa y lo que no.} Las medianas de consumo (155, 1632
  y 3919 TBPD) y de generación (27, 422.5 y 1142.5 BKWH) crecen de A a
  C, mientras que las de exportación (3.3, 6.2 y 0.35 TBPD) no siguen
  ningún orden.
\item
  \textbf{El petróleo queda dentro de A.} Cuatro de los cinco mayores
  exportadores de los 107 países (Emiratos Árabes Unidos, Irak, Nigeria
  y Kuwait) están en A, el mismo grupo de la masa central, es decir que
  la partición no reconoce al brazo de productores que se veía en la
  Pregunta 7.
\item
  \textbf{Cajas de pocos países.} C tiene dos países y B ocho así que
  sus cajas describen casos y no distribuciones.
\end{enumerate}

\textbf{Decisiones.}

\begin{itemize}
\item
  \textbf{Agrupar con las cinco variables sí separa a los exportadores.}
  K-medias sobre las cinco variables estandarizadas forma un grupo B de
  diez países (Angola, Irán, Irak, Kazajistán, Kuwait, México, Nigeria,
  Noruega, Emiratos Árabes Unidos y Venezuela) que reúne a cuatro de los
  cinco mayores exportadores y deja a Canadá con Brasil, India y Japón,
  mientras que con las componentes 2 y 3 el grupo B reúne a los cinco
  mayores exportadores junto a Brasil, Irán, México, Noruega y Venezuela
  y el grupo C pasa a ser Francia, Alemania, Italia y España. Las dos
  particiones tienen grupos de 93, 10 y 4 países y concuerdan poco con
  la de las 57 variables (ARI de 0.244 y 0.405) así que no son un
  ajuste de la misma partición sino otra lectura de los países. Entre
  ellas se prefiere la de las cinco variables por coherencia con la
  decisión de la Pregunta 10 de no agrupar sobre componentes.

  \emph{Respaldo.} El enunciado busca caracterizar a los países «en
  términos de importación, exportación, consumo, producción» y en la
  sesión 5 el profesor advirtió que el agrupamiento queda limitado a las
  variables que se le dan y que la segunda componente guardaba una
  separación que se pierde al agrupar sobre la primera, que es lo que
  ocurre aquí con el eje petrolero.
\item
  \textbf{La partición pedida se mantiene y su sesgo de tamaño se
  reporta como limitación.} La Pregunta 11 pide los grupos de k-medias
  de la Pregunta 10, que explican el 46.6 \% de la inercia de las 57
  variables frente al 36.9 \% de la partición de las cinco variables, de
  modo que la primera describe mejor la base completa pero lo que
  describe es tamaño, mientras que la segunda es la que responde al
  objetivo de separar productores y exportadores. Por esto los boxplots
  se dejan con la partición de k-medias y la de las cinco variables se
  cita como la alternativa que podría usarse si el interés fuera el
  perfil petrolero, con la salvedad de que depende de las cinco
  variables elegidas en la Pregunta 1.

  \emph{Respaldo.} En la sesión 5 se insistió en que un grupo bien
  formado todavía no está caracterizado y que ese paso consiste en
  «entrar a los grupos, mirar cómo son», que es lo que hace esta
  pregunta y lo que deja ver que la partición de k-medias sigue al
  primer eje del ACP.
\end{itemize}

\end{notas}

\imprimirbibliografia
\appendix

\chapter{Documentación de R}\label{documentacion-r}

Se listan las funciones utilizadas en el taller junto con el paquete que
las provee y la pregunta donde se usan, cada una enlazada a su página de
ayuda oficial. El documento se generó con R 4.6.1, \texttt{ggplot2}
4.0.3, \texttt{GGally} 2.4.0, \texttt{factoextra} 2.2.0 y
\texttt{cluster} 2.1.8.2.

\Needspace{4.56in}

\Needspace{5.16in}

\begin{longtable}[]{@{}
  >{\raggedright\arraybackslash}p{(\linewidth - 4\tabcolsep) * \real{0.4444}}
  >{\raggedright\arraybackslash}p{(\linewidth - 4\tabcolsep) * \real{0.2444}}
  >{\raggedright\arraybackslash}p{(\linewidth - 4\tabcolsep) * \real{0.3111}}@{}}
\caption{Funciones de R utilizadas y enlace a su
documentación.}\tabularnewline
\toprule\noalign{}
\begin{minipage}[b]{\linewidth}\raggedright
Función
\end{minipage} & \begin{minipage}[b]{\linewidth}\raggedright
Paquete
\end{minipage} & \begin{minipage}[b]{\linewidth}\raggedright
Uso
\end{minipage} \\
\midrule\noalign{}
\endfirsthead
\toprule\noalign{}
\begin{minipage}[b]{\linewidth}\raggedright
Función
\end{minipage} & \begin{minipage}[b]{\linewidth}\raggedright
Paquete
\end{minipage} & \begin{minipage}[b]{\linewidth}\raggedright
Uso
\end{minipage} \\
\midrule\noalign{}
\endhead
\bottomrule\noalign{}
\endlastfoot
\href{https://stat.ethz.ch/R-manual/R-devel/library/utils/html/read.table.html}{\texttt{read.csv()}}
& utils & Carga \\
\href{https://stat.ethz.ch/R-manual/R-devel/library/utils/html/download.file.html}{\texttt{download.file()}}
& utils & Carga \\
\href{https://stat.ethz.ch/R-manual/R-devel/library/base/html/source.html}{\texttt{source()}}
& base & Datos \\
\href{https://stat.ethz.ch/R-manual/R-devel/library/base/html/grep.html}{\texttt{gsub()}}
& base & Datos \\
\href{https://stat.ethz.ch/R-manual/R-devel/library/stats/html/cor.html}{\texttt{cor()}}
& stats & Datos, P1, P2 \\
\href{https://stat.ethz.ch/R-manual/R-devel/library/base/html/colSums.html}{\texttt{colMeans()}}
& base & P1 \\
\href{https://stat.ethz.ch/R-manual/R-devel/library/stats/html/median.html}{\texttt{median()}}
& stats & P1, P2 \\
\href{https://stat.ethz.ch/R-manual/R-devel/library/stats/html/cor.html}{\texttt{cov()}}
& stats & P1 \\
\href{https://stat.ethz.ch/R-manual/R-devel/library/stats/html/prcomp.html}{\texttt{prcomp()}}
& stats & P3, P4, P7 \\
\href{https://stat.ethz.ch/R-manual/R-devel/library/stats/html/mahalanobis.html}{\texttt{mahalanobis()}}
& stats & P7 \\
\href{https://stat.ethz.ch/R-manual/R-devel/library/stats/html/Chisquare.html}{\texttt{qchisq()}}
& stats & P7 \\
\href{https://stat.ethz.ch/R-manual/R-devel/library/base/html/Log.html}{\texttt{log1p()}}
& base & P7 \\
\href{https://stat.ethz.ch/R-manual/R-devel/library/base/html/scale.html}{\texttt{scale()}}
& base & P9, P10 \\
\href{https://rdrr.io/cran/factoextra/man/facto_summarize.html}{\texttt{facto\_summarize()}}
& factoextra & P3 \\
\href{https://rdrr.io/cran/factoextra/man/fviz_contrib.html}{\texttt{fviz\_contrib()}}
& factoextra & P3 \\
\href{https://rdrr.io/cran/factoextra/man/eigenvalue.html}{\texttt{fviz\_eig()}}
& factoextra & P4 \\
\href{https://rdrr.io/cran/factoextra/man/get_pca.html}{\texttt{get\_pca\_var()}}
& factoextra & P5, P7 \\
\href{https://rdrr.io/cran/factoextra/man/get_pca.html}{\texttt{get\_pca\_ind()}}
& factoextra & P6, P7 \\
\href{https://rdrr.io/cran/factoextra/man/fviz_pca.html}{\texttt{fviz\_pca\_var()}}
& factoextra & P5 \\
\href{https://rdrr.io/cran/factoextra/man/fviz_pca.html}{\texttt{fviz\_pca\_ind()}}
& factoextra & P6, P7, P10 \\
\href{https://rdrr.io/cran/ggrepel/man/geom_text_repel.html}{\texttt{geom\_text\_repel()}}
& ggrepel & P6, P7, P10 \\
\href{https://rdrr.io/cran/GGally/man/ggpairs.html}{\texttt{ggpairs()}}
& GGally & P2 \\
\href{https://stat.ethz.ch/R-manual/R-devel/library/stats/html/kmeans.html}{\texttt{kmeans()}}
& stats & P9, P10, P11 \\
\href{https://stat.ethz.ch/R-manual/R-devel/library/base/html/Special.html}{\texttt{choose()}}
& base & P10 \\
\href{https://stat.ethz.ch/R-manual/R-devel/library/stats/html/dist.html}{\texttt{dist()}}
& stats & P9, P10 \\
\href{https://stat.ethz.ch/R-manual/R-devel/library/stats/html/hclust.html}{\texttt{hclust()}}
& stats & P10 \\
\href{https://stat.ethz.ch/R-manual/R-devel/library/stats/html/cutree.html}{\texttt{cutree()}}
& stats & P10 \\
\href{https://stat.ethz.ch/R-manual/R-devel/library/stats/html/cophenetic.html}{\texttt{cophenetic()}}
& stats & P10 \\
\href{https://stat.ethz.ch/R-manual/R-devel/library/stats/html/rect.hclust.html}{\texttt{rect.hclust()}}
& stats & P10 \\
\href{https://rdrr.io/cran/cluster/man/silhouette.html}{\texttt{silhouette()}}
& cluster & P9, P10 \\
\href{https://stat.ethz.ch/R-manual/R-devel/library/utils/html/stack.html}{\texttt{stack()}}
& utils & P11 \\
\href{https://ggplot2.tidyverse.org/reference/ggplot.html}{\texttt{ggplot()}}
& ggplot2 & P9, P11 \\
\href{https://ggplot2.tidyverse.org/reference/geom_boxplot.html}{\texttt{geom\_boxplot()}}
& ggplot2 & P11 \\
\href{https://ggplot2.tidyverse.org/reference/facet_wrap.html}{\texttt{facet\_wrap()}}
& ggplot2 & P11 \\
\end{longtable}

\chapter{Abreviaturas de las variables}\label{abreviaturas}

El archivo \texttt{R/abreviaturas.R}, que se carga en el chunk de
configuración, trae las dos funciones con las que se acortan los
nombres. La función \texttt{abreviar()} cambia cada fragmento del nombre
original por su abreviatura y quita la unidad del final, mientras que
\texttt{numericas()} arma la matriz \texttt{num} con las columnas
numéricas y el país como nombre de fila.

\Needspace{4.56in}

\begin{Shaded}
\begin{Highlighting}[]
\CommentTok{\# taller{-}02/R/abreviaturas.R}
\CommentTok{\#}
\CommentTok{\# Nombres cortos para las 57 variables de EIA.csv, que llegan a}
\CommentTok{\# tener 60 caracteres y no caben en los gráficos. Se carga desde el}
\CommentTok{\# chunk de configuración y su código se muestra en el Apéndice B.}

\CommentTok{\# Cambia cada fragmento del nombre original por su abreviatura y}
\CommentTok{\# quita la unidad del final (TBPD, BKWH, MK, MMTCD o BB). El orden}
\CommentTok{\# importa: Los fragmentos largos van antes que los cortos que}
\CommentTok{\# contienen.}
\NormalTok{abreviar }\OtherTok{\textless{}{-}} \ControlFlowTok{function}\NormalTok{(v) \{}
\NormalTok{  reglas }\OtherTok{\textless{}{-}} \FunctionTok{c}\NormalTok{(}
    \AttributeTok{Crude\_oil\_including\_lease\_condensate =} \StringTok{"Crudo"}\NormalTok{,}
    \AttributeTok{Petroleum\_and\_other\_liquids =} \StringTok{"PetLiq"}\NormalTok{,}
    \AttributeTok{Refined\_petroleum\_products =} \StringTok{"Refinados"}\NormalTok{,}
    \AttributeTok{Liquefied\_petroleum\_gases\_.LPG. =} \StringTok{"GLP"}\NormalTok{,}
    \AttributeTok{Natural\_gas\_plant\_liquids =} \StringTok{"LGN"}\NormalTok{,}
    \AttributeTok{Bunker\_residual\_fuel\_oil =} \StringTok{"BunkerRes"}\NormalTok{,}
    \AttributeTok{Bunker\_fuel =} \StringTok{"Bunker"}\NormalTok{,}
    \AttributeTok{Distillate\_fuel\_oil =} \StringTok{"Destilado"}\NormalTok{,}
    \AttributeTok{Residual\_fuel\_oil =} \StringTok{"Residual"}\NormalTok{,}
    \AttributeTok{Motor\_gasoline =} \StringTok{"Gasolina"}\NormalTok{,}
    \AttributeTok{Jet\_fuel =} \StringTok{"Jet"}\NormalTok{,}
    \AttributeTok{Kerosene =} \StringTok{"Queroseno"}\NormalTok{,}
    \AttributeTok{Other\_liquids =} \StringTok{"OtrosLiq"}\NormalTok{,}
    \AttributeTok{Other\_refined\_products =} \StringTok{"OtrosRef"}\NormalTok{,}
    \AttributeTok{Refinery\_processing\_gain =} \StringTok{"GananciaRef"}\NormalTok{,}
    \AttributeTok{Hydroelectric\_pumped\_storage\_electricity =} \StringTok{"Bombeo"}\NormalTok{,}
    \AttributeTok{Hydroelectricity =} \StringTok{"Hidro"}\NormalTok{,}
    \AttributeTok{Biomass\_and\_waste\_electricity =} \StringTok{"Biomasa"}\NormalTok{,}
    \AttributeTok{Fossil\_fuels\_electricity =} \StringTok{"Fosil"}\NormalTok{,}
    \AttributeTok{Geothermal\_electricity =} \StringTok{"Geotermica"}\NormalTok{,}
    \AttributeTok{Non.hydro\_renewable\_electricity =} \StringTok{"NoHidro"}\NormalTok{,}
    \AttributeTok{Renewable\_electricity =} \StringTok{"Renov"}\NormalTok{,}
    \AttributeTok{Nuclear\_electricity =} \StringTok{"Nuclear"}\NormalTok{,}
    \AttributeTok{Solar\_electricity =} \StringTok{"Solar"}\NormalTok{,}
    \AttributeTok{Tide\_and\_wave\_electricity =} \StringTok{"Mareas"}\NormalTok{,}
    \AttributeTok{Wind\_electricity =} \StringTok{"Eolica"}\NormalTok{,}
    \AttributeTok{Electricity =} \StringTok{"Elec"}\NormalTok{,}
    \AttributeTok{installed\_capacity =} \StringTok{"cap"}\NormalTok{,}
    \AttributeTok{net\_generation =} \StringTok{"gen"}\NormalTok{,}
    \AttributeTok{net\_consumption =} \StringTok{"cons"}\NormalTok{,}
    \AttributeTok{net\_imports =} \StringTok{"impNeta"}\NormalTok{,}
    \AttributeTok{distribution\_losses =} \StringTok{"perdidas"}\NormalTok{,}
    \AttributeTok{CO2\_emissions =} \StringTok{"CO2"}\NormalTok{,}
    \AttributeTok{consumption =} \StringTok{"cons"}\NormalTok{,}
    \AttributeTok{production =} \StringTok{"prod"}\NormalTok{,}
    \AttributeTok{exports =} \StringTok{"exp"}\NormalTok{,}
    \AttributeTok{imports =} \StringTok{"imp"}\NormalTok{,}
    \AttributeTok{reserves =} \StringTok{"reservas"}\NormalTok{)}
  \ControlFlowTok{for}\NormalTok{ (r }\ControlFlowTok{in} \FunctionTok{names}\NormalTok{(reglas))}
\NormalTok{    v }\OtherTok{\textless{}{-}} \FunctionTok{gsub}\NormalTok{(r, reglas[[r]], v, }\AttributeTok{fixed =} \ConstantTok{TRUE}\NormalTok{)}
  \FunctionTok{sub}\NormalTok{(}\StringTok{"\_(TBPD|BKWH|MK|MMTCD|BB)$"}\NormalTok{, }\StringTok{""}\NormalTok{, v)}
\NormalTok{\}}

\CommentTok{\# Las columnas numéricas, con el país como nombre de fila y los}
\CommentTok{\# nombres abreviados.}
\NormalTok{numericas }\OtherTok{\textless{}{-}} \ControlFlowTok{function}\NormalTok{(d) \{}
\NormalTok{  num }\OtherTok{\textless{}{-}}\NormalTok{ d[}\FunctionTok{vapply}\NormalTok{(d, is.numeric, }\FunctionTok{logical}\NormalTok{(}\DecValTok{1}\NormalTok{))]}
  \FunctionTok{rownames}\NormalTok{(num) }\OtherTok{\textless{}{-}}\NormalTok{ d}\SpecialCharTok{$}\NormalTok{Country}
  \FunctionTok{names}\NormalTok{(num) }\OtherTok{\textless{}{-}} \FunctionTok{abreviar}\NormalTok{(}\FunctionTok{names}\NormalTok{(num))}
\NormalTok{  num}
\NormalTok{\}}
\end{Highlighting}
\end{Shaded}

\Needspace{4.56in}

\Needspace{5.16in}

\begin{longtable}[]{@{}
  >{\raggedright\arraybackslash}p{(\linewidth - 2\tabcolsep) * \real{0.1875}}
  >{\raggedright\arraybackslash}p{(\linewidth - 2\tabcolsep) * \real{0.8125}}@{}}
\caption{Abreviatura y nombre original (con su unidad).}\tabularnewline
\toprule\noalign{}
\begin{minipage}[b]{\linewidth}\raggedright
Abreviatura
\end{minipage} & \begin{minipage}[b]{\linewidth}\raggedright
Variable
\end{minipage} \\
\midrule\noalign{}
\endfirsthead
\toprule\noalign{}
\begin{minipage}[b]{\linewidth}\raggedright
Abreviatura
\end{minipage} & \begin{minipage}[b]{\linewidth}\raggedright
Variable
\end{minipage} \\
\midrule\noalign{}
\endhead
\bottomrule\noalign{}
\endlastfoot
Bunker\_cons & Bunker fuel consumption (TBPD) \\
BunkerRes\_cons & Bunker residual fuel oil consumption (TBPD) \\
Crudo\_exp & Crude oil including lease condensate exports (TBPD) \\
Crudo\_imp & Crude oil including lease condensate imports (TBPD) \\
Crudo\_prod & Crude oil including lease condensate production (TBPD) \\
Crudo\_reservas & Crude oil including lease condensate reserves (BB) \\
Destilado\_cons & Distillate fuel oil consumption (TBPD) \\
Destilado\_prod & Distillate fuel oil production (TBPD) \\
Jet\_cons & Jet fuel consumption (TBPD) \\
Jet\_prod & Jet fuel production (TBPD) \\
Queroseno\_cons & Kerosene consumption (TBPD) \\
Queroseno\_prod & Kerosene production (TBPD) \\
GLP\_cons & Liquefied petroleum gases (LPG) consumption (TBPD) \\
GLP\_prod & Liquefied petroleum gases (LPG) production (TBPD) \\
Gasolina\_cons & Motor gasoline consumption (TBPD) \\
Gasolina\_prod & Motor gasoline production (TBPD) \\
LGN\_prod & Natural gas plant liquids production (TBPD) \\
OtrosLiq\_prod & Other liquids production (TBPD) \\
OtrosRef\_cons & Other refined products consumption (TBPD) \\
OtrosRef\_prod & Other refined products production (TBPD) \\
PetLiq\_CO2 & Petroleum and other liquids CO2 emissions (MMTCD) \\
PetLiq\_cons & Petroleum and other liquids consumption (TBPD) \\
PetLiq\_prod & Petroleum and other liquids production (TBPD) \\
Refinados\_cons & Refined petroleum products consumption (TBPD) \\
Refinados\_prod & Refined petroleum products production (TBPD) \\
GananciaRef & Refinery processing gain (TBPD) \\
Residual\_cons & Residual fuel oil consumption (TBPD) \\
Residual\_prod & Residual fuel oil production (TBPD) \\
Biomasa\_cap & Biomass and waste electricity installed capacity (MK) \\
Biomasa\_gen & Biomass and waste electricity net generation (BKWH) \\
Elec\_perdidas & Electricity distribution losses (BKWH) \\
Elec\_exp & Electricity exports (BKWH) \\
Elec\_imp & Electricity imports (BKWH) \\
Elec\_cap & Electricity installed capacity (MK) \\
Elec\_cons & Electricity net consumption (BKWH) \\
Elec\_gen & Electricity net generation (BKWH) \\
Elec\_impNeta & Electricity net imports (BKWH) \\
Fosil\_cap & Fossil fuels electricity installed capacity (MK) \\
Fosil\_gen & Fossil fuels electricity net generation (BKWH) \\
Geotermica\_cap & Geothermal electricity installed capacity (MK) \\
Geotermica\_gen & Geothermal electricity net generation (BKWH) \\
Bombeo\_cap & Hydroelectric pumped storage electricity installed
capacity (MK) \\
Bombeo\_gen & Hydroelectric pumped storage electricity net generation
(BKWH) \\
Hidro\_cap & Hydroelectricity installed capacity (MK) \\
Hidro\_gen & Hydroelectricity net generation (BKWH) \\
NoHidro\_cap & Non-hydro renewable electricity installed capacity
(MK) \\
NoHidro\_gen & Non-hydro renewable electricity net generation (BKWH) \\
Nuclear\_cap & Nuclear electricity installed capacity (MK) \\
Nuclear\_gen & Nuclear electricity net generation (BKWH) \\
Renov\_cap & Renewable electricity installed capacity (MK) \\
Renov\_gen & Renewable electricity net generation (BKWH) \\
Solar\_cap & Solar electricity installed capacity (MK) \\
Solar\_gen & Solar electricity net generation (BKWH) \\
Mareas\_cap & Tide and wave electricity installed capacity (MK) \\
Mareas\_gen & Tide and wave electricity net generation (BKWH) \\
Eolica\_cap & Wind electricity installed capacity (MK) \\
Eolica\_gen & Wind electricity net generation (BKWH) \\
\end{longtable}

% --- Referencias mínimas (fuente de datos + software usado) ---
\nocite{ideam-datos, r-core, ggplot2, ggextra}
\ifbibimpresa\else\printbibliography[heading=bibintoc]\fi
\end{document}
